% concatenated from main.tex
% \documentclass{amsart}
% \usepackage{graphicx} % Required for inserting images
% \usepackage[english]{babel}

% % Set page size and margins
% % Replace `letterpaper' with`a4paper' for UK/EU standard size
% \usepackage[letterpaper,top=2.5cm,bottom=2.5cm,left=2.25cm,right=2.25cm,marginparwidth=10cm]{geometry}
% \linespread{1.13}
% % Useful packages
% \usepackage{amsmath}
% \usepackage{graphicx}
% \usepackage{amscd,amsthm,amsfonts,graphicx,latexsym,enumitem,imakeidx}
% \usepackage{tikz-cd}
% \usepackage[affil-it]{authblk}
% \usepackage{blindtext}
% \usepackage{listings}
% \usepackage{amssymb,amsfonts}
% \usepackage[all,arc]{xy}
% \usepackage{mathrsfs}
% \usepackage{cite}
% \usepackage{enumitem}
% \usepackage{mathtools}
% \DeclarePairedDelimiter\ceil{\lceil}{\rceil}
% \DeclarePairedDelimiter\floor{\lfloor}{\rfloor}
% \newlist{steps}{enumerate}{1}
% \setlist[steps, 1]{label = Step \arabic*:}
% \usepackage[colorlinks=true, allcolors=blue]{hyperref}
% %\usepackage{refcheck}


\documentclass[10pt]{amsart}
\usepackage{comment}
\usepackage{amsmath,amsfonts,euscript,amscd,amsthm,amssymb,upref,graphics}
%%%%%%%%%%%%%%%%%%%%%%%%%%%%%%%%%%%%%%%%%%%%%%%%%%%%%%%%%%%%%%%%%%%%
\usepackage{mathrsfs,yhmath}
%%%%%%%%%%%%%%%%%%%%%%%%%%%%%%%%%%%%%%%%%%%%%%%%%%%%%%%%%%%%%%%%%%%%
\usepackage[all]{xy}
\usepackage[unicode,pdfencoding=auto,psdextra,hidelinks]{hyperref}
\usepackage{amsrefs}

% Set left margin - The default is 1 inch, so the following 
% command sets a 1.25-inch left margin.
%\setlength{\oddsidemargin}{0.25in}

% Set width of the text - What is left will be the right margin.
% In this case, right margin is 8.5in - 1.25in - 6in = 1.25in.
%\setlength{\textwidth}{6in}

% Set top margin - The default is 1 inch, so the following 
% command sets a 0.75-inch top margin.
%\setlength{\topmargin}{-0.25in}

% Set height of the header
%\setlength{\headheight}{0.3in}

% Set vertical distance between the header and the text
%\setlength{\headsep}{0.2in}

% Set height of the text
%\setlength{\textheight}{8.5in}

% Set vertical distance between the text and the
% bottom of footer
%\setlength{\footskip}{0.5in}

\setlength{\textwidth}{\paperwidth}
\addtolength{\textwidth}{-3in}
\calclayout



\usepackage{latexsym}
\usepackage{graphicx}
\usepackage{tikz}
%\usetikzlibrary{arrows,calc,angles,quotes,positioning,datavisualization, patterns}
\usepackage{caption}
\usepackage{subcaption}
%\input{table}
\usepackage{enumitem}
\usepackage{tikz-cd}


\DeclareMathOperator{\Ht}{ht}	
\DeclareMathOperator{\Lim}{Lim}
\DeclareMathOperator{\Hom}{Hom}
\DeclareMathOperator{\Ext}{Ext}
\DeclareMathOperator{\ext}{ext}
\DeclareMathOperator{\Tor}{Tor}
\DeclareMathOperator{\E}{E}
\DeclareMathOperator{\id}{id}
\DeclareMathOperator{\pd}{pd}
\DeclareMathOperator{\type}{r}
\DeclareMathOperator{\gldim}{gldim}
\DeclareMathOperator{\vdim}{vdim}
\DeclareMathOperator{\depth}{depth}
\DeclareMathOperator{\rank}{rank}
\DeclareMathOperator{\Ker}{Ker}
\DeclareMathOperator{\Coker}{Coker}
\DeclareMathOperator{\Ima}{Im}
\DeclareMathOperator{\Max}{Max}
\DeclareMathOperator{\Min}{Min}
\DeclareMathOperator{\Spec}{Spec}
\DeclareMathOperator{\Ass}{Ass}
\DeclareMathOperator{\Aut}{Aut}
\DeclareMathOperator{\Jac}{Jac}
\DeclareMathOperator{\Supp}{Supp}
\DeclareMathOperator{\Det}{Det}
\DeclareMathOperator{\Red}{Red}
\DeclareMathOperator{\Dim}{dim}
\DeclareMathOperator{\Codim}{codim}
\DeclareMathOperator{\Pic}{Pic}
\DeclareMathOperator{\Num}{Num}
\DeclareMathOperator{\Deg}{deg}
\DeclareMathOperator{\Sing}{Sing}
\DeclareMathOperator{\Div}{Div}
\DeclareMathOperator{\Vol}{Vol}

\newcommand{\cmt}[1]{\textcolor{red}{#1}}

\usepackage{bm}
\usepackage{pgfplots}

%--------Theorem Environments--------
%theoremstyle{plain} --- default
\newtheorem{theo}{Theorem}[section]

\newtheorem{exam}[theo]{Example}
\newtheorem{coro}[theo]{Corollary}
\newtheorem{lemm}[theo]{Lemma}
\newtheorem{prop}[theo]{Proposition}
\newtheorem{claim}[theo]{Claim}
\newtheorem{conj}[theo]{Conjecture}

\theoremstyle{remark}
\newtheorem{rema}[theo]{Remark}
\newtheorem{defi}[theo]{Definition}

\title{Unitriangular factorization of holomorphic symplectic matrices}
\author[G. Huang]{Gaofeng Huang}
\author[F. Kutzschebauch]{Frank Kutzschebauch}
\author[B. Tran]{Phan Quoc Bao Tran}
\address{Mathematisches Institut \\ University of Bern \\  
Bern, Switzerland}
\email{gaofeng.huang@unibe.ch, frank.kutzschebauch@unibe.ch, phan.tran@unibe.ch}

\subjclass{32Q56 (primary); 32Q28, 15A54, 32A17,
20H25 (secondary)}
\keywords{Oka principle, Stein spaces, symplectic matrices, unipotent factorization, elementary
symplectic matrices, rings of holomorphic functions} 


\begin{document}



%\date{October 2024}
\maketitle

\begin{abstract}
We prove that every holomorphic 
symplectic matrix can be factorized as a product of holomorphic unitriangular matrices with respect to the symplectic form $ \left[\begin{array}{ccc} 0 & L_n \\ -L_n & 0\end{array}\right]$ 
where $L$ is the $n \times n$ matrix with $1$ along the skew-diagonal.
Also we prove that holomorphic unitriangular matrices with respect to this symplectic form  are products of not more than $7$ holomorphic unitriangular matrices with respect to the standard symplectic form $\left[\begin{array}{ccc} 0 & I_n \\ -I_n & 0\end{array}\right]$, thus solving an open problem posed in \cite{HKS}.
Combining these two results allows for estimates of the optimal number of factors in the factorization by  holomorphic unitriangular matrices with respect to the standard symplectic form. The existence of that factorization was  obtained earlier by Ivarsson-Kutzschebauch and Schott, however without any estimates. 
Another byproduct of our results is a new, much less technical and more elegant proof of this factorization.
\end{abstract}
%\maketitle

\tableofcontents
\section{Introduction}
Factorization of a given matrix with elements in a field  has been studied widely with many applications to different fields of sciences. The most well-known one is the factorization of a matrix $A \in \mathrm{SL}_n(\mathbb{C})$ into product of some elementary matrices of the form $E + \alpha e_{ij}$, $i \neq j$, called elementary matrices. This factorization is equivalent to factorization into so-called unitriangular matrices (triangular matrices with 1s on the diagonal). Another interesting class of matrices is the class of symplectic matrices.  
Generalizing the case of matrices with entries from a field, it is much more challenging to consider matrices with entries from a ring, the subject of $K$-theory. In this paper we always work
 with  very specific rings, namely the ring of holomorphic functions on a Stein space or the ring of continuous functions on a normal finite dimensional topological space. Recall that the  unitriangular symplectic matrices  with respect to the standard symplectic form
\begin{align*}
    \Omega = \left[\begin{array}{ccc}
        0 & I_n \\
        -I_n & 0
    \end{array}\right]
\end{align*}
are of the following form
\begin{center}
     $M^-(G) = \left[\begin{array}{ccc} I_n & 0 \\ G & I_n\end{array}\right]$ and $M^+(G)=\left[\begin{array}{ccc} I_n & G \\ 0 & I_n\end{array}\right]$, \text{ where $G$ is symmetric.} 
\end{center}
Since the entries of $G$ are holomorphic functions on a reduced Stein space, it can be understood as a holomorphic map:
\begin{center}
    $G: X \to \mathbb{C}^{\frac{n(n+1)}{2}}$.
\end{center}
\iffalse
with 
\begin{center}
    $M^{+}(x) =  \left[\begin{array}{ccc} I_n & A(x) \\ 0 & I_n\end{array}\right]$
\end{center}
with 
\begin{center}
    $A(x) = \left[\begin{array}{ccccc} f_{11}(x) & f_{12}(x) & \cdots &\cdots &f_{1n}(x) \\  
    f_{12}(x) & f_{22}(x) & \ddots & \cdots & f_{2d}(x) \\
    \vdots & \ddots &\ddots &\ddots &\vdots \\
    \vdots & \vdots & \ddots &\ddots & f_{(d-1)d}(x) \\
    f_{1n}(x) & \cdots & \cdots &f_{(d-1)d}(x) & f_{dd}(x)\end{array}\right]$
\end{center}
with $f_{ij}$, $i<j$, are holomorphic maps in $\mathcal{O}(X)$. 
Notice that since the elementary symplectic matrices has the unitriangular form with a symmetric matrix with possibly $\frac{n(n+1)}{2}$ different functions, we can also think the map $M^{+}$ (or $M^-$) as: 
\begin{center}
    $M^+: X \to \mathbb{C}^{\frac{n(n+1)}{2}}$.
\end{center}
\fi
In \cite{IKL} (for $n=2$)  and \cite{Scho} (for general $n$), the following theorem about factorizing matrices into product of  unitriangular symplectic matrices  has been proved:

\begin{theo}\label{old-theo}
There exists a natural
number $K = K(n, d)$ such that given any finite dimensional reduced Stein space $X$ of dimension
$d$ and any null-homotopic holomorphic map $f : X \to \mathrm{Sp}_{2n}(\mathbb{C})$ there exists a holomorphic map:
\begin{center}
$G = (G_1, \cdots, G_K) : X \to (\mathbb{C}^{\frac{n(n+1)}{2}} )^K
$    
\end{center}
such that
\begin{center}
 $f(x) = M^-(G_1(x)) M^+(G_2(x)) 
 \cdots M^{\pm}(G_K(x)),$ \text{ where $x \in X$.}    
\end{center}
\end{theo}
%The factors occurring in this factorization are exactly the unitriangular matrices for the standard symplectic form.
It is immediate that any product of unitriangular symplectic matrices is connected by a path to the constant identity matrix $I_{2n}$, by multiplying the off-diagonal entries by $t \in [0, 1]$, i.e.
\begin{center}
    $\left[\begin{array}{ccc} I_n & tG \\ 0 & I_n\end{array}\right]$ and $\left[\begin{array}{ccc} I_n & 0 \\ tG & I_n\end{array}\right]$.
\end{center}
Therefore the requirement of null-homotopy of the map $f$ is necessary. 


In our paper, we study factorization of a holomorphic symplectic matrix into unitriangular symplectic factors with respect to the symplectic form 
\begin{center}
    $\tilde{\Omega} = \left[\begin{array}{ccc}
        0 & L_n \\
        -L_n & 0
    \end{array}\right]$, where $L$ is the $n \times n$ matrix with $1$ along the skew-diagonal.
\end{center}
The lower unitriangular $\tilde{\Omega}$-symplectic factors are the $\Omega$-symplectic matrices of the form
\begin{center}
    $\left[\begin{array}{ccc} A^{-T} & 0 \\ B & A\end{array}\right],$  $B A^T$ symmetric and $A$ upper unitriangular,
\end{center} 
In order to introduce convenient coordinates $(A,Z) \in \mathbb{C}^{\frac{n(n-1)}{2}} \times  \mathbb{C}^{\frac{n(n+1)}{2}}$ on this set of matrices we  will write them as
\begin{center}
    $M^{-}(A,Z)  = \left[\begin{array}{ccc} I_n & 0 \\ Z & I_n\end{array}\right] \left[\begin{array}{ccc} A^{-T} & 0 \\ 0 & A\end{array}\right]= \left[\begin{array}{ccc} A^{-T} & 0 \\ B & A\end{array}\right],$ $Z = B A^T$ symmetric and $A$ upper unitriangular.
\end{center}
The same for upper unitriangular
\begin{center}
    $ M^{+}(A,Z)  = \left[\begin{array}{ccc} I_n & Z \\ 0 & I_n\end{array}\right] \left[\begin{array}{ccc} A^{-1} & 0 \\ 0 & A^T \end{array}\right] = \left[\begin{array}{ccc} A^{-1} & C \\ 0 & A^{T}\end{array}\right],$ $ Z = C A^{-T}$ symmetric and $A$ upper unitriangular.
\end{center}


Throughout the paper, we call these matrices the \textbf{elementary matrices}. 
%Since the entries of $A, Z$ are holomorphic functions on a reduced Stein space, they can be understood as holomorphic maps:
%\begin{center}
%    $A: X \to \mathbb{C}^{\frac{n(n-1)}{2}}$ and $Z: X \to \mathbb{C}^{\frac{n(n+1)}{2}}$.
%\end{center}
\iffalse
Motivated by the factorization into upper and lower unitrangular matrices, we consider the matrices:
\begin{center}
    $\left[\begin{array}{ccc} A^{-T} & 0 \\ Z & A\end{array}\right] $ and $ \left[\begin{array}{ccc} A^{-1} & Z \\ 0 & A^{T}\end{array}\right]$
\end{center}
those above respectively are: 
\begin{center}
       $M^{-}(A,Z) = \left[\begin{array}{ccc} A^{-T} & 0 \\ Z & A\end{array}\right] = \left[\begin{array}{ccc} I & 0 \\ C & I\end{array}\right]\left[\begin{array}{ccc} A^{-T} & 0 \\ 0 & A\end{array}\right]$, $Z = C.A^{-T}$
    \end{center}
    and
    \begin{center}
    $M^{+}(A,Z)= \left[\begin{array}{ccc} A^{-1} & Z \\ 0 & A^{T}\end{array}\right] =\left[\begin{array}{ccc} A^{-1} & 0 \\ 0 & A^T\end{array}\right]\left[\begin{array}{ccc} I & B \\ 0 & I\end{array}\right]  $, $Z = A^{-1}B$
    \end{center}
 with $B$ and $C$ symmetric, $A$ unipotent upper triangular. These matrices are exactly the unitriangular
symplectic matrices with respect to another
symplectic form $\tilde \Omega$ widely used in $K$-theory; see
Remark \ref{symplectic form}. In this paper, we call them elementary matrices. Notice that we want $M^-$-type matrices to appear in the odd position and $M^+$-type matrices to appear in the even position of the factorization; sometimes we will use the alternative notions: 
\begin{center}
      for $k$ odd, $M_k(A_k,Z_k) = \left[\begin{array}{ccc} I & 0 \\ Z_k & I\end{array}\right]\left[\begin{array}{ccc} A^{-T}_k & 0 \\ 0 & A_k\end{array}\right]$  with $Z_k$ symmetric, $A_k$ upper triangular matrix.
    \end{center}
    and
    \begin{center}
    for $k$ even, $M_k(A_k,Z_k) =\left[\begin{array}{ccc} A^{-1}_k & 0 \\ 0 & A^T_k\end{array}\right]\left[\begin{array}{ccc} I & Z_k \\ 0 & I\end{array}\right]  $ with $Z_k$ symmetric, $A_k$ upper triangular matrix. 
    \end{center}
The choice of uniquely decomposing them for even $i$ into a product in one order and for odd $i$ in the other order is our choice
of coordinates on the set of unitriangular matrices
$\tilde \Omega $ symplectic matrices. By Proposition \ref{one-equa-lemm} both orders are possible. This clever choice
of coordinates on will be important in the reduction of
$2n$ equations to one equation \ref{one-equa-lemm}, the core point and novelty of the paper.
\fi
Our main theorem is a solution to the Gromov-Vaserstein problem for the symplectic form $\tilde \Omega$.

\begin{theo}\label{main-theo} 
There exists a natural
number $K = K(n, d)$ such that given any finite dimensional reduced Stein space $X$ of dimension
$d$ and any null-homotopic holomorphic map $f : X \to \mathrm{Sp}_{2n}(\mathbb{C})$ there exist holomorphic maps:
\begin{center}
$A = (A_1, \cdots, A_K) : X \to (\mathbb{C}^{\frac{n(n-1)}{2}} )^K
$  and $Z = (Z_1, \cdots, Z_K) : X \to (\mathbb{C}^{\frac{n(n+1)}{2}} )^K
$ 
\end{center}
such that
\begin{center}
 $f(x) = M^-(A_1(x), Z_1(x)) M^+(A_2(x), Z_2(x)) 
 \cdots M^{\pm}(A_K(x), Z_K(x)),$ \text{ where $x \in X$.}    
\end{center}
\end{theo}

This theorem together  with our second main result, saying that the unitriangular $\tilde{\Omega}$-symplectic matrices from Theorem \ref{main-theo} are products of not more than 8 
unitriangular $\Omega$-symplectic matrices, gives a new proof of Schott's result.  Moreover our proof is much more elegant and much easier to follow. In fact we use a novel strategy to prove the stratified ellipticity of a certain map, which is the core of the proof. More precisely, we can reduce the number of equations for the fibers of this map, given from the beginning by $2n$ polynomial equations, to one single equation. From this equation
the stratified ellipticity is as easily obtained as in the work of Ivarsson and Kutzschebauch \cite{KI}. In Schott's proof (and in the preceding work of Ivarsson-Kutzschebauch-L\o w)
the reduction is only done from $2n$ to $n$ equations. Studying fibers given by $n$ polynomial equations is the  hard and unpleasant technical part of the papers \cite{IKL}, \cite{Scho}.



We denote by $K_{symp}(n,d)$ the minimal number
of unitriangular factors for $2n \times 2n$ $\Omega$-symplectic matrices with holomorphic functions as entries in any $d$-dimensional Stein space and
$\tilde K_{symp}(n,d)$ the same for  $\tilde \Omega $-symplectic  factors. Observe that the optimal number of factors over the field of complex numbers is classically known in the $\tilde\Omega$-case over fields (and rings of Bass stable rank 1)  to be $4$, however, the optimal number for the $\Omega$-case has been only very recently found to be $5$ by  Jin, Tang and Zhu  in \cite{JTZ}. As  such proofs over the complex numbers usually do, their proof depends on Jordan normal form. It is the fact that Jordan normal form is discontinuous making the problem for continuous and holomorphic matrices so delicate. However, we can reduce the problem to a simpler case by multiplying the unitriangular matrix $A$ with a diagonal matrix $\Lambda$ with pairwise different elements on the diagonal. This new matrix $\Lambda A$ is diagonalizable and the diagonalization process is algebraic. From that point on  we can apply the factorization technique in \cite{JTZ} to get the optimal results as follows:

\begin{theo}\label{main-theo2}
For $n \ge 4$ we have
$$6 \le K_{symp}(n,d) \le 7 \tilde K_{symp}(n,d).$$
For matrices of smaller sizes, there is a better upper bound
\begin{align} \label{K-Ktilde}
    K_{symp}(n,d) \le 4 \tilde K_{symp}(n,d), \quad \text{ for } n = 2 \text{ and } n=3.
\end{align}
\end{theo}
This theorem solves the problem posed in \cite[p. 6]{HKS}.
Moreover, combining Theorem \ref{main-theo2}
with $K$-theoretic results from \cite{VSS},\cite{HKS}
we are in certain cases able to give explicit bounds for the number of factors in Schott's theorem as follows:

\begin{prop} \label{bounds}
    The number of factors for $N \ge n+1$ is bounded by the number for $n$:
    \begin{align} \label{K(N,d)}
        K_{symp}(N,d) \le 7 \tilde{K}_{symp}(n,d).  
    \end{align}
    When the Stein space $X$ is one-dimensional,
    \begin{align}
        5 &\le K_{symp} (n,1)  \le 16 \ \  \text{for} \ \ n=2,3, \label{K(n,1)n=2,3} \\
   6 &\le K_{symp} (n,1)  \le 28 \ \ \text{for } \  n \ge 4.  \label{K(n,1)n>=4}
    \end{align}
    When the Stein space $X$ is two-dimensional,
    \begin{align}   
   5 &\le K_{symp} (n,2)  \le 20 \ \ \text{for} \ \ n=2,3, \label{K(n,2)n=2,3}\\
   6 &\le   K_{symp}(n,2) \le  35  \ \  \text{for } \ n \ge 4. \label{K(n,2)n>=4}
  \end{align}
\end{prop}


\subsection{Strategy of the proof}
    We consider the product of $K$ unitriangular symplectic matrices
    \begin{center}
    $\Psi_K = M^-(A_1,Z_1) M^+(A_2, Z_2) \cdots  M^{\pm}(A_K, Z_{K})$. 
    \end{center} 
    If we naturally identify the set of upper unitriangular matrices with the set $\mathbb{C}^{\frac{n(n-1)}{2}}$, and the set of symmetric matrices with the set $\mathbb{C}^{\frac{n(n+1)}{2}}$, we can consider $\Psi_K$ as a map
    \begin{center}
    $\Psi_K: (\mathbb{C}^{\frac{n(n-1)}{2}} \times \mathbb{C}^{\frac{n(n+1)}{2}})^K \to \mathrm{Sp}_{2n}(\mathbb{C})$.
    \end{center}
     
     %We can identify $M^+(A,Z)$ (or $M^{-}(A,Z)$) as the holomorphic map $\mathbb{C}^{\frac{n(n-1)}{2}} \times \mathbb{C}^{\frac{n(n+1)}{2}}$ with $\frac{n(n-1)}{2}$ variables from $A$ and $\frac{n(n+1)}{2}$ variables from $B$. Thus, if we have a product of $K$ elementary matrices, we have the following holomorphic map: 
    
Given a holomorphic map $f: X \to \mathrm{Sp}_{2n}(\mathbb{C})$ from a finite dimensional Stein space $X$ to the symplectic matrices, we need to find a holomorphic lift $F$ such that the following diagram commutes 
\begin{center}
    \begin{tikzcd}
    & & & (\mathbb{C}^{\frac{n(n-1)}{2}} \times \mathbb{C}^{\frac{n(n+1)}{2}})^K \arrow[dd,"\Psi_K"] \\ \\
    & X \arrow[uurr,"F", dashed] \arrow[rr,"f"] & & \mathrm{Sp}_{2n}(\mathbb{C})
\end{tikzcd}
\end{center}
It is very hard to study the map $\Psi_K$ directly, therefore we choose to consider its composition with the projection $\pi_{2n}$ of a $2n \times 2n$ matrix to its last row
\begin{align}   \label{def-PhiK}
    \Phi_K = \pi_{2n} \circ \Psi_K : (\mathbb{C}^{\frac{n(n-1)}{2}} \times \mathbb{C}^{\frac{n(n+1)}{2}} )^K \to \mathbb{C}^{2n}\backslash \lbrace 0 \rbrace.  
\end{align}
By finding a holomorphic lift in 
\begin{center}
    \begin{tikzcd}
    & & & (\mathbb{C}^{\frac{n(n-1)}{2}} \times \mathbb{C}^{\frac{n(n+1)}{2}})^K \arrow[dd,"\Phi_K"] \\ \\ 
    & X \arrow[rr,"\pi_{2n} \circ f"] \arrow[uurr, dashed] &  & \mathbb{C}^{2n}\setminus \{ 0\}
    \end{tikzcd}
\end{center}
we obtain a product of holomorphic elementary matrices whose last row is the same as the last row of $f$. Using standard linear algebra arguments, Theorem \ref{main-theo} follows by an induction on the size of the matrices. 

To find this holomorphic lift, we will use the Oka principle for stratified elliptic submersions. In fact, after removing the singularity set $S_K$ from its domain, the map $\Phi_K$ becomes such. 
\begin{prop} \label{prop: Phi-K}
The map $\Phi_K: (\mathbb{C}^{\frac{n(n-1)}{2}} \times \mathbb{C}^{\frac{n(n+1)}{2}})^K \setminus S_K \to \mathbb{C}^{2n}\backslash \lbrace 0 \rbrace$ is a stratified elliptic submersion.    
\end{prop}
The continuous lift is constructed using the continuous factorization of symplectic matrices in \cite{IKL}.

The proof of Proposition \ref{prop: Phi-K} uses complete holomorphic vector fields tangent to the fibers of the map $\Phi_K$. We are able to reduce the $2n$ equations describing these fibers, to a single equation. This reduction in the number of defining equations greatly simplifies the search for a set of complete vector fields spanning the tangent space of the fiber at each point. 

\iffalse
In order to prove the Main Theorem \ref{main-theo}, we will apply the Oka principle as in \cite{KI} and \cite{Scho}, based on the existence of the continuous Gromov-Vaserstein solution (Theorem \ref{Cont-GV}). The first important part of any version of the Gromov-Vaserstein problem is to analyze the holomorphic map $\Psi_K$ given by:

The continuous Gromov-Vaserstein solution gives us a natural number $K(n,d)$ and a continuous lift $F: X \to (\mathbb{C}^{\frac{n(n-1)}{2}} \times \mathbb{C}^{\frac{n(n+1)}{2}})^K$, which yields the diagram: 

To prove the holomorphic version of Symplectic Gromov-Vaserstein problem, we want to verify the assumptions of an Oka Principle for the map $\Phi_K$. Unfortunately, $\Phi_K$ is not a submersion everywhere. Indeed, one can find the singular set $S_K$  of this map by explicitly computation (cf. Theorem \ref{sing-locus}) and the best we can have is the open manifold $U_K = (\mathbb{C}^{\frac{n(n-1)}{2}} \times \mathbb{C}^{\frac{n(n+1)}{2}} )^K \backslash S_K$ such that 
\begin{center}
    $\Phi_K|_{U_K}: U_K \to \mathbb{C}^{2n} \backslash \lbrace 0 \rbrace$
\end{center}
is a stratified elliptic submersion (cf. Theorem \ref{theo-surj} and Theorem \ref{ellipticity}).\\
%\begin{theo}
%The singular locus, i.e., the non-submersive locus of $\Phi_K$, is \[S_K = \lbrace (A_K,Z_K)| Z_K \in W^c_K : \mathrm{rank}(W_K(Z_K)) < n; a_{1n,2j} =  \cdots = a_{(n-1)n,2j} = 0 | 1 \leq 2j \leq K \rbrace.\] 
%\end{theo}
%with $W_K = \lbrace Z_K \in (\mathbb{C}^{\frac{n(n+1)}{2}})^K : Z_{2i-1}e_n \neq 0 \xspace$ for some $\xspace 1 \leq i \leq \ceil{\frac{K-1}{2}} \rbrace$, $S_{Z_K} = \lbrace Z_K \in W^c_K : \mathrm{rank}(\mathcal{W}_K) < n \rbrace$ with $\mathcal{W}_K = [Z_2|Z_4|\cdots|Z_{2k}]$ and $k = \ceil{\frac{K-1}{2}}$. 


To prove the stratified ellipticity, we need to construct a \textbf{Stratified spray} of $U_K$ associated with $\Phi_K|_{U_K}$. One way is to find finitely many complete vector fields that are tangent to the fibers and span the tangent space of the fibers. In \cite{KI}, the authors found a way to reduce $n$ polynomial equations that determine the fibers, to one equation. With that, it is much easier to construct complete vector fields that leave the polynomial equation invariant. And with the new factorization for the $\mathrm{Sp}_{2n}(\mathbb{C})$ case, we also have the same property of the fibers of the fibration as in \cite{KI}: That is, every fiber of $\Phi_K|_{U_K}$ can be determined by one equation. This finally leads to the following theorem:

This phenomenon could not happen in \cite{Scho} since it considers different elementary matrices and the author has to work with huge computations on the equations $n$ that determined the fibers. As a consequence, the proof for ellipticity of $\Phi_K$ in \cite{Scho} has to rely on the sophisticated list of complete vector fields on $U_K$ (in \cite{Scho} it is $E_K$). The new way of factorizing matrices therefore gives us a much shorter and elegant proof as in \cite{KI}[Lemma 3.7]. Because we have to work with only one equation, it is much easier to find the complete vector fields that span the tangent vectors at each point in the fibers of $\Phi_K|_{U_K}$.

In the last section, we apply Theorem \ref{main-theo} to show the relation between $ \mathrm{K}_{Sym}(n,d)$ and $\tilde{\mathrm{K}}_{Sym}(n,d)$. A key observation is that when we change the basis to get $\tilde{\Omega}$-symplectic matrix from $\Omega$-symplectic matrix, the elementary matrix appears in the factorization of Theorem \ref{main-theo} becomes unitriangular. Thus, Theorem \ref{main-theo} actually gives us the unitriangular factorization for $\tilde{\Omega}$-symplectic matrix. In order to relate to the unitriangular factorization of $\Omega$-symplectic matrix as in Theorem \ref{old-theo}, we try to find optimal factorization of the new elementary matrices. In \cite{JLX2}, the authors provide a method to study the number of factorization of symplectic matrices in $\mathrm{Sp}_{2n}(\mathbb{C})$ by using Jordan normal form. Unfortunately Jordan normal form is not a "continuous" method, i.e., we could not expect the Jordan decomposition to exist in the ring of matrices with holomorphic function entries. So we need to modify their method a bit. We first factor the matrix $\left[\begin{array}{ccc} A & 0 \\ 0 & A^{-T}\end{array}\right] $ into a product of two matrices, one of them being a diagonalizable matrix in $\mathrm{Mat}(2n,\mathcal{O}(X))$ with pairwise different constant eigenvalues. This matrix can be decomposed into the product of the form $K.\Lambda.K^{-1}$ with $\Lambda$ a diagonal matrix and $K$ an invertible matrix in $\mathrm{Mat}(2n,\mathcal{O}(X))$. From now on we can freely use the factorization described in \cite{JLX2}. This gives us an eight-factors symplectic factorization. Thus, we have the upper bound for the number of factors. For the lower bound, we directly study the equation: 
\begin{center}
    $\left[\begin{array}{ccc} A & B \\ 0 & D\end{array}\right]  = \left[\begin{array}{ccc} I & 0 \\ U & I\end{array}\right].\left[\begin{array}{ccc} I & V \\ 0 & I\end{array}\right].\left[\begin{array}{ccc} I & 0 \\ W & I\end{array}\right].\left[\begin{array}{ccc} I & X \\ 0 & I\end{array}\right]$
\end{center}
and find that they have holomorphic functions solutions for the case $n=2$ and $n=3$, but many cases they have no solution when $n \geq 4$. The reason we consider $4$ for the number of factors is that in $\mathrm{Sp}_{2n}(\mathbb{C})$, the minimal number of symplectic factorizations is $4$. From the fact that matrix in $\mathrm{Sp}_{2n}(\mathcal{O}(X))$ for $n \geq 4$, we can also show that $K_{\mathrm{symp}}(n,d) > 5$. And thus this implies Theorem \ref{main-theo2}.  
\fi

%\section{Preliminaries and conventions}
\iffalse
\subsection{Notation and basic lemmas}
\textbf{Notation}: Let $X$ a Stein manifold, we denote $\mathcal{O}(X)$ be the ring of holomorphic functions on $X$.\\
Let $\mathrm{Sym}(\mathbb{C}^n)$ be the set of symmetric matrices with complex entries. We can identify $\mathrm{Sym}(\mathbb{C}^n)$ with $\mathbb{C}^{\frac{n(n+1)}{2}}$. \\
We denote $\mathrm{Mat}(n,R)$ be the ring of $n \times n$ matrices with entries in the ring $R$. 
\begin{defi}
A $2n \times 2n$ block matrix $H = \left[\begin{array}{ccc} A & B \\ C & D\end{array}\right] \in \mathrm{Mat}(n,\mathcal{O}(X))$ is called \textbf{symplectic} if  $H^T\Omega H = \Omega$ with $\Omega =  \left[\begin{array}{ccc} 0 & I_n \\ -I_n & 0\end{array}\right]$.
\end{defi}
\begin{prop}\label{fact}
Any symplectic matrix of form $H = \left[\begin{array}{ccc} A & B \\ C & D\end{array}\right] \in \mathcal{O}(X)$ has the property:
\begin{enumerate}
    \item $A^TC = C^TA$.
    \item $B^TD=D^TB$
    \item $A^TD-C^TB = I_n$.
\end{enumerate}
\end{prop}
We have useful factorizations of  
%Now we can define the elementary matrices:
%\begin{defi}
%Let $A$ be an upper triangular matrix with entries $1$ on the main diagonal. Symplectic matrices of the form: 
%\begin{center}
%   $\left[\begin{array}{ccc} A^{-T} & 0 \\ Z & A\end{array}\right]$
%\end{center}
%and
%\begin{center}
%    $ \left[\begin{array}{ccc} A^{-1} & Z \\ 0 & A^T\end{array}\right]$
%\end{center}
%are called \textbf{elementary matrices}.
%\end{defi}
\begin{prop}\label{fact-1}
If $H = \left[\begin{array}{ccc} A_1 & B_1 \\ A_2 & B_2\end{array}\right] \in \mathrm{Mat}(2n,\mathcal{O}(X))$ is a symplectic matrix and $A_i, B_i \in \mathrm{Mat}(n,\mathcal{O}(X))$ and $A_1$ is  an invertible matrix, then H has three unique factorizations 
\begin{equation}
    H = \left[\begin{array}{ccc} I & 0 \\ S & I\end{array}\right]\left[\begin{array}{ccc} P & 0 \\ 0 & P^{-T}\end{array}\right]\left[\begin{array}{ccc} I & T \\ 0 & I\end{array}\right],   S = A_2.A_1^{-1}, T=A_1^{-1}.B_1, P = A_1 
\end{equation}
or
\begin{equation}
     H = \left[\begin{array}{ccc} P & 0 \\ 0 & P^{-T}\end{array}\right]\left[\begin{array}{ccc} I & 0 \\ S & I\end{array}\right]\left[\begin{array}{ccc} I & T \\ 0 & I\end{array}\right], S = A^T_1A_2, T = A^{-1}_1.B_1, P = A_1
\end{equation}
or
\begin{equation}
     H = \left[\begin{array}{ccc} I & 0 \\ S & I\end{array}\right]\left[\begin{array}{ccc} I & T \\ 0 & I\end{array}\right]\left[\begin{array}{ccc} P & 0 \\ 0 & P^{-T}\end{array}\right], S = A_2A^{-1}_1, T = B_1A^T_1, P = A_1.
\end{equation}
Where $S$ and $T$ are symmetric and $P$ is nonsingular.
\end{prop}

\begin{proof}
    The proof is the same as \cite{JTZ}[Property 4]. 
\end{proof}
With Proposition \ref{fact}, we are ready to define the elementary matrices:
\begin{defi}\label{def-ele}
     The \textbf{elementary matrices} are the symplectic matrices of the form: 
    \begin{center}
       $M^{-}(A,Z) = \left[\begin{array}{ccc} A^{-T} & 0 \\ Z & A\end{array}\right] $,  $A$ unipotent upper triangular, 
    \end{center}
  or
    \begin{center}
    $M^{+}(A,Z)= \left[\begin{array}{ccc} A^{-1} & Z \\ 0 & A^{T}\end{array}\right]$, $A$ unipotent upper triangular. 
\end{center}

\end{defi}
By the Proposition \ref{fact-1}, we can uniquely factorize $M^-$ and $M^+$ in different ways: 
\begin{coro}\label{fact-imp}
For the elementary matrices, we have:
  \begin{center}
    $M^{-}(A,Z) =\left[\begin{array}{ccc} I & 0 \\ C & I\end{array}\right]\left[\begin{array}{ccc} A^{-T} & 0 \\ 0 & A\end{array}\right]= \left[\begin{array}{ccc} A^{-T} & 0 \\ 0 & A\end{array}\right] \left[\begin{array}{ccc} I & 0 \\ C' & I\end{array}\right]$,  $Z = CA^{-T} = A^{-T}C^{\prime}$ 
\end{center}
and
\begin{center}
    $M^{+}(A,Z) = \left[\begin{array}{ccc} A^{-1} & 0 \\ 0 & A^T\end{array}\right]\left[\begin{array}{ccc} I & B \\ 0 & I\end{array}\right] =  \left[\begin{array}{ccc} I & B^{\prime} \\ 0 & I\end{array}\right] \left[\begin{array}{ccc} A^{-1} & 0 \\ 0 & A^T\end{array}\right]$, $Z = B.A^{-1} = A^{-1}.B^{\prime}$,
\end{center}  
with $B$, $B'$, $C$, $C'$ symmetric and $A$ upper-triangular. 
\end{coro}

%\begin{proof}
% By proposition \ref{fact}, we must have $A^T.C = I_n$, this implies $C = A^{-T}$.   
%\end{proof}
\begin{rema}\label{symplectic form} 
Our term elementary matrices has the following meaning: 
For the standard symplectic form $\Omega =  \left[\begin{array}{ccc} 0 & I_n \\ -I_n & 0\end{array}\right]$ the unitriangular
symplectic matrices (triangular matrices with 1 on the diagonal) are exactly
the matrices considered in the works of Ivarsson,Kutzschebauch,L\o w and Schott namely the matrices of the form
$\left[\begin{array}{ccc} I_n & 0 \\ A & I_n\end{array}\right]$
with $A$ symmetric. 

However $K$-theorists use another symplectic form \begin{center}
    $\tilde{\Omega} =  \left[\begin{array}{ccc} 0 & L_n \\ -L_n & 0\end{array}\right]$
\end{center}
where $L$ is the $n \times n$ matrix with $1$ along the skew-diagonal.

$\tilde{\Omega}$ represents the symplectic form $\Omega$ in another basis.
This new basis can be obtained from the old one by reversing the order of the last $n$ basis elements.  In fact, the above defined elementary matrices are simultaneously upper or lower triangular in this new basis. Moreover they are exactly the unitriangular
 $\tilde{\Omega}$-symplectic matrices. 

%Consider the matrix $\left[\begin{array}{ccc} A & 0 \\ 0 & D\end{array}\right]$ and $D=A^{-T}$. If we assume $A$ is upper triangular, then $D$ is lower triangular. In fact, the above defined elementary matrices are simultaneously upper or lower triangular in another basis. This new basis can be obtained from the old one by reversing the order of the last $n$ basis elements, giving a symplectic form represented by the matrix:
%\begin{center}
%    $\tilde{\Omega} =  \left[\begin{array}{ccc} 0 & L_n \\ -L_n & 0\end{array}\right]$
%\end{center}
%where $L$ is the $n \times n$ matrix with $1$ along the skew-diagonal. Notice that symplectic matrices of type $\left[\begin{array}{ccc} I & 0 \\ C & I\end{array}\right]$ and $\left[\begin{array}{ccc} I & B \\ 0 & I\end{array}\right]$ are remain lower or upper respectively. The elementary matrices $M^{-}(C,A)$ and $M^{+}(B,A)$ become lower or upper unitriangular with respect to $\tilde{\Omega}$. They are in fact exactly the unitriangular $\tilde{\Omega}$-symplectic matrices. matrices with respect The symplectic matrices with respect to $\tilde{\Omega}$ are called $\tilde{\Omega}$-symplectic matrices. 
Unitriangular $\tilde{\Omega}$-symplectic matrices are widely used in K-theory. They have the advantage that they allow to obtain strong results on the number of factors in the factorization as a product of unitriangular matrices (cf. \cite{HKS}[Section 2] for more details). 
\end{rema}



\subsection{Continuous Gromov-Vaserstein problem}
As in \cite{KI} and \cite{Scho}, a priori to the use of Oka principle is the appearance of a continuous solution. Thus, it is necessary to mention the existence of a solution for continuous Gromov-Vaserstein problem in the symplectic case:
\begin{theo}
  There exists a natural number $K(n, d)$ such that given any finite dimensional normal topological space $X$ of dimension $d$ and any null-homotopic continuous mapping $M : X \longrightarrow \mathrm{Sp}_{2n}(\mathbb{C})$ there exist $K$ continuous mappings
\begin{center}
 $G_1, \cdots, G_K : X \longrightarrow \mathbb{C}^{\frac{n(n+1)}{2}}$   
\end{center}
such that
$M(x) = M_1(G_1(x))\cdots M_K(G_K(x))$.
\end{theo}
This theorem immediately implies our needed continuous solution:
\begin{theo}\label{Cont-GV}
    There exists a natural number $K(n, d)$ such that given any finite dimensional normal topological space $X$ of dimension $d$ and any null-homotopic continuous mapping $M : X \longrightarrow \mathrm{Sp}_{2n}(\mathbb{C})$ there exist $K$ continuous mappings
    \begin{center}
       $M =(A_1,Z_1), \cdots, (A_K,Z_K) : X \longrightarrow \mathrm{Sp}_{2n}(\mathbb{C})$ 
    \end{center}
    such that 
    \begin{center}
        $M(x) =   M^-(A_1(x),Z_1(x))\cdots M^-(A_K(x),Z_K(x))$ with $K$ odd 
    \end{center}
    and
    \begin{center}
         $M(x) =   M^-(A_1(x),Z_1(x))\cdots M^+(A_K(x),Z_K(x))$ with $K$ even. 
    \end{center}
\end{theo}
By Corollary, we can factor the map $M$ above as: 
\begin{center}
$M: X \longrightarrow (\mathbb{C}^{\frac{n(n-1)}{2}} \times \mathbb{C}^{\frac{n(n+1)}{2}})^K \longrightarrow \mathrm{Sp}_{2n}(\mathbb{C})$   
\end{center}
\fi

\section{The Oka principle for stratified elliptic submersions}
%\begin{defi}
%Let $X$ and $Z$ be complex spaces. A holomorphic map $h: Z \longrightarrow X$ is said to bea submersion if for each point $z_0 \in Z$ it is locally equivalent via a fiber-preserving biholomorphic map to a projection $p: U \times V \longrightarrow U$, where $U \subset X$ is an open set containing $h(z_0)$ and $V$ is an open set in some $\mathbb{C}^d$.    
%\end{defi}

%\begin{defi}
%Let $h: Z \to X$ be a holomorphic submersion of a complex manifold $Z$ onto a complex manifold $X$. A spray on $Z$ associated with h is a triple $(E, p, s)$, where $p: E \to Z$ is a holomorphic vector bundle and $s: E \to Z$ is a holomorphic map such that for each $z \in Z$ we have
%\begin{enumerate}
%    \item $s(E_z) \subset Z_h(z)$
%    \item $s(0_z) = z$, and
%    \item the derivative $D_s(0_z): T_{0_z}E \to T_zZ$ maps the subspace $Ez \subset T_{0_z}E$ surjectively onto thev ertical tangent space $VT_zZ$.
%\end{enumerate}
%\end{defi}

In this section, we recall the Oka principle that we will use as a crucial tool in our proof of Theorem \ref{main-theo} (see \cite{For1, Gro}), following the monograph of Franc Forstneri{\v c} \cite{For2}. 

Let $h\colon Z \to X$ be a holomorphic submersion of a complex manifold $Z$ onto a complex manifold $X$. For any $x\in X$ the fiber over $x$ of this submersion will be denoted by $Z_x$. At each point $z\in Z$ the tangent space $T_zZ$ contains the vertical tangent space $VT_zZ=\ker Dh$. For holomorphic vector bundles $p\colon E \to Z$ we denote the zero element in the fiber $E_z$ by $0_z$.

% \begin{Def} \label{definspray}
% Let $h\colon Z \to X$ be a holomorphic submersion of a complex manifold $Z$ onto a complex manifold $X$. A spray on $Z$ associated with $h$ is a triple $(E,p,s)$, where $p\colon E\to Z$ is a holomorphic vector bundle and $s\colon E\to Z$ is a holomorphic map such that for each $z\in Z$ we have 
% \begin{itemize}
% \item[(i)]{$s(E_z)\subset Z_{h(z)}$,}
% \item[(ii)]{$s(0_z)=z$, and}
% \item[(iii)]{the derivative $Ds(0_z)\colon T_{0_z}E\to T_zZ$ maps the subspace $E_z\subset T_{0_z}E$ surjectively onto the vertical tangent space $VT_zZ$.}
% \end{itemize} 
% \end{Def} 

% \begin{Rem}
% We will also say that the submersion admits a spray. A spray associated with a holomorphic submersion is sometimes called a (fiber) dominating spray.
% \end{Rem}



\begin{defi}
Let $X$ and $Z$ be complex spaces. A holomorphic map $h: Z \to X$ is said to be a submersion if for each point $z_0\in Z$ it is locally equivalent via a fiber preserving biholomorphic map to a projection $p\colon U\times V \to U$, where $U\subset X$ is an open set containing $h(z_0)$ and $V$ is an open set in some $\mathbb{C}^d$.  
\end{defi}

%We will need to use stratified sprays which are defined as follows.

% \begin{Def} 
% \label{definstratifiedspray}
% We say that a submersion $h\colon Z\to X$ admits stratified sprays if there is a descending chain of closed complex subspaces $X=X_m\supset \cdots \supset X_0$ such that each stratum $Y_k = X_k\setminus X_{k-1}$ is regular and the restricted submersion $h\colon Z|_{Y_k}\to Y_k$ admits a spray over a small neighborhood of any point $x\in Y_k$.  
% \end{Def}
% \begin{Rem}
% We say that the stratification $X=X_m\supset \cdots \supset X_0$ is associated with the stratified spray and vice versa.
% \end{Rem}


\begin{defi}
Let $h\colon Z \to X$ be a holomorphic submersion of a complex manifold $Z$ onto a complex manifold $X$. A fiber spray on $Z$ associated with $h$ is a triple $(E,p,s)$, where $p\colon E\to Z$ is a holomorphic vector bundle and $s\colon E\to Z$ is a holomorphic map such that for each $z\in Z$ we have 
\begin{itemize}
\item[(i)]{$s(E_z)\subset Z_{h(z)}$,}
\item[(ii)]{$s(0_z)=z$, and}
\item[(iii)]{the derivative $Ds(0_z)\colon T_{0_z}E\to T_zZ$ maps the subspace $E_z\subset T_{0_z}E$ surjectively onto the vertical tangent space $VT_zZ$.}
\end{itemize} 
\end{defi}

\begin{rema} \label{VF-spray}
    One way of constructing dominating fiber sprays, as pointed out by Gromov, is to find finitely many $\mathbb{C}$-complete vector fields that are tangent to the fibers and span the tangent space of the fibers at all points in $Z$. One can then use the flows $\varphi_j^t$ of these vector fields $V_j$ to define $s\colon Z\times \mathbb{C}^N\to Z$ via $s(z,t_1,\dots, t_N)=\varphi_1^{t_1}\circ \dots \circ \varphi_N^{t_N}(z)$ which gives a spray.
\end{rema}
 

\begin{defi}
We say that a holomorphic submersion $h: Z \to X$ is \textbf{stratified elliptic} if there is a descending
chain of closed complex subspaces $X = X_m \supset \cdots \supset X_0$ such that each stratum $Y_k = X_k \backslash X_{k-1}$
is regular and the restricted submersion $h: Z|_{Y_k} \to Y_k$ admits a dominating fiber spray over a small neighborhood of any point $x \in Y_k$.
\end{defi}

Let $\pi: E \to B$ be a holomorphic map of a complex space $E$ onto a complex space $B$. (All complex spaces are assumed to be reduced.) Assume that $X$ is a Stein space, $P_0 \subset P$ are compact Hausdorff spaces (the parameter spaces), and $f : P \times X \to B$ is a continuous $X$-holomorphic map, meaning that the map $f(p, \cdot): X \to B$ is holomorphic for every fixed $p \in P$. Assume that $F_0: P \times X \to E$ is a continuous lifting of $f$, i.e., $\pi \circ F_0 = f$, which is $X$-holomorphic over $P_0$. We are looking for a homotopy of liftings $F_t : P \times X \to E$ of $f$ ($t \in [0,1]$) which is fixed on $P_0 \times X$, such that $F_1 = F$ is an $X$-holomorphic lifting. The situation is illustrated with the following diagram:

\[
\begin{tikzcd}
& P_0 \times X \arrow[r, "F_0"] \arrow[d, hook, "incl"'] & E \arrow[d, "\pi"]   \\
& P \times X \arrow[ur, dashed, "F"] \arrow[r, "f"'] & B
\end{tikzcd}
\]



% \begin{defi}
% A holomorphic map $\pi: E \to B$ between reduced complex spaces enjoys the \textbf{Parametric Oka Property with Approximation and Interpolation (POPAI)} if for any collection $(X, X^\circ, K, P, P_0, f, F_0)$, where $X$ is a reduced Stein space, $X^\circ$ is a closed complex subvariety of $X$, $K$ is a compact $\mathcal{O}(X)$-convex subset of $X$, $P_0 \subset P$ are compact Hausdorff spaces, $f : P \times X \to B$ is a continuous $X$-holomorphic map, and $F_0 : P \times X \to E$ is a continuous map such that $\pi \circ F_0 = f$, and $F_0(p, \cdot)$ is holomorphic on $X$ for all $p \in P_0$ and on $K \cap X^\circ$ for all $p \in P$, there exists a homotopy $F_t : P \times X \to E$ such that for all $t \in [0,1]$:
% \begin{itemize}
%   \item[(i)] $\pi \circ F_t = f$,
%   \item[(ii)] $F_t = F_0$ on $(P_0 \times X) \cap (P \times X^\circ)$,
%   \item[(iii)] $F_t$ is $X$-holomorphic on $K$ and uniformly close to $F_0$ on $P \times K$,
%   \item[(iv)] $F_1 : P \times X \to E$ is $X$-holomorphic.
% \end{itemize}
% The map $\pi: E \to B$ enjoys the \textbf{Basic Oka Property with Approximation and Interpolation (BOPAI)} if the above holds when $P$ is a singleton and $P_0 = \emptyset$.
% \end{defi}
We say that the map $\pi: E \to B$ satisfies the \textbf{parametric Oka property} if such a homotopy $F_t$ exists for any data $(X, P_0, P, f, F_0)$. 
It can be illustrated by the following commutative diagram:

\[
\begin{tikzcd}
P_0 \arrow[d, hook, "incl"] \arrow[rr, "F_0"] & & \mathcal{O}(X, E)  \arrow[d, ""] \arrow[rr, hook, "incl"]  & & \mathcal{C}(X, E) \arrow[d, "\pi"] \\
P \arrow[urr, "F_1"] \arrow[urrrr, "\quad F_0"'] \arrow[rr, "f"'] & &\mathcal{O}(X, B) \arrow[rr, hook, "incl"'] & & \mathcal{C}(X, B)
\end{tikzcd}
\]

%Note that a complex manifold $E$ enjoys a certain Oka property if and only if the trivial map $\pi: E \to \{\mathrm{point}\}$ does. For a complex manifold $Y$, all Oka properties for parameter spaces $P_0 \subset P \subset \mathbb{R}^m$ are equivalent, and such $Y$ is an Oka manifold.

Note that the problem of lifting a holomorphic map $f : X \to B$ to a holomorphic map $F : X \to E$ reduces to the problem of finding a section of the pullback of $\pi: E \to B$ by $f$. 

\begin{theo} \label{POP}
    Every stratified elliptic submersion enjoys the parametric Oka property.
\end{theo} 

\section{Proof of main theorem}\label{3}
In this section we explain the inductive proof of a parametric version of Theorem \ref{main-theo}. 

\begin{theo}\label{main-theo-parametric} 
There exists a natural
number $K = K(n, d)$ such that given any compact Hausdorff space $P$, any finite dimensional reduced Stein space $X$, such that $P \times X$ has covering dimension $d$, and any null-homotopic holomorphic map $f : P \times X \to \mathrm{Sp}_{2n}(\mathbb{C})$, there exist $X$-holomorphic maps:
\begin{center}
$A = (A_1, \cdots, A_K) : P \times X \to (\mathbb{C}^{\frac{n(n-1)}{2}} )^K
$  and $Z = (Z_1, \cdots, Z_K) : P \times X \to (\mathbb{C}^{\frac{n(n+1)}{2}} )^K
$ 
\end{center}
such that
\begin{center}
 $f(p, x) = M^-(A_1(p,x), Z_1(p,x)) M^+(A_2(p,x), Z_2(p,x)) 
 \cdots M^{\pm}(A_K(p,x), Z_K(p,x)),$ \text{ where $p \in P, x \in X$.}    
\end{center}
\end{theo}


We abbreviate an element $((A_1,Z_1),(A_2,Z_2),\cdots,(A_K,Z_K)) \in (\mathbb{C}^{\frac{n(n-1)}{2}} \times \mathbb{C}^{\frac{n(n+1)}{2}})^K$ by $(\Vec{A}_K,\Vec{Z}_K )$. %Here $\Vec{Z^{\prime}}_K$ is $(Z^{\prime}_1,Z^{\prime}_2,\cdots,Z^{\prime}_K) \in (\mathbb{C}^{\frac{n(n+1)}{2}})^K$ and $\Vec{A}_K$ is $(A_1,A_2,\cdots,A_K) \in (\mathbb{C}^{\frac{n(n-1)}{2}})^K$.\\
A lengthy but straightforward calculation yields the following.

\begin{prop} \label{SingSet}
    The singularity set $S_K$ of $\Phi_K: (\mathbb{C}^{\frac{n(n-1)}{2}} \times \mathbb{C}^{\frac{n(n+1)}{2}})^K  \to \mathbb{C}^{2n}\backslash \lbrace 0 \rbrace$ defined in \eqref{def-PhiK} is given by
    \begin{align*}
        S_K = \lbrace (\Vec{A}_K,\Vec{Z}_K)|Z_{i} e_n = 0, i<K;  A_j e_n = e_n, 1 < j < K \rbrace.
    \end{align*}
\end{prop}
In other words, in the singular set the last columns of all matrices $Z_i$, except $Z_K$, should be 0. Also, the last columns of all $A_i$, except $A_1$ and $A_K$, should be $e_n$.

The following result will be proven in Section \ref{sec-stra}:
\begin{prop} \label{reg-locus}
For $n \ge 2, K \ge 3$, the map $\Phi_K: (\mathbb{C}^{\frac{n(n-1)}{2}} \times \mathbb{C}^{\frac{n(n+1)}{2}})^K \setminus S_K \to \mathbb{C}^{2n}\backslash \lbrace 0 \rbrace$ is a stratified elliptic submersion.    
\end{prop}


\begin{prop}\label{sup-main} 
Given a null-homotopic, continuous $X$-holomorphic map $f : P \times X \to \mathrm{Sp}_{2n}(\mathbb{C})$, where $P$ is a compact Hausdorff space and $X$ a finite dimensional reduced Stein space. Assume that there exist a natural number $L$ and a continuous lift $F : P \times X \to (\mathbb{C}^{\frac{n(n-1)}{2}} \times \mathbb{C}^{\frac{n(n+1)}{2}})^L$ of $f$ along $\Psi_L$, i.e., $\Psi_L \circ F = f$.
Then there exist a continuous lift $\tilde F : P \times X \to (\mathbb{C}^{\frac{n(n-1)}{2}} \times \mathbb{C}^{\frac{n(n+1)}{2}})^{L+2} \setminus S_{L+2}$ of $f$ along $\Psi_{L+2}$, avoiding the singularity set and a continuous homotopy $h_t: P \times X \to (\mathbb{C}^{\frac{n(n-1)}{2}} \times \mathbb{C}^{\frac{n(n+1)}{2}})^{L+2}\setminus S_{L+2}$ of lifts of $\pi_{2n} \circ f$ along $\Phi_{L+2}$, such that $h_0 = \pi_{2n} \circ \tilde F $ and $h_1$ is $X$-holomorphic.
\end{prop}
\begin{proof}
By assumption we have
\begin{center}
    $F = (F_1, \cdots, F_L) : P \times X \to (\mathbb{C}^{\frac{n(n-1)}{2}} \times \mathbb{C}^{\frac{n(n+1)}{2}} )^L$ with $f = \Psi_L \circ F, \, F_k =  (A_k,Z_k)$.   
\end{center}
We define the map
$\tilde F : P \times X \to (\mathbb{C}^{\frac{n(n-1)}{2}} \times \mathbb{C}^{\frac{n(n+1)}{2}} )^{L+2}$
by
 \begin{center}
  $\tilde F = ((I,I), (I,0), F^{\prime}_1, F_2, \cdots, F_L)$,  where $F_1'=(A_1, Z_1-I)$.
 \end{center}
Since
\begin{center}
    $ \left[\begin{array}{ccc} I & 0 \\ Z_1 & I\end{array}\right]\left[\begin{array}{ccc} A_1^{-T} & 0 \\ 0 & A_1\end{array}\right] = \left[\begin{array}{ccc} I & 0 \\ I & I\end{array}\right]\left[\begin{array}{ccc} I & 0 \\ 0 & I\end{array}\right] \left[\begin{array}{ccc} I & 0 \\ Z_1-I & I\end{array}\right]\left[\begin{array}{ccc} A_1^{-T} & 0 \\ 0 & A_1 \end{array}\right] $
\end{center}
% and
% \begin{center}
%     $F^{\prime}_1 = \left[\begin{array}{ccc} I & 0 \\ Z_1-I & I\end{array}\right]\left[\begin{array}{ccc} A_1^{-T} & 0 \\ 0 & A_1\end{array}\right]$
% \end{center}
%  with the factorization:
%  \begin{center}
%     $F_1 =\left[\begin{array}{ccc} I & 0 \\ I & I\end{array}\right]\left[\begin{array}{ccc} I & 0 \\ 0 & I\end{array}\right] \left[\begin{array}{ccc} I & 0 \\ Z_1-I & I\end{array}\right]\left[\begin{array}{ccc} A_1^{-T} & 0 \\ 0 & A_1 \end{array}\right]$
%  \end{center}
we have $f = \Psi_{L+2} \circ \tilde F$. By Proposition \ref{SingSet}, $\tilde{F}$ has image outside the singularity set $S_{L+2}$. Then Proposition \ref{reg-locus} and Theorem \ref{POP} give a continuous homotopy $h_t: P \times X \to (\mathbb{C}^{\frac{n(n-1)}{2}} \times \mathbb{C}^{\frac{n(n+1)}{2}})^{L+2}\setminus S_{L+2}$ of liftings of $\pi_{2n} \circ f$ along $\Phi_{L+2}$, such that $h_0 = \pi_{2n} \circ \tilde F $ and $h_1$ is $X$-holomorphic.
\end{proof}

\begin{proof}[Proof of Theorem \ref{main-theo-parametric}]
The proof is by induction on $n$. The induction base follows from the factorization for $\mathrm{Sp}_2(\mathbb{C}) = \mathrm{SL}_2(\mathbb{C})$ in \cite{KI}.

 Assume that the theorem is true for $n-1$, we proceed to prove it for $n$. By the solution to the continuous Vaserstein problem, there exist a number $L$ and a continuous lift $F : P \times X \to (\mathbb{C}^{\frac{n(n-1)}{2}} \times \mathbb{C}^{\frac{n(n+1)}{2}})^L$ of $f$ along $\Psi_L$, i.e., $\Psi_L \circ F = f$. Then Proposition \ref{sup-main} yields a continuous homotopy $h_t: P \times X \to (\mathbb{C}^{\frac{n(n-1)}{2}} \times \mathbb{C}^{\frac{n(n+1)}{2}})^{L+2}\setminus S_{L+2}$ of lifts of $\pi_{2n} \circ f$ along $\Phi_{L+2}$, such that $h_0 = \pi_{2n} \circ \tilde F $ and $h_1$ is $X$-holomorphic, where $\tilde{F}$ is as in the proposition. 
 
Since the two symplectic matrices $f(p,x)$ and $\Psi_{L+2}(h_t(p,x))$ have the same last row, we get
\begin{center}
     $\Psi_{L+2}(h_t(p,x))f(p,x)^{-1} = \left[\begin{array}{cccc} A_{t}(p, x)  & 0 & B_{t}(p, x) & b_{1,t}(p,x) \\  a_{t}(p,x) & 1 & b_{2,t}(p,x) & b_{3,t}(p,x) \\
    C_t(p,x) & 0 & D_t(p,x) & d_t(p,x) \\
    0 & 0 & 0 & 1\end{array}\right]$.
\end{center}
where $A_t(p, x), B_t(p, x), C_t(p, x)$ and $D_t(p, x)$ are $(n - 1) \times (n - 1)$ matrices, and the remaining maps are vectors of appropriate dimension. Consider the matrix
\begin{center}
 $\tilde{f}_t(p,x) = \left[\begin{array}{ccc} A_t(p, x) & B_t(p, x) \\  C_t(p, x) & D_t(p, x)\end{array}\right]   \in \mathrm{Sp}_{2n-2}(\mathbb{C})$
\end{center}
At $t=0$ we have from the continuous solution $\Psi_{L+2}(h_0(p, x)) = f(p, x)$, which implies $\tilde{f}_0 = I_{2n-2}$ and thus the holomorphic map
\begin{center}
 $\tilde{f}_1 : P \times X \to \mathrm{Sp}_{2n-2}(\mathbb{C})$  
\end{center}
is null-homotopic. Let $i$ be the canonical inclusion of $ \mathrm{Sp}_{2n-2}(\mathbb{C})$ into $ \mathrm{Sp}_{2n}(\mathbb{C})$ given by
\begin{center}
    $i(\tilde{f}_1(p,x)) = i \left( \left[\begin{array}{ccc} A_t(p, x) & B_t(p, x) \\  C_t(p, x) & D_t(p, x)\end{array}\right] \right) =\left[\begin{array}{cccc} A_{t}(p, x)  & 0 & B_{t}(p, x) & 0 \\  0 & 1 & 0 & 0 \\
    C_t(p,x) & 0 & D_t(p,x) & 0 \\
    0 & 0 & 0 & 1\end{array}\right] $
\end{center}
By the induction hypothesis, we have
\begin{center}
    $i(\tilde{f}_1(p,x)^{-1})=\left[\begin{array}{cccc} \tilde{A}(p, x)  & 0 & \tilde{B}(p, x) & 0 \\  0 & 1 & 0 & 0 \\
    \tilde{C}(p,x) & 0 & \tilde{D}(p,x) & 0 \\
    0 & 0 & 0 & 1\end{array}\right] $
\end{center}
is a finite product of holomorphic elementary symplectic matrices. \\
\\
Then the matrix $\Psi_{L+2}(h_1(p, x))f(p, x)^{-1}i(\tilde{f}_1(p, x)^{-1})$ 
\begin{center}
   $  \left[\begin{array}{cccc} I_{n-1}  & 0 & 0 & b_{1,1}(p,x) \\  a_1(p,x)\tilde{A}(p,x)+b_{2,1}(p,x)\tilde{C}(p,x) & 1 & a_1(p,x)\tilde{B}(p,x)+b_{2,1}(p,x)\tilde{D}(p,x) & b_{3,1}(p,x) \\
   0 & 0 & I_{n-1} & d_1(p,x) \\
    0 & 0 & 0 & 1\end{array}\right]$
\end{center}
is an elementary symplectic matrix, where
\begin{center}
    $a_1(p,x)\tilde{A}(p,x)+b_{2,1}(p,x)\tilde{C}(p,x) = -d_1(p,x)^T$
\end{center}
and
\begin{center}
    $a_1(p,x)\tilde{B}(p,x)+b_{2,1}(p,x)\tilde{D}(p,x) = b_{1,1}(p,x)^T$.
\end{center}
% From this we have:
% \begin{center}
%     $\Psi_L(F_1(p, x))f(p, x)^{-1}i(\tilde{f}(p, x)^{-1})\left[\begin{array}{ccc} I_n & -B(p, x) \\  0 & I_n\end{array}\right] =  \left[\begin{array}{cccc} I_{n-1}  & 0 & 0 & 0 \\ -d_1(p,x)^T  & 1 & 0 & 0 \\
%    0 & 0 & I_{n-1} & d_1(p,x) \\
%     0 & 0 & 0 & 1\end{array}\right]$
% \end{center}
% with 
% \begin{center}
%     $B(p,x) =  \left[\begin{array}{ccc} 0  & -b_{1,1}(p, x) \\  -b_{1,1}(p,x)^T & -b_{3,1}(p,x)-d_{1}(p,x)^Tb_{1,1}(p,x)\end{array}\right]$
% \end{center}
% By the holomorphic factorization in \cite{KI}, we have: 
% \begin{center}
%     $ \left[\begin{array}{cccc} I_{n-1}  & 0 & 0 & 0 \\ -d_1(p,x)^T  & 1 & 0 & 0 \\
%    0 & 0 & I_{n-1} & d_1(p,x) \\
%     0 & 0 & 0 & 1\end{array}\right] = \left[\begin{array}{ccc} H_1(p, x) & 0 \\ 0 & H_1(p, x)^{-T}\end{array}\right]\cdots \left[\begin{array}{ccc} H_{K^{\prime}}(p,x) & 0 \\  0 & H_{K^{\prime}}(p,x)^{-T}\end{array}\right] $
% \end{center}
% Notice that by Proposition \ref{fact-1}, the matrices 
% \begin{center}
%     $\left[\begin{array}{ccc} H_i(p, x) & 0 \\ 0 & H_i(p, x)^{-T}\end{array}\right]$
% \end{center}
% are the symplectic elementary matrices, and thus:
% \begin{center}
%    $ \left[\begin{array}{ccc} H_1(p, x) & 0 \\ 0 & H_1(p, x)^{-T}\end{array}\right]\cdots \left[\begin{array}{ccc} H_{K^{\prime}}(p,x) & 0 \\  0 & H_{K^{\prime}}(p,x)^{-T}\end{array}\right]\left[\begin{array}{ccc} I_n & B(p, x) \\  0 & I_n\end{array}\right]  $
% \end{center}
% is a factorization as in Theorem \ref{main-theo-parametric}. 
This implies $f(p, x)$ has a finite holomorphic symplectic factorization as required. 
\end{proof}
Examining the proof we observe the following.
\begin{rema}
  If $K_{cont}(n,d)$ is such that any holomorphic 
  symplectic matrix of size $2n$ can be factorized by $K_{cont}(n,d)$ continuous elementary symplectic matrices, then we have
  \begin{align}
      K(n,d) \le K_{cont}(n,d) + K(n-1,d) + 3,
    \end{align}   
   Since by Tavgen reduction \cite{HKS}
   \begin{align}
   \quad K_{cont} (n,d) \le K_{cont}(2,d). 
\end{align}
  we deduce $K(n,d) \le (n-1) (K_{cont}(2,d) + 3)$.
\end{rema}

%In this case we can see that the Jacobian is definitely trivial.
%We can split it into two parts:
%\begin{center}
%     $\left[\begin{array}{ccc} 0 \\ -a_{12,3}.z_{2,2} \\ -a_{12,2}z_{2,1}-a_{12,3}z_{3,1} \\ 0 \\ a_{21,3}(z_{1,2}z_{2,1}+z_{22,2}z_{3,1}+z_{3,2}z_{33,1}) \\ 0\end{array}\right] + \left[\begin{array}{ccc} P_1(Z) \\ P_{2}(Z) \\ P_3(Z) \\ z_{11,2}z_{2,1}+z_{1,2}z_{3,1}+z_{2,2}z_{33,1} \\ z_{11,2}z_{2,1}+z_{1,2}z_{3,1}+z_{2,2}z_{33,1} \\ z_{2,2}z_{2,1}+z_{3,2}z_{3,1}+z_{33,2}z_{33,1}+ 1 \end{array}\right]$.
%\end{center}
%From this, we get the Jacobian:
%\begin{center}
%    $[A_K|Z_K]$
%\end{center}
%with $A_K$ ís the Jacobian with partial derivative on $a_{ij,k}$ variables, $Z_K$ ís the Jacobian with partial derivative on $a_{ij,k}$ variables. $Z_K$ is the sum of two the Schott's Jacobian matrix $Z^c_K$ and the jacobian of matrix of the form $Q(A).P(Z)$. Notice that $Z^c_K$ has rank $2n$, i.e., $Z_K = Z^c_K + D$ with $D$ with entries $0$ or of the form $Q(A).P(Z)$. At the entries of where they have the sum of $Q(A).P(Z)$ with entries in $Z^c_K$, the Jacobian at $A_K$ will be nontrivial and it will ensure that  $[A_K|Z_K]$ is a full rank matrix. Thus at $A_K \times \mathbb{C}^{\frac{n(n+1)}{2}} \backslash S_K$ is still a submersion

%If we consider points on $A_K \times Z_K$, we can see that the projection to the last row will have the form:
%\begin{center}
%    $\left[\begin{array}{ccc} P_1(Z) \\ P_{2}(Z) \\ P_3(Z) \\ 0 \\ 0 \\ 1 \end{array}\right]$. 
%\end{center}
%This is the last row of Schott's factorization $Z^c_K$. Thus this implies that $[A_K|Z_K]$ is not a full rank matrix as $A_K = 0$ and $Z_K$ will not be a full rank matrix.  

\section{Stratified ellipticity of $\Phi_K$}\label{sec-stra}

\begin{comment}
    

\subsection{Holomorphic submersive of $\Phi_K$}


%Recall that our elementary matrices have the form below: 
%    \begin{center}
%       $M_K(A,C) = \left[\begin{array}{ccc} A^{-T} & 0 \\ Z & A\end{array}\right] = \left[\begin{array}{ccc} I & 0 \\ C & I\end{array}\right]\left[\begin{array}{ccc} A^{-T} & 0 \\ 0 & A\end{array}\right]$  with $C$ symmetric and $K$ odd
%    \end{center}
%    and
%    \begin{center}
%    $M_K(A,B)= \left[\begin{array}{ccc} A^{-1} & Z \\ 0 & A^{T}\end{array}\right] =\left[\begin{array}{ccc} A^{-1} & 0 \\ 0 & A^T\end{array}\right]\left[\begin{array}{ccc} I & B \\ 0 & I\end{array}\right]  $ with $B$ symmetric and $K$ even. 
%\end{center}
In this section, in order to have a good representation for the proof of submersion and singular locus, we look at the tranpose of the elementary matrices. With suitable coordinates, the elementary matrices can be rewritten to have the form:
 \begin{center}
       $M_K(A_K,Z_K) = \left[\begin{array}{ccc} A^{-1}_K & 0 \\ 0 & A^T_K\end{array}\right] \left[\begin{array}{ccc} I & Z_K \\ 0 & I\end{array}\right] =  \left[\begin{array}{ccc} A^{-1}_K & A^{-1}_K.Z_K \\ 0 & A^T_K\end{array}\right]$, $K$ odd  
    \end{center}
    and
    \begin{center}
    $M_K(A_K,Z_K)=\left[\begin{array}{ccc} I & 0 \\ Z_K & I \end{array}\right] \left[\begin{array}{ccc} A_K^{-T} & 0 \\ 0 & A_K \end{array}\right] =  \left[\begin{array}{ccc} A_K^{-T} & 0 \\ Z_K.A^{-T}_K & A_K \end{array}\right]$, $K$ even
\end{center}
with
\[
A_{K}
=
\begin{bmatrix}
1      & a_{1,2,K}     & \cdots & a_{1,n,K}     \\[6pt]
0      & 1             & \ddots & \vdots     \\[6pt]
\vdots & \ddots        & \ddots & a_{n-1,n,K}   \\[6pt]
0      & \cdots        & 0      & 1
\end{bmatrix}
\]
An upper triangular matrix:
\begin{center}
    $A^{-1}_K = \left[\begin{array}{cccc} 1 & P_{12,K} & \cdots &P_{1n,K} \\  0 & 1 & \ddots & \vdots \\
    \vdots & \ddots & \ddots & P_{(n-1)n,K} \\
    0 & \cdots & 0 & 1\end{array}\right]$
\end{center} with $P_{hn,j}$ are polynomials with $a_{hn,j}$ variables and $Z_K$ a symmetric matrix. 

%We denote: $(\Vec{A}_K,\Vec{Z}_K )$ for the tube $((A_1,Z_1),(A_2,Z_2),\cdots,(A_K,Z_K) \in (\mathbb{C}^{\frac{n(n-1)}{2}} \times \mathbb{C}^{\frac{n(n+1)}{2}})^K$. Here $\Vec{Z}_K$ is $(Z_1,Z_2,\cdots,Z_K) \in (\mathbb{C}^{\frac{n(n+1)}{2}})^K$ and $\Vec{A}_K$ is $(A_1,A_2,\cdots,A_K) \in (\mathbb{C}^{\frac{n(n-1)}{2}})^K$.\\
Recall again that:
\begin{center}
 $\Phi_K = \pi_{2n} \circ \Psi_K : (\mathbb{C}^{\frac{n(n-1)}{2}} \times \mathbb{C}^{\frac{n(n+1)}{2}} )^K \to \mathbb{C}^{2n}\backslash \lbrace 0 \rbrace$,   
\end{center} where $\pi_{2n} : \mathbb{C}^{2n\times 2n} \to \mathbb{C}^{2n}$
is the projection of a $2n \times 2n$ - matrix to its last row. Then after taking transpose, $\Phi_K(\Vec{A}_K,\Vec{Z}_K)$ becomes the projection to the last column with the formula:
\begin{center}
     $\Phi_K(\Vec{A}_K,\Vec{Z}_K)  =M_K(A_K, Z_{K}).\cdots. M_1(A_1,Z_1).e_{2n}$.
\end{center}
From this, we also have a recursive formula: 
\begin{equation}\label{recur-lemma}
 \Phi_K(\Vec{A}_K,\Vec{Z}_K) =M_K(A_K, Z_{K})\Phi_{K-1}(\Vec{A}_{K-1},\Vec{Z}_{K-1}).
\end{equation}

and
\begin{center}
    $S_K = \lbrace (\Vec{A}_K,\Vec{Z}_K)|Z_{i}.e_n = 0, $ $i<K $; $ a_{1n,j} =  \cdots = a_{(n-1)n,j} = 0, 1 < j < K \rbrace$.
\end{center}
%with $\mathcal{W}_K (\Vec{Z}_K) = [Z_2|Z_4|\cdots|Z_{2k}]$  

Equivalently, in the singular set the last columns of all matrices $Z_i$, except $Z_K$, should be 0. Also, the last columns of all $A_i$, except $A_1$ and $A_K$, should be $e_n$. We choose the notation $Z_i.e_n = 0$ just to fit with some computations in the proof of theorem \ref{sing-locus}. We will also sometimes use the term $A_j.e_n \neq e_n$ to say that the last column of $A_j$ is different from $e_n$, that is, there exists $1 \leq h < n$ such that $a_{hn,j} \neq 0$. 
\\
\\
In this subsection we will prove the following theorem:

\begin{theo}\label{sing-locus}
The singular locus, i.e., the non-submersive locus of $\Phi_K$ is $S_K$.

\end{theo}
\begin{rema}
The singular set does not depend on the choice of coordinates. Indeed, the singular set $S_K$ above tells us that for the matrix of the form :
\begin{center}
    $ \left[\begin{array}{ccc} A^{-1}_k & B_k \\ 0 & A^T_k\end{array}\right]$,
\end{center}
or 
\begin{center}
   $ \left[\begin{array}{ccc} A_k^{-T} & 0 \\ C_k & A_k \end{array}\right]$,
\end{center}
$B_k$ and $C_k$ must both have trivial last row and column when $k < K$.
\end{rema}
The ultimate goal is to study the Jacobian of the map $\Phi_K$. Before that, we would recall some tools from \cite{Scho}: 

Let $E_{ij}$ be the $n \times n$ matrix having a $1$ at the entry $(i, j)$ and being zero elsewhere. Then $\Tilde{E}_{ij} =
\frac{1}{1+\delta_{ij}}
(E_{ij} + E_{ji})$ is an elementary symmetric matrix. We will identify $\mathrm{Sym}_n(\mathbb{C})$ and $\mathbb{C}
^{\frac{n(n+1)}{
2}}$ by building a bijection between the
set $\lbrace \Tilde{E}_{ij} : 1 \leq i \leq j \leq n \rbrace$ (which forms the basis of $\mathrm{Sym}_n(\mathbb{C})$). Since the sets $I := \lbrace i \in \mathbb{N} : 1 \leq i \leq
\frac{n(n+1)}{2}
\rbrace$
and $J := \lbrace (i, j) \in \mathbb{N}^2
: 1 \leq i \leq j \leq n \rbrace$ have the same order, we have a bijection $\alpha : I \to J$, which induces an isomorphism $S: \mathbb{C}
^{\frac{n(n+1)}{
2}} \to  \mathrm{Sym}_n(\mathbb{C})$ by defining $S(e_i) = \Tilde{E}_{\alpha(i)}$.

We will show the above theorem by explicitly computing the Jacobian of $\Phi_K$, denoted by $J\Phi_K$. To begin with, we study the Jacobian of elementary matrices. For a fixed $1 \leq i \leq n$, let us define
the mapping $F: \mathbb{C}^n \to \mathbb{C}^{n\times n(n+1)}$ by:
\begin{center}
     $u \mapsto [D_{A_K}(u)|D_{Z_K}(u)]$
\end{center}
with 
\begin{center}
 $D_{A}(u) =  [ E_{1,2}u  \cdots E_{i,j}u \cdots E_{(n-1)n}u| 0_{n \times n}]$ if $A_K$ is upper triangular,    
\end{center}
or 
\begin{center}
 $D_{A}(u) =  [ E_{2,1}u  \cdots E_{j,i}u \cdots E_{n,(n-1)}u | 0_{n \times n}]$ if $A_K$ is lower triangular for $i<j$   
\end{center} and $D_{Z_K}(u) = [\Tilde{E}_{\alpha(1)}u \cdots \Tilde{E}_{\alpha(\frac{n(n+1)}{2})}u]$. For $A^{-1}_K$, we also define $D_{A^{-1}_K}(u)$ as follows:\\
For each matrix $A^{-1} \in \mathrm{SL}_n(\mathbb{C})$, we have the following: 
\begin{center}
    $A^{-1} = \sum_{i<j} P_{ij}(A)E_{ij}$.
\end{center}
Then the derivative 
\begin{center}
    $\frac{\partial}{\partial a_{ij}} A^{-1} = \sum \frac{\partial P_{kh}(A)}{\partial a_{ij}}. E_{kh}$
\end{center}
is an $n \times n$ matrix. From this we define:
\begin{center}
  $D_{A^{-1}_K}(u) := [\frac{\partial}{\partial a_{12}} A^{-1}.u \cdots \frac{\partial}{\partial a_{(n-1)n}} A^{-1}.u|0_n]$  
\end{center}
Here we we only put the partial derivative of $a_{ij}$ with $i < j$. Lower triagular case can be defined the same.
%\subsubsection{Strategy for the proof of submersive}
%In this small section we use some examples to show how the proof works. We first consider the case $K=2$, $N=2$, then $K=3$, $N=2$ to see the potential singular set. \\
%Consider the map:
%    \begin{center}
%        $\Phi_2: \mathbb{C}^{8} \to \mathbb{C}^{4} \backslash \lbrace 0 \rbrace$
%    \end{center}
%as:
%\begin{center}
%    $M_2(Z_2,A_2) M_1(Z_1,A_1).e_{4} = M_2(Z_2,A_2) M_1(Z_1,A_1).\left[\begin{array}{ccc} 0 \\ 0 \\ 0 \\ 1 \end{array}\right]$
%\end{center}
%with 
%\begin{center}
% $ M_{1}(A_{1},Z_{1}) = \left[\begin{array}{ccc} A^{-1}_1 & 0 \\ 0 & A^T_1\end{array}\right] \left[\begin{array}{ccc} I & Z_1 \\ 0 & I\end{array}\right] = \left[\begin{array}{ccccccc} 
%    1 & -a_1  & z_{1,1}-a_1z_{1,2} & z_{1,2}-a_1z_{1,3}  \\ 
%    0 & 1  & z_{1,2} & z_{1,3}  \\
%    0 & 0  & 1 & 0  \\
%    0 & 0  & a_1 & 1 \end{array}\right]$   
%\end{center}
%and 
%\begin{center}
%    $M_2(A_2,Z_2)=\left[\begin{array}{ccc} I & 0 \\ Z_2 & I \end{array}\right] \left[\begin{array}{ccc} A_2^{-T} & 0 \\ 0 & A_2 \end{array}\right] =  \left[\begin{array}{ccc} A_2^{-T} & 0 \\ Z_2.A^{-T}_2 & A_2 \end{array}\right] = \left[\begin{array}{ccccccc} 
%    1 & 0  & 0 & 0 \\ 
%    -a_2 & 1  & 0 & 0  \\
%    z_{2,1}-a_2z_{2,2} & z_{2,2}  & 1 & a_2  \\
%    z_{2,2} - a_2z_{2,3} & z_{2,3}  & 0 & 1 \end{array}\right]$
%\end{center}
%Here 
%\begin{center}
%    $Z_k=  \left[\begin{array}{ccc} z_{k,1} & z_{k,2} \\  z_{k,2} & z_{k,3}\end{array}\right]$.
%\end{center}
%To compute $J\Phi_2$, we need to verify:
%\begin{center}
%    $\frac{\partial}{\partial z_{j,k}} \Phi_2 = \frac{\partial}{\partial z_{j,k}} M_2(Z_2,A_2) M_1(Z_1,A_1).e_{4} $
%\end{center}
%and 
%\begin{center}
%    $\frac{\partial}{\partial a_{k}} \Phi_2 = \frac{\partial}{\partial a_{k}} M_2(Z_2,A_2) M_1(Z_1,A_1).e_{4} $
%\end{center}
%Each are vectors in $\mathbb{C}^4$ and form the Jacobian matrix. We will just do some computations for $z_{2,2}$, $z_{1,2}$, $a_1$ and $a_2$. \\
%\\
%For $z_{1,2}$, we get: 
%\begin{center}
%    $\frac{\partial}{\partial z_{1,2}} \Phi_2 = \frac{\partial}{\partial z_{1,2}} M_2(Z_2,A_2) M_1(Z_1,A_1).e_{4} = M_2(Z_2,A_2).\frac{\partial}{\partial z_{1,1}}(M_1(Z_1,A_1)).e_{4}   $
%\end{center}
%Notice that 
%\begin{center}
%    $\frac{\partial}{\partial z_{1,2}}M_1(Z_1,A_1)=. \left[\begin{array}{ccc} A^{-1}_1 & 0 \\ 0 & A^T_1\end{array}\right] \left[\begin{array}{ccc} I & Z_1 \\ 0 & I\end{array}\right] = \left[\begin{array}{ccccccc} 
%    0 & 0  & -a_1 & 1  \\ 
%    0 & 0  & 1 & 0 \\
%    0 & 0  & 0 & 0  \\
%    0 & 0  & 0 & 0 \end{array}\right]$
%\end{center}
%and we get:
%\begin{center}
%    $\frac{\partial}{\partial z_{1,2}} \Phi_2 = %\left[\begin{array}{ccc} 1 \\ -a_2 \\ z_{2,1}-a_2z_{2,2} \\ 
% z_{2,2}-a_2.z_{2,3} \end{array}\right]$.
%\end{center}
%For $a_1$, we get:
%\begin{center}
%    $\frac{\partial}{\partial a_{1}}M_1(Z_1,A_1)=  \left[\begin{array}{ccccccc} 
%    0 & -1  & -z_{1,2} & -z_{1,3}  \\ 
%    0 & 0  & 0 & 0 \\
%    0 & 0  & 0 & 0  \\
%    0 & 0  & 1 & 0 \end{array}\right]$
%\end{center}
%and 
%\begin{center}
%    $\frac{\partial}{\partial a_{1}} \Phi_2 = 
% \left[\begin{array}{ccc} -z_{1,3} \\ a_2z_{1,3} \\ -
% z_{1,3}z_{2,1}+a_2z_{1,3}z_{2,2} \\ -
% z_{1,3}z_{2,2}+a_2z_{1,3}z_{2,3} \end{array}\right]$.
%\end{center}
%For $z_{2,2}$, we get:
%\begin{center}
%    $\frac{\partial}{\partial z_{2,2}}M_2(Z_2,A_2)=  \left[\begin{array}{ccccccc} 
%    0 & 0  & 0 & 0  \\ 
%    0 & 0  & 0 & 0 \\
%    -a_2 & 1  & 0 & 0  \\
%    1 & 0  & 0 & 0 \end{array}\right]$
%\end{center}
%and we take the product of  
%\begin{center}
    

%$\frac{\partial}{\partial z_{2,2}}M_2(Z_2,A_2)$ and $\Phi_1 = \left[\begin{array}{ccc} z_{1,2}-a_1z_{1,3} \\ z_{1,3} \\ 0 \\ 1 \end{array}\right]$
%\end{center}
%to get:
%\begin{center}
%    $\frac{\partial}{\partial z_{2,2}} \Phi_2 = \left[\begin{array}{ccc} 0 \\ 0\\ -a_2z_{1,2}+a_2a_1z_{1,3}+z_{1,3} \\ z_{1,2}-a_1z_{1,3} \end{array}\right]$.
%\end{center}
%For $a_2$, we get:
%\begin{center}
%    $\frac{\partial}{\partial a_{2}}M_2(Z_2,A_2)=  \left[\begin{array}{ccccccc} 
%    0 & 0  & 0 & 0 \\ 
%    -1 & 0  & 0 & 0 \\
%    -z_{2,2} & 0 & 0 & 1  \\
%    -z_{2,3} & 0 & 0 & 0 \end{array}\right]$
%\end{center}
%and 
%\begin{center}
%    $\frac{\partial}{\partial a_{2}} \Phi_2 = \left[\begin{array}{ccc} 0 \\ -z_{1,2}+a_1z_{1,3} \\ -z_{2,2}z_{1,2}+a_1z_{1,3}z_{2,2}+1 \\ z_{1,2}-a_1z_{1,3} \end{array}\right]$.
%\end{center}
%By the same way of computation, we also get: 
%\begin{center}
%     $\frac{\partial}{\partial z_{1,3}} \Phi_2 = \left[\begin{array}{ccc} -a_1 \\ 1+a_1a_2 \\ -a_1z_{2,1}+(a_1a_2+1)z_{2,2} \\ -a_1z_{2,2} +(a_1a_2+1)z_{2,3}\end{array}\right]$
%\end{center}
%and 
%\begin{center}
%    $\frac{\partial}{\partial z_{2,1}} \Phi_2 = \left[\begin{array}{ccc} 0 \\ 0 \\ z_{1,2}-a_1z_{1,3} \\ 0\end{array}\right]$
%\end{center}
%and 
%\begin{center}
%    $\frac{\partial}{\partial z_{2,3}} \Phi_2 =\left[\begin{array}{ccc} 0 \\ 0 \\ 0 \\  -a_2z_{1,2} +(a_1a_2+1)z_{1,3}\end{array}\right]$.    
%\end{center}
%From all the calculations above, we get $$J \Phi_2$$. It is not hard to see that $J \Phi_2$ is not full rank iff $z_{1,2} = z_{1,3} = 0$.\\
%\\
%Now, we consider the case $K=3$. Now we get: 
%\begin{center}
%    $\Phi_3 = M_3(A_3,Z_3).\Phi_2$.
%\end{center}From the computation method of the case $K=2$, we can split the Jacobian above into two parts. The first part relates to the old variables from $\Phi_2$, another part relates to the new variables from $M_3$. For the first part, we get the matrix:
%\begin{center}
%    $M_3(A_3,Z_3).J\phi_2$
%\end{center}
% denote another part by $\mathcal{A}$. Now, we can see that if $z_{1,2}$ or $z_{1,3}$ are different from $0$, then $M_3(A_3,Z_3).J\phi_2$ is a full rank matrix. So we can assume that we are in the situation $z_{1,2}=z_{1,3} = 0$. Using the same method as in the case $K=2$, we get: 
%\begin{center}
%    $\mathcal{A} = \left[\begin{array}{ccccccc} 
%    -z_{3,2}a_2-z_{3,3} & a_2  & -a_1a_2+1 & -a_3 \\ 
%    0 & 0  & a_2 & 1 \\
%    0 & 0 & 0 & 0  \\
%    a_{2} & 0 & 0 & 0 \end{array}\right] $.
%\end{center}
%From this, we get the Jacobian: 
%\begin{center}
%    $J\Phi_3 = \left[\begin{array}{ccccccc} 
%    *&*& *& -z_{3,2}a_2-z_{3,3} & a_2  & -a_3a_2+1 & -a_3 \\ 
%    *&*& *& 0 & 0  & a_2 & 1 \\
%    *&1& 0& 0 & 0 & 0 & 0  \\
%    *&a_3& 0& a_{2} & 0 & 0 & 0 \end{array}\right] $
%\end{center}
%and by swapping columns, we get:
%\begin{center}
%    $J\Phi_3 = \left[\begin{array}{ccccccc} 
%    *&-a_1& *& a_2 & -a_3a_2+1  & * & -z_{3,2}a_2-z_{3,3} \\ 
%    *&1& *& 0 & a_2  & * & 0 \\
%    *&0& 0& 0 & 0 & 1 & 0  \\
%    *&0& 0& 0 & 0 & a_3 & a_2 \end{array}\right] $.
%\end{center}
%If $a_2 \neq 0$, then $J\Phi_3$ must be full rank. So we might assume that $a_2 = 0$. \\
%Now in this situation we get: 

%    \begin{center}
%    $J\Phi_2= \left[\begin{array}{ccccccc} 
%    1 & -a_1  & 0 & 0 \\ 
%    0 & 1  & 0 & 0 \\
%    z_{2,1} & -a_1z_{2,2}+z_{2,3} & 1 & 0  \\
%    z_{2,2} & -a_1z_{2,2}+z_{2,3} & 0 & 0 \end{array}\right]. $
%\end{center}
%Now denote: 
%\begin{center}
%   $E^*_1 = \left[\begin{array}{ccc} 1 & 0 \\ 0 & 0\end{array}\right] $ 
%\end{center}
%and we can rewrite $J\Phi_2$ as:

%\begin{center}
%    $\left[\begin{array}{ccc} A & 0 \\ B & E^*_1\end{array}\right] $
%\end{center}
%and $J\Phi_3$ as:
%\begin{center}
%     $\left[\begin{array}{ccc} A^{-1}_3A^{-1}_1+A^{-1}_3Z_3B & A^{-1}_3Z_3E^*_1 & A^{-1}_1 \\ A^T_3Z_2A^{-1}_1 & A^T_3E^*_1 & 0\end{array}\right] $
%\end{center}
%We have:
%\begin{center}
%    $A^T_3E^*_1 = \left[\begin{array}{ccc} 1 & 0 \\ a_3 & 1\end{array}\right] \left[\begin{array}{ccc} 1 & 0 \\ 0 & 0\end{array}\right] =\left[\begin{array}{ccc} 1 & 0 \\ a_3 & 0\end{array}\right] $
%\end{center}
%and 
%\begin{center}
%    $A^T_3Z_2A^{-1}_1 = \left[\begin{array}{ccc} 1 & 0 \\ a_3 & 1\end{array}\right] \left[\begin{array}{ccc} z_{2,1} & z_{2,2} \\ z_{2,2} & z_{2,3}\end{array}\right] \left[\begin{array}{ccc} 1 & -a_1 \\ 0 & 1\end{array}\right] =\left[\begin{array}{ccc} 1 & 0 \\ a_3 & 1\end{array}\right] \left[\begin{array}{ccc} z_{2,1} & z_{2,2}-a_1z_{2,1} \\z_{2,2} & z_{2,3}-a_1z_{2,2}\end{array}\right] = \left[\begin{array}{ccc} z_{2,1} & z_{2,2}-a_1z_{2,1} \\z_{2,2}+a_3z_{2,1} &a_3(z_{2,2}-a_1z_{2,1})+ z_{2,3}-a_1z_{2,2}\end{array}\right] $.
%\end{center}
%Then $J\Phi_3$ can be rewritten as:
%\begin{center}
%    $J\Phi_3 =  \left[\begin{array}{ccccccc} 
%    *&*& *& -z_{3,3} & 1  & -a_1 \\ 
%    *&*& *& 0 & 0  & 1 \\
%    z_{2,1}& z_{2,2}-a_1z_{2,1}& 1& 0 & 0 & 0  \\
%    z_{2,2}+a_3z_{2,1}&a_3(z_{2,2}-a_1z_{2,1})+ z_{2,3}-%a_1z_{2,2}& a_3& 0 & 0 & 0  \end{array}\right]$.
%\end{center}
%By reducing the 0-column and applying row operations, we get:
%\begin{center}
%     $J\Phi_3 =  \left[\begin{array}{ccccccc} 
%    *&*& * & 1  & -a_1 \\ 
%    *&*& * & 0  & 1 \\
%    z_{2,1}& z_{2,2}-a_1z_{2,1}& 1 & 0 & 0  \\
%    z_{2,2} & z_{2,3}-a_1z_{2,2}& 0 & 0 & 0  %\end{array}\right]$.
%\end{center}
%Reordering the columns, we have:
%\begin{center}
%     $J\Phi_3 =  \left[\begin{array}{ccccccc} 
%     1  & -a_1 & *  & * & * & -z_{3,3} \\ 
%    0  & 1 & *  &* &* & 0 \\
%     0 & 0 &1  &z_{2,1}& z_{2,2}-a_1z_{2,1} & 0  \\
%     0 & 0 &0 &z_{2,2} & z_{2,3}-a_1z_{2,2}  & %0\end{array}\right]$.
%\end{center}
%It is not hard to see that $J\Phi_3$ is of full rank iff both $z_{2,2}$ and $z_{2,3}$ are not both equal to 0. Thus from the explicit computations for the case $N=2$ and $K=2$ or $K=3$ we see what the singular set $S$ is 
\subsubsection{Proof of theorem \ref{sing-locus}}
Now we can compute the Jacobian:
\begin{lemm}\label{comp-Jac}
    For $u, v \in \mathbb{C}^n$ such that $\left[\begin{array}{ccc} u \\ v \end{array}\right] \in \mathbb{C}^{2n} \backslash \lbrace 0 \rbrace$, the Jacobian of the mapping
    \begin{center}
        $\mathbb{C}^{ \frac{n(n-1)}{2}\times \frac{n(n+1)}{2}} \to \mathbb{C}^{2n} \backslash \lbrace 0 \rbrace, (A_k,Z_k) \mapsto M_k(A_k,Z_k)\left[\begin{array}{ccc} u \\ v \end{array}\right]$
    \end{center}
is given by 
\begin{center}
$\mathcal{A}_k(u,v)=
    \begin{cases}
      \left[\begin{array}{ccc} D_{A^{-1}_k}(u)+D_{A^{-1}_k}(Z_k.v) & A^{-1}_kD_{Z_k}(v) \\  D_{A^{T}_k}(v) & 0\end{array}\right], k=2t+1. \\ \\
      
      \left[\begin{array}{ccc} D_{A^{-T}_k}(u) & 0 \\ D_{A_k}(v)+Z_kD_{A^{-T}_k}(u) & D_{Z_k}(A^{-T}_k.u)  \end{array}\right], k=2t.
    \end{cases}
 $   
\end{center}
\end{lemm}
\begin{proof}
We prove the statement for $k$ odd first. Notice that  $A_k =\sum_{1 \leq i<j \leq n} a_{ij,k}(A)E_{i,j} $, $  A^{-1}_k = \sum_{1 \leq i < j \leq n} P_{ij,k}(A)E_{i,j}$ and $Z_k = \sum_{1\leq i \leq j \leq n} z_{ij}\Tilde{E}_{i,j}$. We then have the partial derivatives:
\begin{align*}
    \frac{\partial}{\partial a_{ij,k}}M_k(A_k,Z_k)  \left[\begin{array}{ccc} u  \\ v \end{array}\right] &= \left[\begin{array}{ccc} \sum_{l<h}\frac{\partial}{\partial a_{ij,k}}P_{lh}(A).E_{l,h} &\sum_{l<h}\frac{\partial}{\partial a_{ij,k}}P_{lh}(A).E_{l,h}.Z_k \\ 0 & E_{j,i} \end{array}\right]\left[\begin{array}{ccc} u  \\ v \end{array}\right] \\
    &= \left[\begin{array}{ccc} \sum_{l<h}\frac{\partial}{\partial a_{ij,k}}P_{lh}(A).E_{l,h}.u + \sum_{l<h}\frac{\partial}{\partial a_{ij,k}}P_{lh}(A).E_{l,h}.(Z_k.v)  \\ E_{j,i}.v \end{array}\right]
\end{align*}
and for $z_{ij}$ we have 
\begin{center}
    $\frac{\partial}{\partial z_{ij,k}}M_k(A_k,Z_k)  \left[\begin{array}{ccc} u  \\ v \end{array}\right] =  \left[\begin{array}{ccc} 0 & A^{-1}_k.\tilde{E}_{i,j} \\ 0 &  0 \end{array}\right]\left[\begin{array}{ccc} u  \\ v \end{array}\right] =  \left[\begin{array}{ccc} A^{-1}_k.\tilde{E}_{i,j}.v \\ 0 \end{array}\right]$
\end{center}
The Jacobian follows from the definition of the maps $D_{Z_k}$ and $D_{A_k}$. We do the same for $k$ even and finish the lemma. 
%\begin{align*}
%    \frac{\partial}{\partial a_{ij,k}}M_k(A_k,Z_k)  \left[\begin{array}{ccc} u  \\ v \end{array}\right] &= \left[\begin{array}{ccc} \sum_{l<h}\frac{\partial}{\partial a_{ij,k}}P_{hl}(A).E_{h,l} &0 \\ Z_{k}(\sum_{l<h}\frac{\partial}{\partial a_{ij,k}}P_{hl}(A).E_{h,l}) & E_{i,j} \end{array}\right]\left[\begin{array}{ccc} u  \\ v \end{array}\right] \\ &= \left[\begin{array}{ccc} \sum_{l<h}\frac{\partial}{\partial a_{ij,k}}P_{hl}(A).E_{h,l}.u   \\ E_{i,j}.v + Z_k(\sum_{l<h}\frac{\partial}{\partial a_{ij,k}}P_{hl}(A).E_{h,l})u\end{array}\right]
%\end{align*}
%and for $z_{ij}$ we have 
%\begin{center}
%    $\frac{\partial}{\partial z_{ij,k}}M_k(A_k,Z_k)  \left[\begin{array}{ccc} u  \\ v \end{array}\right] =  \left[\begin{array}{ccc} 0 & 0 \\ \tilde{E}_{i,j}.A^{-T}_k &  0 \end{array}\right]\left[\begin{array}{ccc} u  \\ v \end{array}\right] =  \left[\begin{array}{ccc} 0 \\ \tilde{E}_{i,j}(A^{-T}_k.u) \end{array}\right]$
%    \end{center}
\end{proof}
\begin{coro}\label{Jacobian}
    The Jacobian $J\Phi_1$ is given by $A_1(e_{2n})$. For $K 
    \geq 2$, the Jacobian $J\Phi_K$ of $\Phi_K$ is given by
    \begin{center}
$J\Phi_K(\Vec{A}_K,\Vec{Z}_K) =
 \left[\begin{array}{ccc} M_K(A_K,Z_K).J\Phi_{K-1}(\Vec{A}_{K-1},\Vec{Z}_{K-1}) & \mathcal{A}_K(\Phi_{K-1}(\Vec{A}_{K-1},\Vec{Z}_{K-1})) \end{array}\right]$.
\end{center}
\end{coro}
From Corollary \ref{Jacobian}, we can see that the situation that makes $J\Phi_K(\Vec{A}_K,\Vec{Z}_K)$ a non-full rank will rely on  $J\Phi_{K-1}(\Vec{A}_{K-1},\Vec{Z}_{K-1})$ and $\mathcal{A}_K(\Phi_{K-1}(\Vec{A}_{K-1},\Vec{Z}_{K-1}))$. 
\begin{lemm}\label{lemm-sing1}
Let $(\Vec{A}_K,\Vec{Z}_K) \in (\mathbb{C}^{\frac{n(n-1)}{2}} \times \mathbb{C}^{\frac{n(n+1)}{2}} )^K$ and assume that there is $1 \leq k \leq[\frac{K-1}{2}]$ such that $Z_{2i-1}e_n = 0$ for all $1 \leq i \leq k$ and $a_{hn,2j} = 0$ for all $h<n$ and $j \leq k$. Then $\Phi_t (\Vec{A}_t,\Vec{Z}_t ) = e_{2n}$ for all $1 \leq t \leq 2k$.
\end{lemm}
\begin{proof}
 We prove the lemma by induction. With $n = 1$ then we have $\Phi_1(\Vec{A}_1,\Vec{Z}_1) = M_1(A_1,Z_1).e_{2n} =  \left[\begin{array}{ccc} Z_1.e_n \\ A^T_1.e_n \end{array}\right] $. Since we assume $a_{hn,2j} = 0$ for all $h<n$ and $j\leq k$, then $A^T_1.e_n = e_n$. Thus $M_1(A_1,Z_1).e_{2n} =  \left[\begin{array}{ccc} 0 \\ e_n \end{array}\right] = e_{2n}$. We assume that the statement is true for $t<2k$. Then we have:
 \begin{center}
     $\Phi_t (\Vec{A}_{t+1},\Vec{Z}_{t+1} ) = M_{t+1} (A_{t+1},Z_{t+1} )\Phi_{t}(\Vec{A}_t,\Vec{Z}_{
t}) = M_{t+1} (A_{t+1},Z_{t+1} )e_{2n}$
 \end{center}
 If $t+1$ is even then it is easy. If $t+1$ is odd then 
 \begin{center}
    $M_{t+1} (A_{t+1},Z_{t+1} )e_{2n} = \left[\begin{array}{ccc} Z_{t+1}.e_n \\ A^T_{t+1}.e_n \end{array}\right]= \left[\begin{array}{ccc} 0 \\ e_n \end{array}\right] = e_{2n}$.
 \end{center}
\end{proof}
We give a simple example to illustrate the computation of the lemma above: 
%\begin{exam}\label{simp-exam1}
%We give a simple example
%\begin{center}
% $\Phi_3 = \pi_{2n} \circ \Psi_3 : (\mathbb{C}^{\frac{n(n-1)}{2}} \times \mathbb{C}^{\frac{n(n+1)}{2}} )^3 \to \mathbb{C}^{2n} \backslash \lbrace 0 \rbrace$,   
%\end{center}    
%with the matrix form:
%\begin{center}
%     $\left[\begin{array}{ccc} A & B \\ C & D\end{array}\right] = \left[\begin{array}{ccc} A^{-T}_1 & 0 \\ Z_1A^{-T}_1 & A_1\end{array}\right] \left[\begin{array}{ccc} A^{-1}_2 & A^{-1}_2Z_2 \\ 0 & A^T_2\end{array}\right] \left[\begin{array}{ccc} A^{-T}_3 & 0 \\ Z_3A^{-T}_3 & A_3\end{array}\right]$
%\end{center}
%We consider the projection to the last row:
%\begin{center}
%    $\left[\begin{array}{ccc} A^{-1}_3 & A^{-1}_3Z_3 \\ 0 & A^T_3\end{array}\right] \left[\begin{array}{ccc} A^{-T}_2 & 0 \\ Z_2A^{-T}_2 & A_2\end{array}\right] \left[\begin{array}{ccc} A^{-1}_1 & A^{-1}_1Z_1 \\ 0 & A^T_1\end{array}\right]e_{2n}$ (*)
%\end{center}
%Suppose 
%\begin{center}
%    $A_k = \left[\begin{array}{ccc}1 & a_{12,k} & a_{13,k} \\ 0 & 1 & a_{23,k} \\ 0 & 0 & 1\end{array}\right]$.
%\end{center}
%For $Z_k$, we define:
%\begin{center}
%    $Z_k = \left[\begin{array}{ccc}z_{11,k} & z_{1,k} & z_{2,k} \\ z_{1,k} & z_{22,k} & z_{3,k} \\ z_{2,k} & z_{3,k} & z_{33,k}\end{array}\right]$.
%\end{center}
%Suppose now $a_{23,k} = a_{13,k} = 0$ for $k$ even and $z_{33,k} = z_{2,k}= z_{3,k} = 0$ for $k$ odd. In this case: 
%\begin{center}
 %   $A^{-1}_k.Z_k = \left[\begin{array}{ccc}1 & a_{12,k} & a_{13,k} \\ 0 & 1 & a_{23,k} \\ 0 & 0 & 1\end{array}\right]. \left[\begin{array}{ccc}z_{11,k} & z_{1,1} & 0 \\ z_{1,1} & z_{22,k} & 0 \\ 0 & 0 & 0\end{array}\right] = \left[\begin{array}{ccc}* & * & 0 \\ * & * & 0 \\ * & * & 0\end{array}\right]$.
%\end{center}Then (*) is equal to:
%\begin{center}
   % $\left[\begin{array}{ccc} A^{-1}_3 & A^{-1}_3Z_3 \\ 0 & A^T_3\end{array}\right] \left[\begin{array}{ccc} A^{-T}_2 & 0 \\ Z_2A^{-T}_2 & A_2\end{array}\right] \left[\begin{array}{ccc} A^{-1}_1 & A^{-1}_1Z_1 \\ 0 & A^T_1\end{array}\right]e_{2n}= \left[\begin{array}{ccc} A^{-1}_3 & A^{-1}_3Z_3 \\ 0 & A^T_3\end{array}\right] \left[\begin{array}{ccc} A^{-T}_2 & 0 \\ Z_2A^{-T}_2 & A_2\end{array}\right]e_{2n}$.
%\end{center}
%Since $a_{23,2}=a_{13,2} = 0$ and $A^{-1}_3.Z_3$ has trivial last column, we have $\Phi_3(\Vec{A}_3,\Vec{Z}_3)$ is $e_{2n}$.
%\end{exam}
Denote: 
\begin{center}
   $E^* =   \left[\begin{array}{ccccccc}  1 & 0 & \cdots &0 &0 & \cdots &0\\ 
  0 & 1 & \cdots & 0 &0 & \cdots &0\\
     0 & 0 & \cdots & 1 &0  & \cdots &0\\
    0 & 0 & \cdots & 0 &0  & \cdots &0\end{array}\right]$   
\end{center} which is an $n \times \frac{n(n+1)}{2}$ matrix. The following lemma is important for us to compute the Jacobian of $\Phi_K$:
\begin{lemm}\label{lemm-sing2}
We have $ D_{A^{-T}}(e_n) = D_{A^T}(e_n) = 0$ and $D_{A}(e_n)= E^*$ .     
\end{lemm}
\begin{proof}
Indeed, we know that:
 \begin{center}
  $D_{A_K}(e_n) =  [ E_{12}e_n  \cdots E_{ij}e_n \cdots E_{(n-1)n}e_n|0_{n}]$ with $i < j$.  
 \end{center}
 Since only
 \begin{center}
  $E_{in}.e_n = \left[\begin{array}{ccc} 0 \\ \cdots \\ 0 \\ 1 \\ \cdots \\ 0 \end{array}\right] \neq 0$ with $1$ at the $i$-th row,  
 \end{center}
 then by the definition of $D_{A_S}$, it must have the form above. 
   
\end{proof}

%Now we can start the proof of theorem \ref{sing-locus}. Our strategy is to begin with proving the submersive of $\Phi_K$ on a big open subset, called $W_K$ (which we will define it in the first step). Then we will prove the submersive on  $ (\mathbb{C}^{\frac{n(n-1)}{2}\times\frac{n(n+1)}{2}})^K \backslash S_K $ and point out that this is the smooth locus of $\Phi_K$.

Now we come to the proof of the theorem \ref{sing-locus}. Here we give a proof by induction:
\begin{proof}[Proof of the theorem \ref{sing-locus}:]
We start with the case $K=2$. Then the jacobian $J\Phi_2$ is:
\begin{center}
    $\left[\begin{array}{ccc} M_2(A_2,Z_2).J\Phi_{1}(\Vec{A}_{1},\Vec{Z}_{1}) & \mathcal{A}_2(\Phi_{1}(\Vec{A}_{1},\Vec{Z}_{1})) \end{array}\right]$.
\end{center}
Notice that $J\Phi_{1}(\Vec{A}_{1},\Vec{Z}_{1}) =\mathcal{A}_1(e_{2n}) $. Notice that by lemma \ref{comp-Jac}, $\mathcal{A}_1(e_{2n})$ has the form :
\begin{center}
    $\left[\begin{array}{ccc} D_{A^{-1}_1}(Z_1.e_n) & A^{-1}_1D_{Z_1}(e_n) \\  D_{A^{T}_1}(e_n) & 0\end{array}\right]$
\end{center}
by lemma \ref{lemm-sing2}, $D_{A^{T}_1}(e_n) = 0$, we then get: 
\begin{center}
    $\left[\begin{array}{ccc} D_{A^{-1}_1}(Z_1.e_n) & A^{-1}_1D_{Z_1}(e_n) \\  0 & 0\end{array}\right]$.
\end{center}
Thus $M_2(A_2,Z_2).J\Phi_{1}(\Vec{A}_{1},\Vec{Z}_{1})$ has the form:
\begin{center}
    $\left[\begin{array}{ccc} A^{-T}_2D_{A^{-1}_1}(Z_1.e_n) & A^{-T}_2A^{-1}_1D_{Z_1}(e_n) \\  Z_2A^{-T}_2D_{A^{-1}_1}(Z_1.e_n) & Z_2A^{-T}_2A^{-1}_1D_{Z_1}(e_n)\end{array}\right]$.
\end{center}
Denote $\Phi_S(\vec{A}_S,\vec{Z}_S)= (P^S_f,P^S_s)$. Now for $\mathcal{A}_2(P^1_f,P^1_s)$ we have:
\begin{center}
    $ \left[\begin{array}{ccc} D_{A^{-T}_2}(P^1_f) & 0 \\ D_{A_2}(P^1_s)+Z_2D_{A^{-T}_2}(P^1_f) & D_{Z_2}(A^{-T}_2.P^1_f)  \end{array}\right]$
\end{center}
It is not hard to see that for $Z_1$, if its last column is different from 0, i.e., $Z_1.e_n \neq 0$, then $P^1_f \neq 0$ and $D_{Z_2}(A^{-T}_2.P^1_f)$ is of full rank. Together with $A_2^{-T}A^{-1}_1D_{Z_1}(e_n)$ which is of full rank, we get $J\Phi_2$ also a matrix with full rank. On the other hand, if $Z_1.e_n = 0$, then $P^1_f = 0_n$ and $\mathcal{A}_2(P^1_f,P^1_s)$ is:
\begin{center}
     $ \left[\begin{array}{ccc} 0 & 0 \\ D_{A_2}(e_n) & 0  \end{array}\right]$
\end{center}
Notice that $D_{A_2}(e_n) = E^*$, so it is not a matrix of full rank. Thus, this implies $J\Phi_2$ not a matrix of full-rank matrix. \\
Now in the next step, we prove the theorem for the case $K=3$.
The jacobian $J\Phi_3$ has the form: 
\begin{center}
     $\left[\begin{array}{ccc} M_3(A_3,Z_3).J\Phi_{2}(\Vec{A}_{2},\Vec{Z}_{2}) & \mathcal{A}_3(\Phi_{2}(\Vec{A}_{2},\Vec{Z}_{2})) \end{array}\right]$.
\end{center}
If $J\Phi_{2}(\Vec{A}_{2},\Vec{Z}_{2})$ is of full rank then $J\Phi_3$ is of full rank, so we assume that we are in the situation where $Z_1.e_n = 0$ (singular locus of $J\Phi_2$). \\
For $\mathcal{A}_3(\Phi_{2}(\Vec{A}_{2},\Vec{Z}_{2}))$, it has the form:
\begin{center}
    $  \left[\begin{array}{ccc} D_{A^{-1}_3}(P^2_f)+D_{A^{-1}_3}(Z_3.P^2_s) & A^{-1}_3D_{Z_3}(P^2_s) \\  D_{A^{T}_3}(P^2_s) & 0\end{array}\right]$.
\end{center}
In case $A_2.e_n \neq e_n$, more precisely: 
\begin{center}
    $\Phi_2(\vec{A}_2,\vec{Z}_2) =  \left[\begin{array}{ccc} 0   \\  A_2.e_n \end{array}\right] $.
\end{center}
Suppose the element $a_{hn,2}$ of the last column of $A_2$ is different from 0. Then $D_{A^T_{3}}(P^2_s)$ has the form: 
 \begin{center}
   $ \left[\begin{array}{ccccccc} 0 & \cdots  & 0 & 0 & 0 & \cdots &0 \\ 
    0 & \cdots & * & * & 0 & \cdots & * \\
    0 & \cdots & * & * & 0 & \cdots & * \\
    0 & \cdots & * & * & a_{hn,2} & \cdots & *\end{array}\right]$. 
 \end{center} 
 Indeed, consider the product
 \begin{center}
    $v_{nh}=E_{nh}.\left[\begin{array}{ccc} 0 \\ \cdots \\ 0 \\ a_{hn,2} \\ \cdots \\ 1 \end{array}\right] = \left[\begin{array}{ccccccc} 0 & \cdots  & 0 & 0 & 0 & \cdots &0 \\ 
    0 & \cdots & 0 & 0 & 0 & \cdots & 0 \\
    0 & \cdots & 0 & 0 & 0 & \cdots & 0 \\
    \cdots & \cdots & \cdots & \cdots & \cdots & \cdots & \cdots \\
    0 & \cdots & 0 & 1 & 0 & \cdots & 0\end{array}\right].\left[\begin{array}{ccc} 0 \\ \cdots \\ 0 \\ a_{hn,S} \\ \cdots \\ 1 \end{array}\right] = \left[\begin{array}{ccc} 0 \\ \cdots \\ 0 \\ 0 \\ \cdots \\ a_{hn,2} \end{array}\right]$.
 \end{center}
We know that $J\Phi_3$ has the form: 
\begin{center}
    $\left[\begin{array}{cccc}  *&*& D_{A^{-1}_{3}}(Z_{3}.P^{2}_s)  & A^{-1}_{3}D_{Z_{3}}(P^{2}_s) \\  *& A^T_{3}D_{A_2}(e_n)& D_{A^T_{3}}(P^{2}_s) & 0 \end{array}\right]$. 
\end{center}
Since $D_{A_2}(e_n) = E^*$, the matrix $A^T_{3}D_{A_2}(e_n)$ is the lower triangular matrix with the 0 vector as the last column. By swapping this 0 column with $v_{nh}$, we get a rank n matrix. Combined with $A^{-1}_3D_{Z_3}(P^2_S)$, we get a full-rank matrix. \\
Thus, we can suppose that $A_{2}.e_n = e_n$ and assume that $Z_2.e_n \neq 0$. Then $J\Phi_3$ has the form: 
\begin{center}
    
    $\left[\begin{array}{cccc}  *&*& D_{A^{-1}_{3}}(Z_{3}.e_n)  & A^{-1}_{3}D_{Z_{3}}(e_n) \\  A^T_3.Z_2A^T_2A^{-1}_1D_{Z_1}(e_n)& A^T_{3}D_{A_2}(e_n)& 0 & 0 \end{array}\right]$. 

\end{center}
Consider the product $\left[\begin{array}{cccc}  I_n&0 \\  0& A^{-T}_{3}\end{array}\right].J\Phi_3$, it is:
\begin{center}
    $J = \left[\begin{array}{cccc}  *&*& D_{A^{-1}_{3}}(Z_{3}.e_n)  & A^{-1}_{3}D_{Z_{3}}(e_n) \\  Z_2A^T_2A^{-1}_1D_{Z_1}(e_n)& D_{A_2}(e_n)& 0 & 0 \end{array}\right]$.
\end{center}
since  $\left[\begin{array}{cccc}  I_n&0 \\  0& A^{-T}_{3}\end{array}\right]$ is of full rank, the matrix $J$ must have the same rank with $J\Phi_3$. Since $A^{-1}_3D_{Z_3}(e_n)$ is of full rank, we will focus on $Z_2A^T_2A^{-1}_1D_{Z_1}(e_n)$ and $D_{A_2}(e_n)$. Since $D_{A_2}(e_n)$ is $E^*$, the rank of $J$ based on the last row of $Z_2$, i.e., $\rank J < 2n $ iff $Z_2.e_n = 0$. Thus we finish the case $K=3$.\\
Now we assume that the theorem is true for $K-1$, and we will prove that it is also true for $K$. First assuming that the theorem is true for $K-1$ an odd number, we will prove that it is also true for $K$ an even number.\\
We also assume that we are in the singular locus of $(\vec{A}_{K-1},\vec{Z}_{K-1})$. The jacobian $J\Phi_K$ has the form: 

    \begin{center}
     $\left[\begin{array}{ccc} M_K(A_K,Z_K).J\Phi_{K-1}(\Vec{A}_{K-1},\Vec{Z}_{K-1}) & \mathcal{A}_K(\Phi_{K-1}(\Vec{A}_{K-1},\Vec{Z}_{K-1})) \end{array}\right]$.
\end{center}
For $\mathcal{A}_K(\Phi_{K-1}(\Vec{A}_{K-1},\Vec{Z}_{K-1}))$, it has the form:
\begin{center}
    $ \left[\begin{array}{ccc} D_{A^{-T}_k}(P^{K-1}_f) & 0 \\ D_{A_k}(P^{K-1}_s)+Z_kD_{A^{-T}_k}(P^{K-1}_f) & D_{Z_k}(A^{-T}_k.P^{K-1}_f)  \end{array}\right] $
\end{center}
Assume $Z_{K-1}.e_n \neq 0$,.i.e, $P^{K-1}_f \neq 0$ and $P^{K-1}_s = e_n$. This implies $D_{Z_k}(A^{-T}_k.P^{K-1}_f)$ a full-rank matrix. \\
Notice that by Lemma \ref{comp-Jac}:
\begin{center}
    $\mathcal{A}_{K-1}(\Phi_{K-2}(\Vec{A}_{K-2},\Vec{Z}_{K-2})) =  \mathcal{A}_{K-1}(P^{K-2}_f,P^{K-2}_s) =  \left[\begin{array}{ccc} * & A^{-1}_{K-1}D_{Z_{K-1}}(P^{K-2}_s) \\  D_{A^{T}_{K-1}}(P^{K-2}_s) & 0\end{array}\right].$
\end{center}
This implies
\begin{center}
    $J\Phi_{K-1}(\Vec{A}_{K-1},\Vec{Z}_{K-1}) = \left[\begin{array}{ccc}  *& A^{-1}_{K-1}D_{Z_{K-1}}(P^{K-2}_s) \\  *& 0 \end{array}\right]$.
\end{center}
Then for $M_{K}(A_{K},Z_{K}).J\Phi_{K-1}(\Vec{A}_{K-1},\Vec{Z}_{K-1})$, it has the form:
\begin{center}
    $ \left[\begin{array}{ccc} A_{K}^{-T} & 0 \\ Z_{K}.A^{-T}_{K} & A_{K} \end{array}\right].  \left[\begin{array}{ccc}  *& A^{-1}_{K-1}D_{Z_{K-1}}(P^{K-2}_s) \\  *& 0 \end{array}\right] =  \left[\begin{array}{ccc}  *&  A_{K}^{-T}A^{-1}_{K-1}D_{Z_{K-1}}(P^{K-2}_s) \\  *&   *\end{array}\right]$
\end{center}

Hence $\left[\begin{array}{ccc} M_K(A_K,Z_K).J\Phi_{K-1}(\Vec{A}_{K-1},\Vec{Z}_{K-1}) & \mathcal{A}_K(\Phi_{K-1}(\Vec{A}_{K-1},\Vec{Z}_{K-1})) \end{array}\right]$ has the form :
\begin{center}
    $ \left[\begin{array}{cccc} * & A^{-T}_{K}A^{-1}_{K-1}D_{Z_{K-1}}(P^{K-2}_s) & * & 0  \\ * & * & * & D_{Z_{K}}(A^{-T}_{K}.P^{K-1}_f)\end{array}\right]$.
\end{center}
Since $P^{K-2}_s = e_n$, we must have $A^{-T}_{K}A^{-1}_{K-1}D_{Z_{K-1}}(P^{K-2}_s)$ a full-rank matrix. Thus, this implies $J\Phi_K$ a full-rank matrix.\\
From now on we can assume that $Z_{K-1}.e_n = 0$. Assume that $A_{K-1}e_n \neq e_n$. Then we have the following:
\begin{center}
   $\mathcal{A}_{K-2}(\Phi_{K-3}(\Vec{A}_{K-3},\Vec{Z}_{K-3}))= \mathcal{A}_{K-2}(e_{2n}) = \left[\begin{array}{ccc}0 & 0 \\ E^* & 0  \end{array}\right]$ 
\end{center} and thus  \begin{align*}
      M_{K-1}(A_{K-1},Z_{K-1}).J\Phi_{K-2}(\Vec{A}_{K-2},\Vec{Z}_{K-2}) =  & \left[\begin{array}{ccc} A^{-1}_{K-1} & A^{-1}_{K-1}.Z_{K-1} \\ 0 & A^T_{K-1}\end{array}\right]\left[\begin{array}{cccc} * &0 & 0 \\ * & D_{A_{K-2}}(e_n) & 0  \end{array}\right]\\
=&\left[\begin{array}{cccc} * &A^{-1}_{K-1}.Z_{K-1}.E^* & 0 \\ * & A^T_{K-1}E^* & 0  \end{array}\right].
 \end{align*}
 Also, we have: 
 \begin{center}
     $\mathcal{A}_{K-1}(\Phi_{K-2}(\Vec{A}_{K-2},\Vec{Z}_{K-2})) =  \mathcal{A}_{K-1}(P^{K-2}_f,P^{K-2}_s) =  \left[\begin{array}{ccc} * & A^{-1}_{K-1}D_{Z_{K-1}}(e_n) \\  D_{A^{T}_{K-1}}(e_n) & 0\end{array}\right].$
 \end{center}
 This implies $J\Phi_{K-1}$ has the form: 
 \begin{center}
      $\left[\begin{array}{cccc} * &A^{-1}_{K-1}.Z_{K-1}.E^* & A^{-1}_{K-1}D_{Z_{K-1}}(e_n)  \\ * & A^T_{K-1}E^* & 0  \end{array}\right]$.
 \end{center}
 Now $J\Phi_{K}(\Vec{A}_{K},\Vec{Z}_{K})$ has the form: 
 
\begin{center}
$\left[\begin{array}{cccc} * &A^{-T}_{K}A^{-1}_{K-1}.Z_{K-1}.E^* & A^{-T}_{K}A^{-1}_{K-1}D_{Z_{K-1}}(e_n) & 0 \\ * & Z_K.A^{-T}_{K}A^{-1}_{K-1}.Z_{K-1}.E^*+A_{K}A^T_{K-1}E^* & *  & E^*  \end{array}\right]$ 
\end{center} By rows and columns operations, the Jacobian matrix is equivalent to
\begin{center}
  $\left[\begin{array}{cccc} * &A^{-T}_{K}A^{-1}_{K-1}.Z_{K-1}.E^* & A^{-T}_{K}A^{-1}_{K-1}D_{Z_{K-1}}(e_n) & 0 \\ * & A_{K}A^T_{K-1}E^* & *  & E^*  \end{array}\right]$  
\end{center} Since we assume $A_{K-1}$ is a matrix with $a_{hn,K-1} \neq 0$ for some $h < n$, the matrix $A_{K}A^T_{K-1}E^*$ will have a column at position $h$ such that the last element is nonzero. Thus, the matrix: 
\begin{center}
    $\left[\begin{array}{cccc} A_{K}A^T_{K-1}E^*  & E^*  \end{array}\right]$
\end{center} has rank $n$. Hence, the matrix
\begin{center}
   $\left[\begin{array}{cccc} * &A^{-T}_{K}A^{-1}_{K-1}.Z_{K-1}.E^* & A^{-T}_{K}A^{-1}_{K-1}D_{Z_{K-1}}(e_n) & 0 \\ * & A_{K}A^T_{K-1}E^* & *  & E^*  \end{array}\right]$
\end{center} 
is full rank as $A^{-T}_{K}A^{-1}_{K-1}D_{Z_{K-1}}(e_n)$ is of rank $n$. Thus, if $A_{K-1}.e_n \neq e_n$, then the matrix $J\Phi_K$ has full rank. \\
Our last step is to prove that if $A_{K-1}.e_n = e_n$ and $Z_{K-1}.e_n =  0$, $J\Phi_K$ is not a matrix of full rank.  Now let $K=2k+2$ and $\left[\begin{array}{ccc}I_n & 0 \end{array}\right]J\Phi_{S} = H_S$ and $\left[\begin{array}{ccc}0 & I_n \end{array}\right]J\Phi_{S} = G_S$. For $K-1= 2k+1$, Since $J\Phi_{2k}(\Vec{A}_{2k},\Vec{Z}_{2k}) =\left[\begin{array}{ccc} H_{2k} \\ G_{2k} \end{array}\right] $, we have: 
 \begin{center}
     $ \left[\begin{array}{ccc} M_{2k+1}(A_{2k+1},Z_{2k+1}).J\Phi_{2k}(\Vec{A}_{2k},\Vec{Z}_{2k}) \end{array}\right] =  \left[\begin{array}{ccc} A^{-1}_{2k+1}H_{2k}+A^{-1}_{2k+1}.Z_{2k+1}.G_{2k}  \\ 
 A^{T}_{2k+1}G_{2k} \end{array}\right]$.
 \end{center}
 For $\mathcal{A}_{2k+1}(e_{2n}) $, we have:
\begin{center}
    $\mathcal{A}_{2k+1}(e_{2n}) = \left[\begin{array}{ccc} D_{A^{-1}_{2k+1}}(Z_{2k+1}.e_n) & A^{-1}_{2k+1}D_{Z_{2k+1}}(e_n) \\  0 & 0\end{array}\right]$.
\end{center}
With the fact that $Z_{2k+1}.e_n = 0$, we have:
\begin{center}
     $J\Phi_{2k+1}(\Vec{A}_{2k+1},\Vec{Z}_{2k+1})= 
    \left[\begin{array}{cccc}  A^{-1}_{2k+1}H_{2k}+A^{-1}_{2k+1}.Z_{2k+1}.G_{2k} &A^{-1}_{2k+1}D_{Z_{2k+1}}(e_n) \\ 
 A^{T}_{2k+1}G_{2k} & 0 \end{array}\right]$.
\end{center}
Now we compute $  J\Phi_{2k+2}(\Vec{A}_{2k+2},\Vec{Z}_{2k+2})$.
For $ \mathcal{A}_{2k+2}(\Phi_{2k+1}(\Vec{A}_{2k+1},\Vec{Z}_{2k+1}))$, it has the form:
    \begin{center}
        $\mathcal{A}_{2k+2}(e_{2n})=
       \left[\begin{array}{ccc} 0 & 0  \\  D_{A_{2k+2}}(e_n) = E^* & 0\end{array}\right]$.
    \end{center}
Then by Lemma \ref{comp-Jac}, we can write $J\Phi_{2k+2}(\Vec{A}_{2k+2},\Vec{Z}_{2k+2})$ as:
\begin{center}
     $
 \left[\begin{array}{cccc}   A_{2k+2}^{-T}\mathcal{D}_{2k+2} &A_{2k+2}^{-T}\mathcal{E}_{2k+2}  & 0\\ 
 Z_{2k+2}.A^{-T}_{2k+2}.\mathcal{D}_{2k+2}+A_{2k+2}A^{T}_{2k+1}G_{2k} & Z_{2k+2}.A^{-T}_{2k+2}.\mathcal{E}_{2k+2}  & E^*\end{array}\right]$
\end{center}
with $\mathcal{D}_{2k+2} =  A^{-1}_{2k+1}H_{2k}+A^{-1}_{2k+1}.Z_{2k+1}.G_{2k}$, $\mathcal{E}_{2k+2} =  A^{-1}_{2k+1}D_{Z_{2k+1}}(e_n)$. By row operations, we have:
\begin{center}
    $\mathcal{F} =
 \left[\begin{array}{cccc}   A_{2k+2}^{-T}\mathcal{D}_{2k+2} &A_{2k+2}^{-T}\mathcal{E}_{2k+2} & 0\\ 
 A_{2k+2}A^{T}_{2k+1}G_{2k} &  0 &  E^*\end{array}\right]$. 
\end{center}
Now we have: 
\begin{center}
    $G_{2k} = \left[\begin{array}{cccc} 
 Z_{2k}.A^{-T}_{2k}.\mathcal{D}_{2k}+A_{2k}A^{T}_{2k-1}G_{2k-2} & Z_{2k}.A^{-T}_{2k}.\mathcal{E}_{2k}  & E^*\end{array}\right]$
\end{center}
The rank of $G_{2k}$ is always higher than or equal to $n-1$ since $\rank E^* = n-1$. We can see that $\rank G_{2k} =n$ if $  (Z_{2k}.A^{-T}_{2k}.\mathcal{E}_{2k})$ or $Z_{2k}.A^{-T}_{2k}.\mathcal{D}_{2k}+A_{2k}A^{T}_{2k-1}G_{2k-2}$ have a nontrivial last row. Also, since $A^{-T}_{2k}.\mathcal{E}_{2k}$ is full rank, the non-triviality of the last row of $  (Z_{2k}.A^{-T}_{2k}.\mathcal{E}_{2k})$ is the same as $Z_{2k}$, we then have:
\begin{center}
    $\rank G_{2k} =  \rank \left[\begin{array}{cccc} 
 Z_{2k}.A^{-T}_{2k}.\mathcal{D}_{2k}+A_{2k}A^{T}_{2k-1}G_{2k-2} & Z_{2k}  & E^*\end{array}\right]$.
\end{center}
Since $Z_{2k}$ has the trivial last row and the matrix $A_{2k-1}$ has the last column as $e_n$, we get: 
\begin{center}
    $\rank G_{2k} =  \rank \left[\begin{array}{cccc} 
G_{2k-2} & Z_{2k}  & E^*\end{array}\right]$.
\end{center}
Then, we get the following: 
\begin{center}
 $\mathrm{rank}(G_{2k}) =\rank ( Z_{2k}|Z_{2k-2}  
|\cdots|Z_2|E^*)$.   
\end{center}
Since $(\vec{A}_{K-1},\vec{Z}_{K-1}$) is in the singular locus, we must have that all $Z_{2i}$ must have trivial last row, and hence $G_{2k}$ have trivial last row, and since $A^T_{2k+1}$ has $e_n$ as its last row, $A^T_{2k+1}G_{2k}$ must have trivial last row.\\
Now, assume $K-1$ is even, then we can repeat the proof of case $K = 3$. Indeed, since we are in the singular locus, by Lemma \ref{lemm-sing1} we have $\Phi_S(\vec{A}_S,\vec{Z}_S) = e_n$ when $S \leq K-2$. Thus, the proof can be repeated as the proof for $K=3$.
\end{proof}
\end{comment}

Recall that $S_K$ denotes the singularity set of the map $\Phi_K$ described in Proposition \ref{SingSet}. From \cite[Theorem 3.5]{Scho} we have the following. 
\begin{lemm}
    For $n \ge 2, K \ge 3$, the submersion $\Phi_K : (\mathbb{C}^{\frac{n(n-1)}{2}} \times \mathbb{C}^{\frac{n(n+1)}{2}})^K \setminus S_K \to \mathbb{C}^{2n}\setminus \{0\}$ is surjective. 
\end{lemm}


We start the proof of Proposition \ref{reg-locus}, namely the stratified ellipticity of the submersion $\Phi_K$ outside the singularity set $S_K$. The proof uses the method in \cite{KI}. Over each stratum of $\mathbb{C}^{2n}\setminus \{0\}$ in an appropriate stratification, we are able to reduce the $2n$ defining polynomial equations of the fibers of $\Phi_K$ to one single polynomial equation $p(x_1, \cdots , x_m) = a$. It turns out that the polynomial $p$ will be no more than linear in each of the variables $x_i, i = 1, \dots, m$. Therefore, the vector fields 
\begin{align*}
    V_{ij,p} = \frac{\partial p}{\partial x_i}\frac{\partial}{\partial x_j} - \frac{\partial p}{\partial x_j}\frac{\partial}{\partial x_i}, \quad  1 \leq i < j \leq m
\end{align*}
will be globally integrable \cite[Lemma 5.2]{KI} and span the tangent space of the fiber $\{ p = a \}$ at all smooth points \cite[Lemma 5.3]{KI}. The dominating fiber spray is then constructed as in Remark \ref{VF-spray}. 



Recall that
\begin{center}
       For $k$ even, $M^+(A_k,Z_k) =\left[\begin{array}{ccc} I & Z_k \\ 0 & I\end{array}\right] \left[\begin{array}{ccc} A^{-1}_k & 0 \\ 0 & A^T_k\end{array}\right]  $, 
    \end{center}
    and
    \begin{center}
    for $k$ odd, $M^-(A_k,Z_k)=\left[\begin{array}{ccc} I & 0 \\ Z_k & I \end{array}\right] \left[\begin{array}{ccc} A_k^{-T} & 0 \\ 0 & A_k \end{array}\right] $.
\end{center}
 For $K$ even, we have:
 \begin{center}
     $\Phi_K((A_1,Z_1),\cdots (A_K,Z_K)) = P^K =  e^T_{2n} M^-(A_1,Z_1) M^+(A_2,Z_2) \cdots M^-(A_{K-1},Z_{K-1}) M^+(A_K,Z_K)$.
 \end{center}
 Hence we have a recursive formula:
 \begin{center}
     $P^K = P^{K-1} M^+(A_K,Z_K)$
 \end{center}
Denote $P^K_f=(P^K_1,\cdots,P^K_n)$, $P^K_s=(P^K_{n+1},\cdots,P^K_{2n})$. Then the above equation can be writen as
\begin{center}
    $P^K \left[\begin{array}{ccc} A_K & 0 \\ 0 & A^{-T}_K\end{array}\right] = P^{K-1} \left[\begin{array}{ccc} I_n & Z_K \\ 0 & I_n\end{array}\right]$
\end{center}
and this is equivalent to
\begin{equation}\label{equa-2}
    \begin{cases}
    P^K_f A_K = P^{K-1}_{f} \\
      P^K_s A^{-T}_K = P^{K-1}_s + P^{K-1}_f Z_K 
\end{cases}
\end{equation}
similarly we have
\begin{center}
     $P^{K-1} = P^{K-2} M^-(A_{K-1},Z_{K-1})$
\end{center}
which is equivalent to
\begin{center}
    $P^{K-1} \left[\begin{array}{ccc} A^T_{K-1} & 0 \\ 0 & A^{-1}_{K-1}\end{array}\right] = P^{K-2} \left[\begin{array}{ccc} I_n & 0 \\ Z_{K-1} & I_n\end{array}\right]$
\end{center}
and
\begin{equation}\label{equa-3}
    \begin{cases}
    P^{K-1}_f A^T_{K-1} = P^{K-2}_f+P^{K-2}_s Z_{K-1}  \\
      P^{K-1}_s A^{-1}_{K-1} = P^{K-2}_{s}
\end{cases}
\end{equation}

We define the following stratification of $\mathbb{C}^{2n} \setminus \{0 \}$
\begin{align*}
    G_i &=  \lbrace (a,b) \in \mathbb{C}^n \times \mathbb{C}^n : a_1 =  \cdots = a_{i-1} =0, a_i \neq 0 \rbrace , \quad i = 1, \dots, n \\
    N_i &= \lbrace (0,b) \in \mathbb{C}^n\times \mathbb{C}^n: b_n = \cdots = b_{n-i+1} =  0, b_{n-i} \neq 0  \rbrace, \quad i = 0, \dots, n-1
\end{align*}

\begin{lemm}\label{red-1}
Let $(a,b)\in G_i \subset \mathbb{C}^{2n}\setminus \{0\}, i = 1, \dots, n$. Then the fiber $(P^K)^{-1}(a,b)$ given by the set of $2n$ equations $P^K(A,Z) = (a, b)$ can be described by one equation, namely $P^K_i = a_i$.
\end{lemm}
\begin{proof}
    Step 1: Observe that for $j>i$
    \begin{align*}
        P^{K-1}_j = P^K_f A_K e_j = P^K_j+\sum_{k=1}^{j-1} a_{kj,K}P_k^K 
        = P^K_j+a_{ij,K}P_i^K+\sum_{k=i+1}^{j-1} a_{kj,K}P_k^K. 
    \end{align*}
Over $G_i$ we have $P^{K}_i = a_i \neq 0$, thus we can use these equations to express $a_{ij,K}$
\begin{align*}
    P^{K-1}_j &= P^K_j + a_{ij,K} a_i + \sum_{k=i+1}^{j-1} a_{kj,K}P_k^K, \\
    \implies a_{ij,K} &=(P^{K-1}_j-P^K_j-\sum_{k=i+1}^{j-1} a_{kj,K}P_k^K)/a_i \quad \text{for } j = i+1, \dots, n.
\end{align*}
Step 2: We now use the second set of equations in (\ref{equa-2}) to express $n$ variables in $Z_K$. Over $G_i$, we have $P^{K-1}_i = P^K_i = a_i \neq 0$. Let $\alpha = P^K_s A^{-T}_K-P^{K-1}_s$.
% \begin{center}
%     $\alpha = P^K_s A^{-T}_K-P^{K-1}_s$.
% \end{center}
By (\ref{equa-2}) we have $\alpha = P^{K-1}_f Z_K$, that is
\begin{center}
    $\alpha_j = \alpha e_j = P^{K-1}_f Z_K e_j = \sum^n_{l=1} P^{K-1}_l z_{lj,K} $ for $j = 1, \cdots,n$.
\end{center}
When $j \neq i$
\begin{center}
    $z_{ij,K}= (\alpha_j-\sum^n_{l=1,l\neq i}z_{lj,K}P^{K-1}_l)/a_i $
\end{center}
and
\begin{center}
   $z_{ii,K} = (\alpha_i-\sum^n_{l=1,l\neq i}z_{il,K}P^{K-1}_l)/a_i$ (by symmetry of $Z_K$)
\end{center}
and then we have
\begin{center}
    $z_{ii,K}= \frac{1}{a_i}(\alpha_i-\sum^n_{l=1,l\neq i}P^{K-1}_l(\frac{1}{a_i}(\alpha_i-\sum^n_{k=1,k\neq i}z_{ki,K}P^{K-1}_k)))$.
\end{center}
It remains to consider the following system of equations
\begin{equation}
    \begin{cases}
    P^K_1 = P^{K-1}_1=0 \\
    \quad \quad \, \vdots \\
      P^K_{i-1}= P^{K-1}_{i-1} =  0\\
      P^K_i = P^{K-1}_{i} = a_i \neq 0
\end{cases}
\end{equation}
Step 3: Now we use the first set of equation in (\ref{equa-3}) to express $(i-1)$ variables of $A_{K-1}$ as follows:
\begin{center}
    $P^{K-1}_f A^T_{K-1} = P^{K-2}_f+P^{K-2}_s Z_{K-1} =:  \beta$
\end{center}
and
\begin{center}
    $\beta_j =  P^{K-1}_f A^T_{K-1} e_j = P^{K-1}_j + \sum^n_{k=j+1} P^{K-1}_k a_{jk,K-1}$.
\end{center}
For $j =  1,2,\cdots,i-1$
\begin{align*}
    \beta_j &= P^{K-1}_j + a_i a_{ji, K-1} + \sum^n_{k=j+1, k \neq i} P^{K-1}_k a_{jk,K-1} \\
    \implies a_{ji,K-1}  &= (\beta_j-P^{K-1}_j-\sum^n_{k=j+1,k \neq i}P^{K-1}_ka_{jk,K-1})/a_i.
\end{align*}
\end{proof}


\begin{lemm}\label{red-2}
     Let $(0,b)\in N_i \subset \mathbb{C}^{2n}\setminus \{0\}, i = 0, \dots, n-1$. Then the fiber $(\Phi_K|_{U_K})^{-1}(0,b)$ given by the set of $2n$ equations $P^K(A,Z)=(0,b)$ can be described by one equation, namely $P^K_{2n-i}=b_{n-i}$.
\end{lemm}
\begin{proof}%[Proof of Lemma \ref{red-2}:]
Recall the inductive formula
\begin{align*}
    P^K = P^{K-1} M^+(A_K,Z_K) = P^{K-1} \left[\begin{array}{ccc} I & Z_K \\ 0 & I\end{array}\right] \left[\begin{array}{ccc} A^{-1}_K & 0 \\ 0 & A^T_K\end{array}\right].
\end{align*}
We consider the fiber of $(0,b)$ in $N_i$. Then (\ref{equa-2}) becomes
\begin{align*}
    \begin{cases}
    0 = P^{K-1}_{f} \\
    b A^{-T}_K = P^{K-1}_s  
    \end{cases}
\end{align*}
which is equivalent to
\begin{align} \label{P_f=0}
    \begin{cases}
    0 = P^{K-1}_{f}  \\
    b  = P^{K-1}_s  A^{T}_K
    \end{cases}
\end{align}
Writing out the second equation, we have
\[
    ( b_1, \dots, b_{n-i}, 0, \dots, 0) = P^{K-1}_s 
    \begin{bmatrix}
     1 & 0 & \cdots &0 \\  a_{12,K} & 1 & \ddots & 0 \\
    \vdots & \ddots & \ddots & \vdots \\
    a_{1n,K} & \cdots & a_{n-1,n,K} & 1
    \end{bmatrix}
\]
The last $n+i+1$ components yield
\begin{align} \label{n+i+1}
    P^{K-1}_{2n-i}=b_{n-i}, \, P^{K-1}_{2n-i+1}=0, \, \dots, P^{K-1}_{2n-1} = 0, \, P^{K-1}_{2n} = 0.
\end{align}
Since $b_{n-i}$ is nonzero, we can use the first $n-i-1$ equations to express the variables $a_{k,n-i,K}, k = 1, \dots, n-i-1$ in $A_K$ as in Step 1 of the proof of Lemma \ref{red-1}.

We write $P^{K-1} = P^{K-2} M^-(A_{K-1}, Z_{K-1})$
\begin{center} \label{P^K-1}
    $P^{K-1} = P^{K-2} \left[\begin{array}{ccc} I_n & 0 \\ Z_{K-1} & I_n\end{array}\right] \left[\begin{array}{ccc} A^{-T}_{K-1} & 0 \\ 0 & A_{K-1}\end{array}\right] = P^{K-2} \left[\begin{array}{ccc} A^{-T}_{K-1} & 0 \\ 0 & A_{K-1}\end{array}\right] \left[\begin{array}{ccc} I_n & 0 \\ \tilde{Z}_{K-1} & I_n\end{array}\right]$,
\end{center}
where $\tilde{Z}_{K-1}:= A_{K-1}^{-1} Z_{K-1} A_{K-1}^{-T}$ is also symmetric. 


\begin{center}
    $P^{K-1} \left[\begin{array}{ccc} I_n & 0 \\ -\tilde{Z}_{K-1} & I_n\end{array}\right] = P^{K-2} \left[\begin{array}{ccc} A^{-T}_{K-1} & 0 \\ 0 & A_{K-1}\end{array}\right]$
\end{center}
equivalently
\begin{equation}
\begin{cases}
    P^{K-1}_f - P^{K-1}_s \tilde{Z}_{K-1} &= P^{K-2}_f A^{-T}_{K-1} \\
    P^{K-1}_{s} A^{-1}_{K-1} &= P^{K-2}_{s} 
\end{cases}  
\end{equation}
Therefore, the first set of equations from \eqref{P_f=0} and the equations \eqref{n+i+1} are equivalent to 
\begin{align*}
    \begin{cases}
        - P^{K-1}_s \tilde{Z}_{K-1} &= P^{K-2}_f A^{-T}_{K-1} \\
        P^{K-1}_{2n-i} &= b_{n-i} \\
        \sum^{j-1}_{k=1}d_{kj}P^{K-1}_{n+k} & = P^{K-2}_{n+j}, \quad j = n-i+1, \cdots, n  
    \end{cases}
\end{align*}
where we choose 
\begin{center}
     $A^{-1}_{K-1}  = \left[\begin{array}{cccc} 1 & d_{12} & \cdots & d_{1n} \\  0 & 1 & \ddots & \vdots \\
    \vdots & \ddots & \ddots & d_{n-1,n} \\
    0 & \cdots & 0 & 1\end{array}\right]$.
\end{center}
A closer inspection on the last set of equations, together with \eqref{n+i+1}, gives
\begin{align}
    d_{n-i,j} b_{n-i} + \sum^{n-i-1}_{k=1}d_{kj}P^{K-1}_{n+k} = P^{K-2}_{n+j}  \quad \text{ for } j = n-i+1, \dots, n. 
\end{align}
Since $b_{n-i}$ is nonzero, we can use these equations to express $d_{n-i,j}$
\[
     d_{n-i,j} =  (P^{K-2}_{n+j} - \sum^{n-i-1}_{k=1}d_{kj}P^{K-1}_{n+k})/b_{n-i}  \quad  \text{ for } j = n-i+1, \dots, n. 
\]
Since $P^{K-1}_s \neq 0$, the first set of equations $- P^{K-1}_s \tilde{Z}_{K-1} = P^{K-2}_f A^{-T}_{K-1}$ can be used to express $n$ variables in $\tilde{Z}_{K-1}$ as in Step 2 of proof of Lemma \ref{red-1}. 

The introduction of $\tilde{Z}_{K-1}$ and $d_{ij}, i>j$ is another choice of coordinates on the set of unitriangular $\tilde{\Omega}$-symplectic matrices. It allows to reduce the set of $2n$ equations in a simple way to one single equation. 
\end{proof}

For $K$ odd the reduction to one equation is very similar. The stratification of $\mathbb{C}^{2n} \setminus \{0\}$ is as follows:
\begin{align*}
    &\lbrace (a,b) \in \mathbb{C}^n\times \mathbb{C}^n: b_1 = \cdots = b_{i-1} =  0, b_{i} \neq 0  \rbrace \quad \text{ for } i = 1, \dots, n, \\
    &\lbrace (a,0) \in \mathbb{C}^n\times \mathbb{C}^n: a_n = \cdots = a_{n-i+1} =  0, a_{n-i} \neq 0  \rbrace \quad \text{ for } i = 0, \dots, n-1.
\end{align*}

%First consider the case $i=1$, then on the stratum $G_1 \backslash G_2$ we have: 
%\begin{center}
%    $\Phi_K: \Phi^{-1}_K(G_1\backslash G_{2}) \to G_1\backslash G_{2}$
%\end{center}
%Indeed, the fibration above is bi-holomorphic to the fibration: 
%\begin{center}
%       \begin{tikzcd}
% \lbrace (X,a,b) \in \mathbb{C}^M \times G_1\backslash G_{2}; P_{K,1}(X) = a_1 \rbrace  \arrow[d,"pr"] \\
% G_1\backslash G_{2}.
%\end{tikzcd}
%$pr: (X,\omega,a,b) \mapsto (a,b)$
%\end{center}
%Through the bi-holomorphic map  (also as a graph):
%\begin{center}
%    $\psi_1: \Phi^{-1}_K(G_1\backslash G_{2}) \to \Phi^{-1}_K(G_1\backslash G_{2}) \times \mathbb{C}^n $
%\end{center}
%such that $\psi_1(x) = (x,f_1(x))$ with:
%\begin{center}
%    $f_i=(a_{12,K},\cdots,a_{1n,K},z_{11,K},\cdots,z_{1n,K}): \mathbb{C}^M \to \mathbb{C}^{n-1}\times \mathbb{C}^n$.
%\end{center}
%Here, 
%\begin{center}
%$\mathbb{C}^M = \left( \mathbb{C}^{\frac{n(n+1)}{2}} \times \mathbb{C}^{\frac{n(n-1)}{2}}  \right )^{K-1}\times \mathbb{C}^{\frac{n(n-1)}{2}} \times \mathbb{C}^{\frac{(n-1)(n-2)}{2}} $.    
%\end{center}
%Since $P_{1}^K$ is no more than linear in $z_{1j}$ and $a_{1k}$, by Lemma \ref{lemm-vectorfield1} and Lemma \ref{lemm-vectorfield2}, we have the global fields $V_{ij,P_{i,K}}$, $\frac{\partial}{\partial \omega_l}$ and $\frac{\partial}{\partial \nu_h}$ span the tangent vectors at each fibers. Thus these vector fields give us a spray over $G_1\backslash G_{2}$.\\ \\
%Now consider the case $i=2$, on the stratum $G_1 \backslash G_2$ we have: 
%\begin{center}
%    $\Phi_K: \Phi^{-1}_K(G_2\backslash G_{3}) \to G_2\backslash G_{3}$
%\end{center}
%and the fibration above is bi-holomorphic to the fibration: 
%\begin{center}
%       \begin{tikzcd}
% \lbrace (X,a,b) \in \mathbb{C}^M \times G_1\backslash G_{2}; P_{K,2}(X) = a_2 \rbrace \times \mathbb{C}^{n-1}_{\omega} \arrow[d,"pr"] \\
% G_1\backslash G_{2}
%\end{tikzcd}
%$pr: (X,\omega,a,b) \mapsto (a,b)$
%\end{center}
%through the bi-holomorphic map  (also as a graph):
%\begin{center}
%    $\psi_2: \Phi^{-1}_K(G_2\backslash G_{3}) \to \Phi^{-1}_K(G_2\backslash G_{3}) \times \mathbb{C}^n $
%\end{center}
%such that $\psi_2(x) = (x,f_2(x))$ with:
%\begin{center}
%    $f_i=(a_{23,K},\cdots,a_{2n,K},z_{21,K},\cdots,z_{2n,K}): \mathbb{C}^M \to \mathbb{C}^{n-2}\times \mathbb{C}^n$.
%\end{center}
%Here
%\begin{center}
 %$\mathbb{C}^M = \left( \mathbb{C}^{\frac{n(n+1)}{2}} \times \mathbb{C}^{\frac{n(n-1)}{2}}  \right )^{K-1}\times \mathbb{C}^{\frac{n(n-1)}{2}} \times \mathbb{C}^{\frac{(n-2)(n-3)}{2}} $   
%\end{center}
%and $ \mathbb{C}^{n-1}_{\omega} =  \mathbb{C}^{n-1}_{a_{12,K},a_{13,K}, \cdots a_{1n,K}}$ since we have free variables $a_{12,K},a_{13,K}, \cdots, a_{1n,K}$. By Lemma \ref{lemm-vectorfield1} and Lemma \ref{lemm-vectorfield2}, we have that the global fields $V_{ij,P_{i,K}}$, $\frac{\partial}{\partial a_{1j}}$ span the tangent vectors at each fiber. Thus, these vector fields give us a spray over $G_2\backslash G_{3}$.
%\\
% In the stratum $G_i$ with $i\geq 1$ we have constructed a stratum that admits a spray:
% \begin{center}
%     $\Phi_K: \Phi^{-1}_K(G_i\backslash G_{i+1}) \to G_i\backslash G_{i+1}$.
% \end{center}
% This fibration is biholomorphic to the fibration: 
% \begin{center}
%        \begin{tikzcd}
%  \lbrace (X,a,b) \in \mathbb{C}^M \times G_i\backslash G_{i+1}; P_{K,i}(X) = a_i \rbrace \times \mathbb{C}^{N}_{\omega}  \arrow[d,"pr"] \\
%  G_i\backslash G_{i+1}
% \end{tikzcd}
% $pr: (X,\omega,a,b) \mapsto (a,b)$
% \end{center}
% through the biholomorphic map (also as a graph):
% \begin{center}
%     $\psi_i: \Phi^{-1}_K(G_i\backslash G_{i+1}) \to \Phi^{-1}_K(G_i\backslash G_{i+1}) \times \mathbb{C}^n $
% \end{center}
% such that $\psi_i(x) = (x,f_i(x))$ with:
% \begin{center}
%     $f_i=(a_{i(i+1),K},\cdots,a_{in,K},z_{i1,K},\cdots,z_{in,K}): \mathbb{C}^M \to \mathbb{C}^{n-i-1}\times \mathbb{C}^n$.
% \end{center}
% By Lemma \ref{lemm-vectorfield1} and Lemma \ref{lemm-vectorfield2}, we have that the global fields $V_{ij,P_{i,K}}$ and $\frac{\partial}{\partial \omega_l}$ span the tangent vectors at each fiber. Thus, the flows of these vector fields give us a spray over $G_i\backslash G_{i+1}$.\\
% Now, from the lemma \ref{red-2}, we have shown that in $N_i$ we can reduce everything to one equation, which is the same as the stratum $N_i$. The rest is the analog as of the proof for stratum $N_i$.  \\
% Now we come to the case $K$ odd. We first choose the appropriated coordinates: 
% \begin{center}
%        For $k$ even, $M_k(A_k,Z_k) =\left[\begin{array}{ccc} I & Z_k \\ 0 & I\end{array}\right] \left[\begin{array}{ccc} A^{-1}_k & 0 \\ 0 & A^T_k\end{array}\right]  $, 
%     \end{center}
%     and
%     \begin{center}
%     for $k$ odd, $M_k(A_k,Z_k)=\left[\begin{array}{ccc} I & 0 \\ Z_k & I \end{array}\right] \left[\begin{array}{ccc} A_K^{-T} & 0 \\ 0 & A_k \end{array}\right] $.
% \end{center}
%  For $K$ odd, we have:
%  \begin{center}
%      $\Phi_K((A_1,Z_1),\cdots (A_K,Z_K)) = P^K =  e^T_{2n}M_1(A_1,Z_1) M_2(A_2,Z_2) \cdots M_{K-1}(A_{K-1},Z_{K-1}) M_K(A_K,Z_K)$
%  \end{center}
% Let $A^{-1}_k = D_k$. By the same argument above, we have:
% \begin{equation}\label{equa-4}
%     \begin{cases}
%     P^{K}_f D^{-T}_{K} = P^{K-1}_f+P^{K-1}_s Z_{K} \\
%       P^{K}_s D_{K} = P^{K-1}_{s}
% \end{cases}
% \end{equation}
% and
% \begin{equation}\label{equa-5}
%     \begin{cases}
%     P^{K-1}_f D^{-1}_{K-1} = P^{K-2}_{f} \\
%       P^{K-1}_s D^{T}_{K-1} = P^{K-2}_s+P^{K-2}_f Z_{K-1} .
% \end{cases}
% \end{equation}
% Now consider the stratification 
% \begin{center}
%     $G_i = \lbrace (a,b) \in \mathbb{C}^n\times \mathbb{C}^n: b_1 = \cdots = b_{i-1} =  0, b_{i} \neq 0  \rbrace$.
% \end{center}

% By the same proof as Lemma \ref{red-1}, we get for $j = n-i+1, \cdots, n$:
% \begin{center}
%     $d_{(n-i)j,K} =  \frac{1}{P^{K}_{2n-i}}(P^{K-1}_{n+j}-P^{K}_{n+j}-\sum^{j-1}_{k=1,k \neq n-i}d_{kj,K}P^{K}_{n+k})$.
% \end{center}
% Now using the first Equation of (\ref{equa-4}), we can express $n$ variables of $Z_K$, as in Lemma \ref{red-1}. For $i \neq j$:
% \begin{center}
%     $z_{ij,K}= \frac{1}{P^{K-1}_{2n-i}}(-\sum^{n-1}_{l=0,l\neq i-1}z_{(n-l)j,K}P^{K-1}_{2n-l}) $
% \end{center}
% and
% \begin{center}
%    $z_{ii,K} = \frac{1}{P^{K-1}_{2n-i}}(-\sum^{n}_{l=1,l\neq i}z_{li,K}P^{K-1}_{n+l})$ (by symmetry of $Z_K$)
% \end{center}
% and then we have:
% \begin{center}
%     $z_{ii,K}= \frac{1}{P^{K-1}_{2n-i}}(-\sum^{n-1}_{l=0,l\neq i}P^{K-1}_{n+l}(\frac{1}{P^{K-1}_{2n-i}}(-\sum^{n}_{k=1,k\neq n-i}z_{ki,K}P^{K-1}_{n+k})))$.
% \end{center}
% Now we express $(i-1)$ variables of $D_{K-1}$ from the second equation of (\ref{equa-5}):
% \begin{center}
%     $P^{K-1}_s D^T_{K-1} = P^{K-2}_s+P^{K-2}_f Z_{K-1} :=  \gamma$
% \end{center}
% and
% \begin{center}
%     $\gamma_j =  P^{K-1}_s D^T_K e_j = P^{K-1}_{2n-j} + \sum^{n}_{k=j} P^{K-1}_{n+k} d_{kj,K-1}$.
% \end{center}
% For $j =0,1,\cdots,n-i-1$:
% \begin{center}
%     $d_{ij,K-1}  = \frac{1}{P^{K-1}_{2n-i}}(\gamma_j-P^{K-1}_{2n-j}-\sum^n_{k=j+1,k \neq n-i}P^{K-1}_{n+k}d_{kj,K-1}).  $
% \end{center}
% Now we consider the
% \begin{center}
%     $N:= \lbrace (a,0) \in \mathbb{C}^n \times \mathbb{C}^n\rbrace$ 
% \end{center}
% As in the proof of \ref{red-2}, we will change to new coordinates with: 
% \begin{center}
% for $k$ even, $M_k(A_k,Z_k)=\left[\begin{array}{ccc} I & Z_k \\ 0 & I\end{array}\right] \left[\begin{array}{ccc} A^{-1}_k & 0 \\ 0 & A^T_k\end{array}\right]  $.
% \end{center}
% Now we have: 
% \begin{center}
    

% $P^{K-1} = P^{K-2} \left[\begin{array}{ccc} A^{-1}_{K-1} & 0 \\ 0 & A^T_{K-1}\end{array}\right]\left[\begin{array}{ccc} I & Z_{K-1} \\ 0 & I\end{array}\right]  $
% \end{center}
% and we can rewrite as: 
% \begin{center}
%     $P^{K-1}.\left[\begin{array}{ccc} I_n & B \\ 0 & I_n\end{array}\right] = P^{K-2}.\left[\begin{array}{ccc} D & 0 \\ 0 & D^{-T}\end{array}\right]$
% \end{center}
% Where $B=(b_{ij})$ a symmetric matrix and $ A^{-1}_{K} = D$

% This is equivalent to
% \begin{equation}\label{equa-6}
%    \begin{cases}
%     P^{K-1}_s+P^{K-1}_fB = P^{K-2}_s.D^{-T}  \\
%       P^{K-1}_{f}= P^{K-2}_{f}D.
% \end{cases}  
% \end{equation}
% We express by using the set of the second equation of (\ref{equa-6}), just as we did in Lemma \ref{red-1} for the first equation of (\ref{equa-2}) we get the expression of elements in $D$. After that we can express the elements in $B$ by using the same way as in \ref{red-1} (which uses the second equation of (\ref{equa-2})) by using the first equation of (\ref{equa-6}). Then we finish the reduction to one equation in the case $K$ odd. 

% The construction of stratified spray is exactly the same as $K$ even case.

\section{Optimal Factorization of elementary matrices}
In this section we prove that the elementary matrices can be written as a product of unitriangular symplectic matrices with respect to the standard symplectic form. Most importantly, the number of factors is independent of the size of the matrices, i.e., independent of $n$. 
%\subsection{General factorization of elementary matrices}
% In this section, we prove that the symplectic matrices of the form $\left[\begin{array}{ccc} A & B \\ 0 & D\end{array}\right]$ would have at most 8 factors. Now we apply the  proposition \ref{fact-1}  to show that the number of factorization is 8. Notice that our $A_1$ is the triangular nilpotent matrix.


Let $T^+_n(\mathbb{C})$ ($T^-_n(\mathbb{C})$) denote the set of upper (respectively lower) triangular matrices in $\mathrm{GL}_n(\mathbb{C})$, and $UT^+_n(\mathbb{C})$ ($UT^-_n(\mathbb{C})$) the set of upper (resp.\ lower) unitriangular matrices in $\mathrm{GL}_n(\mathbb{C})$. Denote by $Sym_n(\mathbb{C})$ the space of symmetric square matrices of size $n$. 

\begin{lemm}\label{fact-ele-lemm}
Let $A: X \to T^\pm_n(\mathbb{C})$ be a holomorphic map such that the diagonal entries $\lambda_1, \dots, \lambda_n$ are pairwise distinct constants. Then $A$ is holomorphically diagonalizable by a unitriangular matrix, i.e., there exists a holomorphic map $K: X \to UT^\pm_n(\mathbb{C})$ such that $A = K\Lambda K^{-1}$, $\Lambda = \mathrm{diag}(\lambda_1, \dots, \lambda_n)$.   
\end{lemm}
\begin{proof}
    Let us consider the case $A = (a_{ij})$ is upper triangular. 
    Let $K$ be an upper unitriangular matrix and $\alpha_{ij}, 1 \le i < j \le n$ the entries of $K$ which are above the diagonal. The set of linear equations given by $A K = K \Lambda$ can be uniquely solved for $\alpha_{i, i+1}$ using the equation at position $(i, i+1)$
    \[
        a_{i, i+1} + \lambda_{i+1} \alpha_{i, i+1} = \lambda_i \alpha_{i, i+1}.
    \]
    After having solved the first superdiagonal, continue with the second superdiagonal, and so on. 
\end{proof}

\begin{theo}\label{fact-ele}
    Given holomorphic maps $A: X \to UT^+_n(\mathbb{C})$ and $Z: X \to Sym_n(\mathbb{C})$, there exist holomorphic maps $Z_i, \tilde{Z}_i : X \to Sym_n(\mathbb{C}), i = 1, \dots, 7$ such that
    \begin{align*}
        M^-(A, Z) &= 
        \begin{bmatrix} I & 0 \\ Z & I   \end{bmatrix}
        \begin{bmatrix} A^{-T} & 0 \\ 0 & A  \end{bmatrix}
        =
        \begin{bmatrix} I & 0 \\ Z_1 & I   \end{bmatrix}
        \begin{bmatrix} I & Z_2 \\ 0 & I   \end{bmatrix}
        \cdots 
        \begin{bmatrix} I & 0 \\ Z_7 & I   \end{bmatrix} \\
        &=
        M^-(Z_1) M^+(Z_2) \cdots M^-(Z_7)
        \quad  \text{ and } \\
        M^+(A, Z) &= 
        \begin{bmatrix} I & Z \\ 0 & I   \end{bmatrix}
        \begin{bmatrix} A^{-1} & 0 \\ 0 & A^{T}  \end{bmatrix}
        =
        \begin{bmatrix} I & \tilde{Z}_1 \\ 0 & I   \end{bmatrix}
        \begin{bmatrix} I & 0 \\ \tilde{Z}_2 & I   \end{bmatrix}
        \cdots 
        \begin{bmatrix} I & \tilde{Z}_7 \\ 0 & I   \end{bmatrix} \\
        &=
        M^+(\tilde{Z}_1) M^-(\tilde{Z}_2) \cdots M^+(\tilde{Z}_7).
    \end{align*}
\end{theo}
\begin{proof}%[Proof of Theorem \ref{fact-ele}:]
First we make the diagonal entries of $A$ pairwise distinct by multiplying it with a constant diagonal matrix $\Lambda$ with such property
\begin{center}
    $\left[\begin{array}{ccc} A & 0 \\ 0 & A^{-T}\end{array}\right] = \left[\begin{array}{ccc} \Lambda & 0 \\ 0 & \Lambda^{-1}\end{array}\right]
    \left[\begin{array}{ccc} \Lambda^{-1}A & 0 \\ 0 & \Lambda A^{-T}\end{array}\right] 
    =  
    \left[\begin{array}{ccc} \Lambda & 0 \\ 0 & \Lambda^{-1}\end{array}\right]
    \left[\begin{array}{ccc} \Lambda^{-1}A & 0 \\ 0 & (\Lambda^{-1} A)^{-T}\end{array}\right]$.
\end{center}
The constant-valued factor decomposes as
\[
    \begin{bmatrix} \Lambda & 0 \\ 0 & \Lambda^{-1}   \end{bmatrix}
    = 
    \begin{bmatrix} I & 0 \\ -\Lambda^{-1}  & I  \end{bmatrix}
    \begin{bmatrix} I & \Lambda - I \\ 0 & I  \end{bmatrix}
    \begin{bmatrix} I & 0 \\ I & I  \end{bmatrix}
    \begin{bmatrix} I &  \Lambda^{-1} - I \\ 0 & I  \end{bmatrix}.
\]
By Lemma \ref{fact-ele-lemm} there exists a holomorphic map $K : X \to UT^+_n(\mathbb{C})$ so that $\Lambda^{-1}A = K\Lambda^{-1} K^{-1}$. 
Write $P_1 = KK^T$, $P_2 =  K^{-T} \Lambda K^{-1}$, then $\Lambda^{-1}A = P_1 P_2$ and $(\Lambda^{-1} A)^{-T} = P_1^{-1} P_2^{-1}$. 
As in \cite[Lemma 1]{JLX2}, we can write the second factor, starting with an upper unitriangular matrix, as
\begin{center}
    $\left[\begin{array}{ccc} P_1 P_2 & 0 \\ 0 & P_1^{-1} P_2^{-1}\end{array}\right] 
    =  
    \left[\begin{array}{ccc} I & -P_1 \\ 0 & I\end{array}\right]
    \left[\begin{array}{ccc} I & 0 \\ P_1^{-1} - P_2 & I\end{array}\right]
    \left[\begin{array}{ccc} I & P_2^{-1} \\ 0 & I \end{array}\right]
    \left[\begin{array}{ccc} I & 0 \\ P_2 P_1 P_2 - P_2 & I\end{array}\right]$.
\end{center}
Combining these two factorizations, taking into account the lower triangular form of the first factor in $M^-$, gives a decomposition of $M^-$ into $7$ factors. For $M^+$ observe that
\[
    \begin{bmatrix} \Lambda & 0 \\ 0 & \Lambda^{-1}   \end{bmatrix}
    = 
    \begin{bmatrix} I & \Lambda - I \\ 0 & I  \end{bmatrix}
    \begin{bmatrix} I & 0 \\ I & I  \end{bmatrix}
    \begin{bmatrix} I &  \Lambda^{-1} - I \\ 0 & I  \end{bmatrix}
    \begin{bmatrix} I & 0 \\ -\Lambda  & I  \end{bmatrix} \quad \text{ and }
\]
\begin{center}
    $\left[\begin{array}{ccc} P_1 P_2 & 0 \\ 0 & P_1^{-1} P_2^{-1} \end{array}\right] 
    =  
    \left[\begin{array}{ccc} I & 0 \\ P_1^{-1}P_2^{-1}P_1^{-1}-P_1^{-1} & I\end{array}\right]\left[\begin{array}{ccc} I & P_1 \\ 0 & I\end{array}\right]\left[\begin{array}{ccc} I & 0 \\ P_2-P_1^{-1} & I\end{array}\right]\left[\begin{array}{ccc} I & -P_2^{-1} \\ 0 & I \end{array}\right]$. \qedhere
\end{center}
% \[
%     \begin{bmatrix} I & \Lambda - I \\ 0 & I  \end{bmatrix}
%     \begin{bmatrix} I & 0 \\ I & I  \end{bmatrix}
%     \begin{bmatrix} I &  \Lambda^{-1} - I \\ 0 & I  \end{bmatrix}
%     = 
%     \begin{bmatrix} \Lambda & 0 \\ I & \Lambda^{-1}    \end{bmatrix}
% \]
\end{proof}
% \begin{rema}
%     There is also a direct proof for the diagonalization of the uppertriangular matrix $A: X \to \mathrm{SL}_{n}(\mathbb{C})$ with constant and pairwise different eigenvalues. Consider the equation: 
%     \begin{center}
%         $D.A = A.D$
%     \end{center}
%     Since $A$ is upper triangular with constant and pairwise different elements on the diagonal, it is easy to see that $D$ should also be upper-triangular. First, we choose on the diagonal of $D$ the holomorphic functions $e^{f_1}, e^{f_2}, \cdots, e^{f_n}$ with $f_n \in \mathcal{O}(X)$. Now we prove by induction that we can also construct such invertible $D$. For $n=2$, consider
%     $D = \left[\begin{array}{ccc} e^{f_1} & d \\ 0 & e^{f_2}\end{array}\right]$. Now we have an equation:
%     \begin{center}
%         $\left[\begin{array}{ccc} e^{f_1} & d \\ 0 & e^{f_2}\end{array}\right].\left[\begin{array}{ccc} a_{11} & f_{12} \\ 0 & a_{22}\end{array}\right] = \left[\begin{array}{ccc} a_{11} & 0 \\ 0 & a_{22}\end{array}\right].\left[\begin{array}{ccc} e^{f_1} & d \\ 0 & e^{f_2}\end{array}\right]. $
%     \end{center}
%     From this equation we get: 
%     \begin{equation}
%         e^{f_1}f_{12}+d.a_{22} = a_{11}d
%     \end{equation}
%     we get: $d= \frac{e^{f_1}.f_{12}}{a_{12}-a_{22}}$, which is a well-defined holomorphic function. \\
%     Now we assume that the construction is good in dimension $n-1$ and we will prove that it is also good in dimension $n$. Consider the product:
%     \begin{center}
%         $\left[\begin{array}{ccccc} e^{f_1} & d_{1,2} & \cdots &d_{1,n} \\  0 & e^{f_2} & \cdots & d_{2,n} \\
%     \vdots & \vdots  & \ddots & \vdots \\
%     0 & \cdots & \cdots & e^{f_n}\end{array}\right] \left[\begin{array}{ccccc} a_1 & a_{1,2} & \cdots &a_{1,n} \\  0 & a_2 & \cdots & a_{2,n} \\
%     \vdots & \vdots  & \ddots & \vdots \\
%     0 & \cdots & \cdots & a_n\end{array}\right]= \left[\begin{array}{ccccc} a_1 & 0 & \cdots &0 \\  0 & a_2 & \cdots & 0 \\
%     \vdots & \vdots  & \ddots & \vdots \\
%     0 & \cdots & \cdots & a_n\end{array}\right] \left[\begin{array}{ccccc} e^{f_1} & d_{1,2} & \cdots &d_{1,n} \\  0 & e^{f_2} & \cdots & d_{2,n} \\
%     \vdots & \vdots  & \ddots & \vdots \\
%     0 & \cdots & \cdots & e^{f_n}\end{array}\right]$.
%     \end{center}
%     By induction hypothesis, we can find the holomorphic function $d_{i,j}$, based on $e^{f_k}$ and $a_{i,j}$ with $i < j\leq n-1$. For $d_{1,n}$, we have the equation:
%     \begin{center}
%         $e^{f_1}a_1+d_{1,2}a_{2,n}+\cdots+d_{1,n}a_n = a_1d_{1,n}$
%     \end{center}
%     which we will get:
%     \begin{center}
%         $d_{1,n}= \frac{e^{f_1}a_1+d_{1,2}a_{2,n}+\cdots+d_{1,n-1}.a_{n-1,n}}{a_1-a_n}$
%     \end{center}
%     since all $d_{i,j}$ with $i<j \leq n-1$ are holomorphic functions of $X$ by induction hypothesis and $a_1-a_n \neq 0$ by the fact that every constants on the diagonal are pairwise different (and also different from $0$), we get a well-defined holomorphic function $d_{1,n}$ on $X$. For $d_{i,n}$ we can do the same and all of them are holomorphic functions based on $a_{i,j}$, $d_{i,j}$ and $e^{f_k}$. Thus we complete the proof.
% \end{rema}

% \subsection{Factorization of elementary matrices with four factors}
% In this section, we study the optimal factorization of matrices in $\mathrm{Sp}_{2n}(\mathcal{O}(X))$, inspired by \cite{JLX2}.
% In \cite{JLX2}, the authors proved that non-singular symplectic matrices with entries in a field can be factorized into the product of four elementary symplectic matrices. In this subsection, we check if some special symplectic matrices with holomorphic entries can be factorized into four elementary symplectic matrices. First, consider an elementary symplectic matrix in Definition \ref{def-ele}: $S = \left[\begin{array}{ccc} A & B \\ 0 & D\end{array}\right]  $ with $A$ upper unitriangular, $B$ and $D$ are matrices such that for $J = \left[\begin{array}{ccc} 0 & I \\ -I & 0\end{array}\right]$, they satisfy the matrix equation:
% \begin{center}
%     $\left[\begin{array}{ccc} A & B \\ 0 & D\end{array}\right].\left[\begin{array}{ccc} 0 & I \\ -I & 0\end{array}\right].\left[\begin{array}{ccc} A^T & 0 \\ B^T & D^T\end{array}\right] = \left[\begin{array}{ccc} 0 & I \\ -I & 0\end{array}\right]$
% \end{center}
% and this gives an equation:
% \begin{center}
%     $\left[\begin{array}{ccc} -BA^T+AB^T & AD^T \\ -DA^T & 0\end{array}\right] = \left[\begin{array}{ccc} 0 & I \\ -I & 0\end{array}\right]$
% \end{center}
% and this should give $D = (A^{-1})^T$ and $AB^T = BA^T$

\begin{rema}\label{fact-num-theo1}
    (a) When $n=2$ or $n=3$ in Theorem \ref{fact-ele}, direct computation shows that $4$ instead of $7$ factors suffice: $M^-(A, Z) = 
        M^-(Z_1) M^+(Z_2) M^-(Z_3) M^+(Z_4)$ and $M^+(A, Z) = 
        M^+(\tilde{Z}_1) M^-(\tilde{Z}_2) M^+(\tilde{Z}_3) M^-(\tilde{Z}_4)$. \\
    (b) For $n \ge 4$, direct computation shows that $5$ factors are not sufficient. Thus the number of factors for factoring $M^\pm(A, Z)$ as a product of elementary symplectic matrices with respect to the standard symplectic form is at least $6$. 
\end{rema}
% \begin{proof}
% Let $A$, $B$, $C$ and $D$ are $3 \times 3$ matrices satisfy above condition. We want the factorization:
% \begin{center}
%     $S = \left[\begin{array}{ccc} I & 0 \\ U & I\end{array}\right].\left[\begin{array}{ccc} I & V \\ 0 & I\end{array}\right].\left[\begin{array}{ccc} I & 0 \\ W & I\end{array}\right].\left[\begin{array}{ccc} I & X \\ 0 & I\end{array}\right]$
% \end{center}
% with $U,V,W,X$ are symmetric and this yields a matrices equation:
% \begin{center}
%     $\left[\begin{array}{ccc} A & B \\ 0 & D\end{array}\right] = \left[\begin{array}{ccc} I + VW & X+V(I+WX) \\ U+(I+UV)W & *\end{array}\right] $
    
% \end{center}
% and from this we have a system of equation: 
% \begin{equation}\label{equa-10}
%     \begin{cases}
%       A = I + VW\\
%       U.A+W = 0\\
%       AX+V = B\\
%       WA^{-1} \text{ is symmetry} \\
%       A^{-1}V \text{ is symmetry}
%     \end{cases}.
% \end{equation}
% Suppose $A^{-1} =  \left[\begin{array}{ccc}1 & a_1 & a_2 \\ 0 & 1 & a_3 \\ 0 & 0 & 1\end{array}\right]$.\\
% We have that 
% \begin{center}
%   $W.A^{-1} = \left[\begin{array}{ccc}x_1 & y_1 & y_2 \\ y_1 & x_2 & y_3 \\ y_2 & y_3 & x_3\end{array}\right].\left[\begin{array}{ccc}1 & a_1 & a_2 \\ 0 & 1 & a_3 \\ 0 & 0 & 1\end{array}\right] = \left[\begin{array}{ccc}x_1 & a_1.x_1+y_1 & a_2.x_1+a_3.y_1+y_2 \\ y_1 & y_1.a_1+x_2 & y_1.a_2+a_3x_2+y_3 \\ y_2 & a_1.y_2+y_3 & a_2.y_2 + a_3.y_3+x_3\end{array}\right]$   
% \end{center}
% symmetry implies:
% \begin{equation}
%     \begin{cases}
%      a_1.x_1+y_1 = y_1 \Rightarrow x_1 = 0 \\
%       a_1.x_1+a_3.y_1 = 0 \Rightarrow y_1 = 0\\
%       y_2.a_1 = a_3.x_2
% \end{cases}
% \end{equation}

% Now set $y_2 = a_3.f$ and $x_2=a_1.f$, we have a matrix:
% \begin{center}
%     $W = \left[\begin{array}{ccc}0 & 0 & a_3.f \\ 0 & a_1.f & y_3 \\ a_3.f & y_3 & x_3\end{array}\right]$.
% \end{center}
% We also have that: 
% \begin{center}
%   $A^{-1}.V = \left[\begin{array}{ccc}1 & a_1 & a_2 \\ 0 & 1 & a_3 \\ 0 & 0 & 1\end{array}\right].\left[\begin{array}{ccc}s_1 & t_1 & t_2 \\ t_1 & s_2 & t_3 \\ t_2 & t_3 & s_3\end{array}\right] = \left[\begin{array}{ccc}s_1 + a_1.t_1+a_2.t_2 & t_1+a_1s_2+a_2.t_3 & t_2+a_1.t_3+a_2s_3 \\ t_1+a_3.t_2 & s_2+a_3.t_3 & t_3+a_3.s_3 \\ t_2 & t_3 & s_3\end{array}\right]$   
% \end{center}
% symmetry implies:
% \begin{equation}
%     \begin{cases}
%      s_3 = t_3 = 0\\
%       t_2.a_3 = a_1.s_2
% \end{cases}
% \end{equation}
% let $t_2 = a_1.g$, $s_2 = a_3.g$, we have:
% \begin{center}
%     $V = \left[\begin{array}{ccc}s_1 & t_1 & a_1.g \\ t_1 & a_3.g & 0 \\ a_1.g & 0 & 0\end{array}\right]$.
% \end{center}
% Now from equation (10), we have:
% \begin{center}
%     $V.W = \left[\begin{array}{ccc}s_1 & t_1 & a_1.g \\ t_1 & a_3.g & 0 \\ a_1.g & 0 & 0\end{array}\right].\left[\begin{array}{ccc}0 & 0 & a_3.f \\ 0 & a_1.f & y_3 \\ a_3.f & y_3 & x_3\end{array}\right] = \left[\begin{array}{ccc}0 & -a_1 & -a_2+a_1.a_3 \\ 0 & 0 & -a_3 \\ 0 & 0 & 0\end{array}\right]$.
% \end{center}
% This equivalences to:
% \begin{center}
%     $ \left[\begin{array}{ccc}a_1.a_3g.f & a_1f.t_1+a_1.g.y_3 & a_3.f.s_1+t_1.y_3+a_1.g.x_3 \\ 0 & a_1.a_3.gf & a_3f.t_1+a_3g.y_3 \\ 0 & 0 & a_1.a_3.g.f\end{array}\right] =  \left[\begin{array}{ccc}0 & -a_1 & -a_2+a_1.a_3 \\ 0 & 0 & -a_3 \\ 0 & 0 & 0\end{array}\right]$
% \end{center}
% This implies a system of equation: 
% \begin{equation}
%     \begin{cases}
%      g.f = 0\\
%      a_1ft_1+a_1gy_3 = -a_1\\
%      a_3.f.s_1+t_1.y_3+a_1.g.x_3 = -a_2 + a_1.a_3 \\
%      a_3f.t_1+a_3g.y_3 = -a_3
% \end{cases}
% \end{equation}
% From this, we can take $f = 0$, $g=1$, $y_3 =  -1$, $t_1 = -a_2$, $x_3=a_3$. Then for any $B$ satisfies the condition $AB^T=BA^T$, this gives us a solution for the factorization problem of $M^-(A,B)$ matrix. For the elementary $M^+(A,C)$, the proof go exactly the same as $M^+(A,C)^T = M^{-}(A,C)$.
% \end{proof}
% \subsubsection{$4 \times 4$ case:}
% Now, we will see that the solution type in $4 \times 4$ case is unable to occur.
% \begin{prop}\label{fact-num-prop}
% The $8 \times 8$ elementary matrices $M^+(A,B)$, $M^-(A,C)$ could not always factorize in product of 4 elementary matrices of type $M^+(I,Z)$ and $M^-(I,Z)$.      
% \end{prop}
% \begin{proof}
% Define by the same way as before:
% \begin{center}
%     $W =  \left[\begin{array}{cccc}x_1 & y_1 & y_2 & y_4\\ y_1 & x_2 & y_3 & y_5\\ y_2 & y_3 & x_3 & y_6\\ y_4 & y_5 & y_6 & x_4\end{array}\right]$ and $V = \left[\begin{array}{cccc}s_1 & t_1 & t_2 & t_4 \\ t_1 & s_2 & t_3 & t_5 \\ t_2 & t_3 & s_3 & t_6 \\ t_4 & t_5 & t_6 & s_4\end{array}\right]$
% \end{center}
% it is easy to see that the symmetry condition in (\ref{equa-10}) gives:
% \begin{center}
%     $W = \left[\begin{array}{cccc}0 & 0 & 0 & y_4\\ 0 & 0 & y_3 & y_5\\ 0 & y_3 & x_3 & y_6\\ y_4 & y_5 & y_6 & x_4\end{array}\right]$
% \end{center}
% with the constrains: 
% \begin{equation}
%     \begin{cases}
%      a_1.y_4 = a_6.y_3 \\
%       a_2.y_4+a_4.y_5 = a_5.y_3+a_6.x_3
% \end{cases}
% \end{equation}
% Now let $y_4 = a_6.f$ and $y_3 = a_1.f$, we get: 
% \begin{center}
%     $W = \left[\begin{array}{cccc}0 & 0 & 0 & a_6.f\\ 0 & 0 & a_1.f & y_5\\ 0 & a_1.f & x_3 & y_6\\ a_6.f & y_5 & y_6 & x_4\end{array}\right]$
% \end{center}
% Now, by doing exactly the same way, we have: 
% \begin{center}
%     $V =  \left[\begin{array}{cccc}s_1 & t_1 & t_2 & t_4 \\ t_1 & s_2 & t_3 & 0 \\ t_2 & t_3 & 0 & 0 \\ t_4 & 0 & 0 & 0\end{array}\right]$
% \end{center}
% with constrains:
% \begin{equation}
%     \begin{cases}
%      a_1.t_3 = a_6.t_4 \\
%       a_5.t_4+a_4.t_2 = a_1.s_2+a_6.t_3
% \end{cases}
% \end{equation}
% now let $t_4 = a_1.g$, $t_3 = a_6.g$, we get 
% \begin{center}
%     $V =  \left[\begin{array}{cccc}s_1 & t_1 & t_2 & a_1.g \\ t_1 & s_2 & a_6.g & 0 \\ t_2 & a_6.g & 0 & 0 \\ a_1.g & 0 & 0 & 0\end{array}\right]$.
% \end{center}
% Using (\ref{equa-10}) again, we get: 
% \begin{align*}
%     V.W =  &\left[\begin{array}{cccc}s_1 & t_1 & t_2 & a_1.g \\ t_1 & s_2 & a_6.g & 0 \\ t_2 & a_6.g & 0 & 0 \\ a_1.g & 0 & 0 & 0\end{array}\right].\left[\begin{array}{cccc}0 & 0 & 0 & a_6.f\\ 0 & 0 & a_1.f & y_5\\ 0 & a_1.f & x_3 & y_6\\ a_6.f & y_5 & y_6 & x_4\end{array}\right]\\ 
%      =& \left[\begin{array}{cccc}0 & -a_1 & a_1.a_4-a_2 & -a_1.a_4.a_6+a_2.a_6+a_1.a_5-a_3 \\ 0 & 0 & -a_4 & a_4.a_6-a_5 \\ 0 & 0 & 0 & -a_6 \\ 0 & 0 & 0 & 0\end{array}\right]=A-I.
% \end{align*}
% This equivalence to:
% \begin{center}
%     $V.W =  \left[\begin{array}{cccc}a_1a_6fg & a_1ft_2+a_1gy_5 & t_1a_1f+t_2x_3+a_1gy_6 & a_6f.s_1+t_1y_5+t_2y_6+a_1gx_4 \\ 0 & a_1a_6fg & a_1fs_2+a_6gx_3 & t_1a_6f+s_2y_3+y_6ga_6 \\ 0 & 0 & a_1a_6fg &  a_6ft_2+a_6gy_5 \\ 0 & 0 & 0 & a_1a_6fg\end{array}\right]= \left[\begin{array}{cccc}0 & -a_1 & a_1.a_4-a_2 & -a_1.a_4.a_6+a_2.a_6+a_1.a_5-a_3 \\ 0 & 0 & -a_4 & a_4.a_6-a_5 \\ 0 & 0 & 0 & -a_6 \\ 0 & 0 & 0 & 0\end{array}\right]$.
% \end{center}
% this gives a system of equations:
% \begin{equation}
%     \begin{cases}
%      fg = 0\\
%       a_1ft_2+a_1gy_5 = -a_1\\
%       a_1fs_2+a_6gx_3 = -a_4\\
%       a_6ft_2+a_6gy_5 = -a_6\\
%       t_1a_1f+t_2x_3+a_1gy_6 = a_1.a_4-a_2\\
%       t_1a_6f+s_2y_3+y_6ga_6 = a_4.a_6-a_5\\
%       a_6f.s_1+t_1y_5+t_2y_6+a_1gx_4 = -a_1.a_4.a_6+a_2.a_6+a_1.a_5-a_3
% \end{cases}.
% \end{equation}
% Suppose $f=0$, we have:
% \begin{equation}
%     \begin{cases}
%      gy_5 = -1\\
%      a_6gx_3 = -a_4\\
%      gy_5 = -1\\
%      t_2x_3+a_1gy_6 = a_1.a_4-a_2\\
%      s_2y_3+y_6ga_6 = a_4.a_6-a_5\\
%      t_1y_5+t_2y_6+a_1gx_4 = -a_1.a_4.a_6+a_2.a_6+a_1.a_5-a_3
% \end{cases}.
% \end{equation}
% from here we can see clearly that we could not have any holomorphic solution with variables in $a_i$ since there is an equation $a_6gx_3 = -a_4$. Thus, we finish the proposition for $M^+(A,B)$, and the case $M^-(A,C)$ follows directly.  
% \end{proof}
% \subsubsection{General n:}
% In this case we prove a theorem generalizing from Proposition \ref{fact-num-prop}:
% \begin{prop}\label{fact-num-prop2}
%   The elementary matrices $M^+(A,B)$, $M^-(A,C) \in \mathrm{Sp}_{2n}(O(X))$ with $n \geq 5$ could not always factorize in product of 4 elementary matrices of type $M^+(I,Z)$ and $M^-(I,Z)$.   
% \end{prop}
% \begin{proof}
%  For simplicity, we consider these symmetric matrices:

% \begin{center}
%     $W =  \left[\begin{array}{cccc}w_{1,1} & w_{1,2} & \cdots & w_{1,n}\\ w_{2,1} & w_{2,2} & \cdots & w_{2,n}\\ \vdots & \vdots & \ddots & \vdots \\ w_{n,1} & w_{n,2} & \cdots & w_{n,n}\end{array}\right]$ and $V = \left[\begin{array}{cccc}v_{1,1} & v_{1,2} & \cdots & v_{1,n}\\ v_{2,1} & v_{2,2} & \cdots & v_{2,n}\\ \vdots & \vdots & \ddots & \vdots \\ v_{n,1} & v_{n,2} & \cdots & v_{n,n}\end{array}\right]$
% \end{center}
% and the unitriangular matrix:

% \begin{center}
%     $A^{-1} = \left[\begin{array}{ccccc} 1 & a_{1,2} & a_{1,3}& \cdots &a_{1,n} \\  0 & 1 & a_{2,3}& \cdots & a_{2,n} \\
%     0 & 0 & 1 & \cdots &  a_{3,n} \\
%     \vdots & \vdots & \vdots & \ddots & \vdots \\
%     0 & \cdots & 0 & 0 & 1\end{array}\right]$.
% \end{center}
% Consider the multiplication matrix $A^{-1}W$ (which is symmetry):
% \begin{center}
%     $W.A^{-1} = \left[\begin{array}{cccc}w_{1,1} & w_{1,2} & \cdots & w_{1,n}\\ w_{2,1} & w_{2,2} & \cdots & w_{2,n}\\ \vdots & \vdots & \ddots & \vdots \\ w_{n,1} & w_{n,2} & \cdots & w_{n,n}\end{array}\right]\left[\begin{array}{ccccc} 1 & a_{1,2} & a_{1,3}& \cdots &a_{1,n} \\  0 & 1 & a_{2,3}& \cdots & a_{2,n} \\
%     0 & 0 & 1 & \cdots &  a_{3,n} \\
%     \vdots & \vdots & \vdots & \ddots & \vdots \\
%     0 & \cdots & 0 & 0 & 1\end{array}\right] $
% \end{center}
% Which is the matrix:
% \begin{center}
%     $ \left[\begin{array}{cccccc}w_{1,1} & a_{1,2}w_{1,1}+w_{1,2} & a_{1,3}w_{1,1}+a_{2,3}w_{1,2}+w_{1,3} & a_{1,4}w_{1,1}+\cdots+w_{1,4} & \cdots & \sum_{i=1}^{n-1} a_{i,n}w_{1,i}+w_{1,n}\\ w_{2,1} & a_{1,2}.w_{2,1}+w_{2,2} & a_{1,3}w_{2,1}+a_{2,3}w_{2,2}+w_{2,3} & a_{1,4}w_{2,1}+\cdots+w_{2,4} &\cdots & \sum_{i=1}^{n-1} a_{i,n}w_{2,i}+w_{2,n}\\  w_{3,1} & a_{1,2}.w_{3,1}+w_{3,2} & a_{1,3}w_{3,1}+a_{2,3}w_{3,2}+w_{3,3} & a_{1,4}w_{3,1}+\cdots+w_{3,4}& \cdots & \sum_{i=1}^{n-1} a_{i,n}w_{3,i}+w_{3,n} \\ w_{4,1} & a_{1,2}.w_{4,1}+w_{4,2} & a_{1,3}w_{4,1}+a_{2,3}w_{4,2}+w_{4,3} & a_{1,4}w_{4,1}+\cdots+w_{4,4}& \cdots & \sum_{i=1}^{n-1} a_{i,n}w_{4,i}+w_{4,n} \\
%     \vdots & \vdots & \vdots & \vdots & \ddots & \vdots \\ w_{n,1} & a_{1,2}w_{n,1}+ w_{n,2} & a_{1,3}w_{n,1}+a_{2,3}w_{n,2} + w_{n,3} & a_{1,4}w_{n,1}+\cdots+w_{n,4}& \cdots & \sum_{i=1}^{n-1} a_{i,n}w_{n,i}+w_{n,n}\end{array}\right]$
% \end{center}
% More precisely, at position $(i,j)$, we have the holomorphic function:
% \begin{center}
%     $\sum_{k=1}^{j-1} a_{k,j}.w_{i,k}+w_{i,j}$
% \end{center}
% and from the symmetric, we have for $(i,j)$ the equation: 
% \begin{center}
%     $\sum_{k=1}^{j-1} a_{k,j}.w_{i,k}+w_{i,j} = \sum_{k=1}^{i-1}a_{k,i}.w_{j,k}+w_{j,i} $.
% \end{center}
% Starting with $(1,2)$, we have: 
% \begin{center}
%     $a_{1,2}.w_{1,1}+w_{1,2} = w_{2,1}$
% \end{center}
% which gives us $w_{1,1} = 0$ for general $a_{1,2}$. Continue the process for $(1,3)$, we have: 
% \begin{center}
%     $a_{1,3}w_{1,1}+a_{2,3}w_{1,2}+w_{1,3} = w_{3,1}$.
% \end{center}
% Together with $w_{1,1} = 0$, we get $w_{1,2} = 0$. Continue the process we will get $w_{1,i} = 0$ for $1 \leq i < n$.\\
% Next, consider position $(2,3)$, we have:
% \begin{center}
%     $a_{1,2}w_{3,1}+w_{3,2} = a_{1,3}w_{2,1}+a_{2,3}w_{2,2}$
% \end{center}
% which gives us $w_{2,2} =0$. Continue this process we get $w_{2,i} = 0$ for $1 \leq i < n-1$. Continuing this process, we will get $w_{i,j} = 0$ for all $i+j \leq n$. Now, we consider the position $(n,2)$: 
% \begin{center}
%     $a_{1,2}w_{n,1} + w_{n,2} = \sum_{i=1}^{n-1} a_{i,n}w_{2,i}+w_{2,n}$
% \end{center}
% Notice that $w_{2,i} = 0$ for all $i \leq n-2$, so we get:
% \begin{center}
%     $a_{1,2}w_{n,1} = a_{n-1,n}w_{2,n-1}$.
% \end{center}
% Let $w_{n,1} = a_{n-1,n}.f_1$, we get $w_{2,n-1} = a_{1,2}.f_1$. Now for position $(n-i,i+2)$, we have:
% \begin{center}
%     $\sum_{k=1}^{i+1}a_{k,i+2}w_{n-i,k}+w_{n-i,i+2} = \sum_{k=1}^{n-i-1}a_{k,n-i}w_{i+2,k}+w_{i+2,n-i} $.
% \end{center}
% From this equation, we get: 
% \begin{center}
%    $a_{i+1,i+2}w_{n-i,i+1} = a_{n-i-1,n-i}w_{i+2,n-i-1}$. 
% \end{center}
% Let $w_{n-i,i+1}=a_{n-i-1,n-i}.f_{i+1}$, then we get $w_{i+2,n-i-1}=a_{i+1,i+2}.f_{i+1}$.\\
% \\
% For $i=1$, we have: 
% \begin{center}
%    $w_{n-1,2}=a_{n-2,n-1}.f_{2}$ and $w_{3,n-2} = a_{2,3}.f_2$. 
% \end{center}
% This gives: 
% \begin{center}
%     $a_{n-2,n-1}.f_2 = a_{1,2}f_1 $.
% \end{center}
% Now we take $f_1 = a_{n-2,n-1}.f$, then $f_2 = a_{1,2}f$. This implies:
% \begin{center}
% $w_{n,1}=a_{n-1,n}.a_{n-2,n-1}f$, $w_{n-1,2}=a_{n-2,n-1}.a_{1,2}f$ and $w_{3,n-2} = a_{2,3}.a_{1,2}f$. 
% \end{center}
% Thus from the symmetric condition of (\ref{equa-10}), we have the matrix:
% \begin{center}
%     $W = \left[\begin{array}{cccccc}0 & \cdots & 0 & 0 & 0 &a_{n-1,n}.a_{n-2,n-1}f\\ 0 & \cdots & 0 & 0 & a_{n-2,n-1}.a_{1,2}f & w_{2,n}\\  0 & \cdots & 0 & a_{2,3}.a_{1,2}f & w_{3,n-1} & w_{3,n} \\
%     \vdots & \vdots & \vdots & \vdots & \vdots & \vdots \\ 0 & a_{n-2,n-1}.a_{1,2}f & w_{n-1,3} & \cdots & w_{n-1,n-1} & w_{n-1,n} \\ a_{n-1,n}.a_{n-2,n-1}f & w_{n,2} & w_{n,3} & \cdots & w_{n,n-1} &w_{n,n}\end{array}\right] $
% \end{center}
% By the same method, we get:
% \begin{center}
%     $V = \left[\begin{array}{cccccc}v_{1,1} & v_{1,2} & v_{1,3} & \cdots & v_{1,n-1} & a_{1,2}a_{2,3}g\\ v_{2,1} & v_{2,2} & \cdots & v_{2,n-2}  & a_{2,3}.a_{n-1,n}g &0\\ v_{3,1} & \cdots & v_{3,n-3} & a_{n-2,n-1}a_{n-1,n} g & 0 & 0 \\ \vdots & \vdots & \vdots & \vdots & \vdots & \vdots \\ w_{n-1,1} & a_{2,3}.a_{n-1,n}g & 0 & \cdots & 0 & 0 \\ a_{1,2}a_{2,3}g & 0 & 0 & \cdots & 0 & 0\end{array}\right] $
% \end{center}
% From the matrices equation $VW = A - I$, we get the system of equations:
% \begin{equation}
%     \begin{cases}
%      fg = 0\\
%       a_{n-2,n-1}.a_{1,2}fv_{1,n-1}+a_{1,2}a_{2,3}g.w_{n,2} = -a_{1,2}\\
%       a_{2,3}.a_{1,2}fv_{2,n-2}+a_{2,3}.a_{n-1,n}gw_{n-1,3} = -a_{2,3}\\
%       a_{n-1,n}.a_{n-2,n-1}fv_{n-1,1}+a_{2,3}.a_{n-1,n}gw_{2,n} = -a_{n-1,n}\\
%       a_{n-2,n-1}a_{1,2}f.v_{n-2,2}+a_{n-2,n-1}a_{n-1,n}g.w_{3,n-1} = -a_{n-2,n-1}\\
     
% \end{cases}
% \end{equation}
% Let $f =0$, we get:
% \begin{equation}
%     \begin{cases}
%       a_{2,3}g.w_{n,2} = -1\\
%       a_{n-1,n}gw_{n-1,3} = -1\\
%       a_{2,3}gw_{2,n} = -1\\
%       a_{n-1,n}g.w_{3,n-1} = -1.\\
     
% \end{cases}
% \end{equation}
% These equations show us that we could not have $g$, $w_{n,2}$ and $w_{n-1,3}$ such that they are holomorphic functions with variables $a_{i,j}$. Thus this implies the proof for the case $M^{+}(A,B)$, and the case $M^{-}(A,B)$ follows directly. 
% \end{proof}
% From Proposition \ref{fact-num-prop} and Proposition \ref{fact-num-prop2}, we have a theorem for possibility of factorizing elementary matrices $M^+(A,B)$ and $M^-(A,C)$ into 4 elementary matrices of type $M^+(I,Z)$ and $M^-(I,Z)$.
% \begin{theo}\label{fact-num-theo}
%  The elementary matrices $M^+(A,B)$, $M^-(A,C) \in \mathrm{Sp}_{2n}(O(X))$ with $n \geq 4$ could not always factorize in product of 4 elementary matrices of type $M^+(I,Z)$ and $M^-(I,Z)$.    
% \end{theo}
% \subsection{Factorization of elementary matrices with five factors}
% In this section, we will try to find optimal number of factorization for $H$. First, we have that matrices of the form $\left[\begin{array}{ccc} A & 0 \\ 0 & A^{-T}\end{array}\right] $ in general could not factorize into product of 5 elementary matrices.
% \begin{theo}\label{fact-num-theo2}
% The symplectic matrices of the form $\left[\begin{array}{ccc} A & 0 \\ 0 & A^{-T}\end{array}\right] $ could not always be factorized into product of 5 elementary matrices of type $M^+(I,Z)$ and $M^-(I,Z)$. 
% \end{theo}
% \begin{proof}
% Suppose we have that:
% \begin{center}
%     $ \left[\begin{array}{ccc} A & 0 \\ 0 & A^{-T}\end{array}\right]= \left[\begin{array}{ccc} I & 0 \\ U & I\end{array}\right]\left[\begin{array}{ccc} I & V \\ 0 & I\end{array}\right]\left[\begin{array}{ccc} I & 0 \\ W & I\end{array}\right]\left[\begin{array}{ccc} I & X \\ 0 & I\end{array}\right]\left[\begin{array}{ccc} I & 0 \\ Y & I\end{array}\right]$.
% \end{center}
% Then this is equivalence to 
% \begin{center}
%     $ \left[\begin{array}{ccc} A & 0 \\ -A^T.Y & A^{-T}\end{array}\right]=\left[\begin{array}{ccc} A & 0 \\ 0 & A^{-T}\end{array}\right]\left[\begin{array}{ccc} I & 0 \\ -Y & I\end{array}\right]= \left[\begin{array}{ccc} I & 0 \\ U & I\end{array}\right]\left[\begin{array}{ccc} I & V \\ 0 & I\end{array}\right]\left[\begin{array}{ccc} I & 0 \\ W & I\end{array}\right]\left[\begin{array}{ccc} I & X \\ 0 & I\end{array}\right]$
% \end{center}
% It is easy to see that $\left[\begin{array}{ccc} A & 0 \\ -A^T.Y & A^{-T}\end{array}\right]$ is a symplectic matrix. This means $\left[\begin{array}{ccc} A & 0 \\ -A^T.Y & A^{-T}\end{array}\right]$ can be factorized into product of 4 elementary matrices, which contradicts Theorem \ref{fact-num-theo}.
% \end{proof}

\begin{proof}[Proof of Theorem \ref{main-theo2}]
By Theorem \ref{fact-ele}, $\mathrm{K}_{symp}(n,d) \leq 7 \tilde{\mathrm{K}}_{symp}(n,d)$. 
Then by Remark \ref{fact-num-theo1} (b), $6 \leq \mathrm{K}_{symp}(n,d)$.
In the case $n \leq 3$, by Remark \ref{fact-num-theo1} (a), $\mathrm{K}_{symp}(n,d) \leq 4\tilde{\mathrm{K}}_{symp}(n,d)$. 
\end{proof}

\begin{proof}[Proof of Proposition \ref{bounds}]
    The inequality \eqref{K(N,d)} follows from \cite[Theorem 2.4]{HKS} and Theorem \ref{main-theo2}.
    By \cite[Theorem 1]{HKS}, 
    \[
        \tilde{K}_{symp}(n,1) = 4 \text{ and }\tilde{K}_{symp}(n,2) \le 5 \text{ for all } n.
    \]
    Together with the inequalities \eqref{K-Ktilde}, this shows the upper bounds in \eqref{K(n,1)n=2,3}, \eqref{K(n,1)n>=4}, \eqref{K(n,2)n=2,3} and \eqref{K(n,2)n>=4}. 
     The lower bounds in \eqref{K(n,1)n=2,3} and \eqref{K(n,2)n=2,3} hold even for matrices with constant entries \cite{JTZ}, while the lower bounds in \eqref{K(n,1)n>=4} and \eqref{K(n,2)n>=4} follow from Theorem \ref{main-theo2}. 
\end{proof}

\section*{Funding}

All authors were partially supported by Schweizerischer Nationalfonds grant 200021-207335.

%\nocite{*}

\begin{thebibliography}{9}
\bibitem[For10]{For1} F. Forstneri{\v c}, \textit{The Oka principle for sections of stratified fiber bundles}, Pure Appl.
Math. Q. \textbf{6} (2010), no. 3, 843–874. 


\bibitem[For17]{For2} F. Forstneri{\v c}, \textit{Stein manifolds and holomorphic mappings}, Second edition, Springer, Cham, 2017.


\bibitem[Gro89]{Gro} M. Gromov, \textit{Oka’s principle for holomorphic sections of elliptic bundles}, J. Amer. Math. Soc. \textbf{2} (1989), no. 4, 851–897.

\bibitem[HKS22]{HKS}
G. Huang, F. Kutzschebauch and J. Schott. \textit{Factorization of holomorphic matrices and Kazhdan's property (T)}, Bull. Sci. Math. \textbf{190} (2024), 103376.


\bibitem[IK12]{KI}
B. Ivarsson, F. Kutzschebauch. \textit{Holomorphic factorization of mappings into }$SL_n(\mathbb{C})$. Ann. Math. \textbf{175} (2012), no. 1., 45 - 69.

\bibitem[IKL22]{IKL}
B. Ivarsson, F. Kutzschebauch, E. Løw. \textit{Holomorphic factorization of mappings into }$Sp_4(\mathbb{C})$, Anal. PDE, \textbf{16} (2023), no. 1, 233–277.

\bibitem[JTZ20]{JTZ} 
    P. Jin, Y. Tang, A. Zhu, 
    \textit{Unit triangular factorization of the matrix symplectic group}. SIAM J. Matrix Anal. Appl. \textbf{41} (2020), no. 4, 1630–1650.
    
\bibitem[JLX22]{JLX2}
    P. Jin, Z. Lin, B. Xiao, 
    \textit{Optimal unit triangular factorization of symplectic matrices}, Linear Algebra Appl. \textbf{650} (2022), 236–247.
   
% \bibitem[Lei18]{lei2}  
%      J. Leiterer,
%     \textit{Similarity of holomorphic matrices on 1-dimensional Stein spaces}, arXiv:1710.10969, 2018. 
% \bibitem[Lei23]{lei1}
%     J. Leiterer,
%     \textit{On the Jordan structure of holomorphic matrices}, Revue Roumaine Mathematiques Pures Appliquees, 2023.  

\bibitem[Scho24]{Scho}
    J. Schott,
    \textit{Holomorphic factorization of mappings into the symplectic group}. J. Eur. Math. Soc. (2025), published online first.
    
\bibitem[VSS11]{VSS} 
    N. A. Vavilov, A. V. Smolenskiĭ, B. Sury, 
    \textit{Unitriangular factorizations of Chevalley groups}. 
    Zap. Nauchn. Sem. S.-Peterburg. Otdel. Mat. Inst. Steklov. (POMI) 
    \textbf{388} (2011), 17–47, 309–310; translation in
    J. Math. Sci. (N.Y.) \textbf{183} (2012), no. 5, 584–599.

\end{thebibliography}
\end{document}
